%% GENERATED by generate_papers.py -- edit the bank_*.py files, not this file.
\documentclass[11pt,a4paper]{article}
\usepackage[a4paper, top=20mm, bottom=20mm, left=18mm, right=18mm, headheight=14pt]{geometry}
\usepackage[T1]{fontenc}
\usepackage{lmodern}
\usepackage{amsmath, amssymb, mathtools}
\usepackage{braket}
\usepackage{array, tabularx, booktabs, multicol}
\usepackage{enumitem}
\usepackage[dvipsnames]{xcolor}
\usepackage[most]{tcolorbox}
\usepackage{needspace}
\usepackage{fancyhdr, lastpage}
\usepackage{pifont}
\usepackage{tikz}
\usepackage[hidelinks]{hyperref}
\newcolumntype{L}{>{\raggedright\arraybackslash}X}
\definecolor{brand}{HTML}{1F4E79}
\definecolor{okgreen}{HTML}{1B7F3B}
\setlength{\parindent}{0pt}
\setlength{\parskip}{3pt}
\newcommand{\optlabel}[1]{\textbf{(#1)}}
\newcommand{\correctlabel}[1]{\textcolor{okgreen}{\textbf{(#1)}\,\ding{51}}}
\newenvironment{qitem}[1]{\begin{list}{}{\setlength{\leftmargin}{3.3em}\setlength{\labelwidth}{3.1em}%
  \setlength{\labelsep}{0.2em}\setlength{\topsep}{3pt}\setlength{\parsep}{3pt}}\item[\textbf{Q.#1}\hfill]}{\end{list}}
\newcommand{\qsep}{\par\vspace{2pt}{\color{black!15}\hrule}\vspace{4pt}}
\newcommand{\qmeta}[2]{{\footnotesize\color{black!60}\textsc{#1} \;$\cdot$\; #2}\par}
\newtcolorbox{passagebox}[1]{enhanced, breakable, colback=brand!4, colframe=brand!70, fonttitle=\bfseries,
  title={#1}, left=6pt, right=6pt, top=4pt, bottom=4pt}
\newtcolorbox{sectionbanner}{enhanced, colback=brand, colframe=brand, coltext=white, fontupper=\large\bfseries,
  left=8pt, right=8pt, top=6pt, bottom=6pt, sharp corners}
\newtcolorbox{infobox}[1]{enhanced, breakable, colback=black!3, colframe=black!45, fonttitle=\bfseries,
  title={#1}, left=6pt, right=6pt, top=4pt, bottom=4pt}
\newtcolorbox{solbox}[1]{enhanced, breakable, colback=okgreen!3, colframe=okgreen!60!black, fonttitle=\bfseries,
  title={Correct answer: (#1) \quad Detailed solution}, left=6pt, right=6pt, top=3pt, bottom=4pt, before skip=4pt, after skip=6pt}
\newcommand{\bub}[1]{\tikz[baseline=-0.65ex]\node[draw, circle, inner sep=0pt, minimum size=4.6mm, font=\scriptsize]{#1};}
\allowdisplaybreaks

\begin{document}
\pagestyle{fancy}\fancyhf{}
\fancyhead[L]{\small\color{black!70}Grover's Algorithm --- Assessment}
\fancyhead[R]{\small\color{black!70}\textbf{Answer Book}}
\fancyfoot[C]{\small Page \thepage\ of \pageref*{LastPage}}
\renewcommand{\headrulewidth}{0.4pt}

\thispagestyle{empty}
\begin{center}
{\color{brand}\rule{\linewidth}{1.2pt}}\\[6pt]
{\Huge\bfseries\color{brand} Grover's Algorithm}\\[4pt]
{\Large Assessment of quantum search}\\[4pt]
{\large\bfseries ANSWER BOOK}\\[2pt]
{\color{brand}\rule{\linewidth}{1.2pt}}
\end{center}
This book accompanies the separate \emph{Question Paper} (80 questions, 3 hours, 240 marks). It contains
\begin{enumerate}[leftmargin=1.4em, itemsep=0pt, topsep=2pt]
\item the \textbf{answer key} with the topic tested by every question (page~\pageref{key}), and
\item the \textbf{Answer Bank}: every question restated with the correct option marked \textcolor{okgreen}{\ding{51}} and a
      complete solution: concepts required, relevant formulae, step-by-step explanation, why each other option is
      incorrect, and a key takeaway (from page~\pageref{bank}).
\end{enumerate}
\begin{infobox}{Marking scheme (every question)}
Each question has \textbf{four} options (A), (B), (C) and (D). \textbf{ONLY ONE} of these four options is correct.\\[2pt]
\begin{tabular}{@{}ll@{}}
\emph{Full marks}: & $+3$ if ONLY the correct option is chosen;\\
\emph{Zero marks}: & $\phantom{+}0$ if none of the options is chosen (question unanswered);\\
\emph{Negative marks}: & $-1$ in all other cases (wrong option, or more than one option chosen).
\end{tabular}
\end{infobox}

\begin{infobox}{How to evaluate}
For each section count $C$ = number of correct answers and $W$ = number of questions answered wrongly (or with more than
one bubble darkened). Section marks $= 3C - W$; unanswered questions score $0$. The total is out of $240$.
Every numerical answer was re-computed independently (NumPy, including gate-by-gate simulation of the Grover circuits, and CUDA-Q running the notebook's own kernels on
the CPU simulator) by \texttt{verify\_answers.py} in the source folder.
\end{infobox}
\begin{infobox}{Conventions (identical to the Question Paper)}

\begin{itemize}[leftmargin=1.2em, itemsep=1pt, topsep=2pt]
\item $n$ qubits, $N = 2^n$ items, $M$ marked items ($M = 1$ unless stated otherwise). A basis state is written $\ket{x} = \ket{x_1x_2\dots x_n}$:
      \textbf{qubit 1 is the most significant bit} and is written first. (The accompanying notebook writes \texttt{N\_QUBITS} for $n$ and $M = 2^{N}$ for the number of items, and $O$ for the oracle.)
\item $\ket{s} = H^{\otimes n}\ket{0\dots0} = \frac{1}{\sqrt N}\sum_x\ket{x}$. The oracle is $U_\omega\ket{x} = (-1)^{f(x)}\ket{x}$ with $f(x) = 1$ exactly on the marked items; the diffusion is $D = 2\ket{s}\!\bra{s} - I$; one iteration is $G = DU_\omega$ (oracle first).
\item $\sin\theta = \sqrt{M/N}$ with $0 < \theta \le \frac\pi2$; $\;\ket{m}$ and $\ket{r}$ are the uniform superpositions of the marked and the unmarked items. Then
      $\displaystyle G^k\ket{s} = \sin\big((2k+1)\theta\big)\ket{m} + \cos\big((2k+1)\theta\big)\ket{r}$, so $P_k = \sin^2\big((2k+1)\theta\big)$ and
      $k_{\rm opt} = \operatorname{round}\big(\frac{\pi}{4\theta} - \frac12\big)$.
\item \textbf{Textbook circuit:} $H^{\otimes n}$, then $k$ times: [$X$ on each qubit where $x^\ast$ has a $0$; a multi-controlled $Z$ on all $n$ qubits; the same $X$ gates] (no ancillas) followed by [$H^{\otimes n}$; $X^{\otimes n}$; a multi-controlled $Z$; $X^{\otimes n}$; $H^{\otimes n}$]; finally measure all qubits.
\item $H = \frac{1}{\sqrt2}\begin{psmallmatrix}1&1\\1&-1\end{psmallmatrix}$; $X, Y, Z$ are the Pauli matrices; $\ket{\pm} = \frac{1}{\sqrt2}(\ket0 \pm \ket1)$. Two-qubit matrices use the basis order $\ket{00}, \ket{01}, \ket{10}, \ket{11}$.
\item All measurements are in the computational basis unless stated otherwise; $\log$ means $\log_2$. A numerical value given in a question may be used without a calculator.
\end{itemize}
\end{infobox}
\newpage
\section*{Answer key}\label{key}
{\small Answer distribution: (A) 20, (B) 20, (C) 20, (D) 20.}

\needspace{14\baselineskip}
\subsection*{Section A --- Mathematics and Search Foundations}
{\small\renewcommand{\arraystretch}{1.15}
\begin{tabularx}{\linewidth}{@{}r c L @{\hspace{1em}} r c L@{}}
\toprule
\textbf{Q.} & \textbf{Ans.} & \textbf{Topic tested} & \textbf{Q.} & \textbf{Ans.} & \textbf{Topic tested}\\ \midrule
1 & \textbf{B} & Triple-angle identity & 11 & \textbf{A} & The signed sum $N - 2M$ \\
2 & \textbf{C} & Reflection matrix $2uu^{T} - I$ & 12 & \textbf{B} & Normalising a state vector \\
3 & \textbf{D} & Two reflections make a rotation & 13 & \textbf{B} & Square-root scaling \\
4 & \textbf{D} & Eigenvalues of a rotation & 14 & \textbf{C} & Complex inner product \\
5 & \textbf{A} & Powers of a complex phase & 15 & \textbf{C} & $\sin^2 3\theta$ as a polynomial \\
6 & \textbf{D} & Optimal integer $k$; why $-\frac12$ matters & 16 & \textbf{D} & Traces of projectors and reflections \\
7 & \textbf{A} & Inversion about the mean & 17 & \textbf{C} & Rounding rule vs.\ $\arcsin x > x$ \\
8 & \textbf{C} & Classical search without replacement & 18 & \textbf{A} & Spectrum of a reflection \\
9 & \textbf{A} & Binary vs. linear search & 19 & \textbf{D} & First maximum and period of $p(k)$ \\
10 & \textbf{B} & Counting solutions of a Boolean function & 20 & \textbf{B} & Tolerance window around the peak \\
\bottomrule
\end{tabularx}}

\needspace{14\baselineskip}
\subsection*{Section B --- Quantum Computing Foundations}
{\small\renewcommand{\arraystretch}{1.15}
\begin{tabularx}{\linewidth}{@{}r c L @{\hspace{1em}} r c L@{}}
\toprule
\textbf{Q.} & \textbf{Ans.} & \textbf{Topic tested} & \textbf{Q.} & \textbf{Ans.} & \textbf{Topic tested}\\ \midrule
21 & \textbf{B} & Uniform state, partial outcome & 31 & \textbf{A} & Overlap $\cos\theta$ \\
22 & \textbf{A} & Tensor-product amplitude & 32 & \textbf{D} & State-vector memory \\
23 & \textbf{C} & Hadamard layer signs & 33 & \textbf{A} & Diffusion for one qubit \\
24 & \textbf{C} & Phase kickback with $\ket{-}$ & 34 & \textbf{D} & $\ket{s}$ is fixed by $D$ \\
25 & \textbf{D} & No kickback with $\ket{+}$ & 35 & \textbf{C} & Single-qubit gate actions \\
26 & \textbf{C} & $X$ conjugation selects the item & 36 & \textbf{B} & Product vs.\ entangled states \\
27 & \textbf{D} & CCZ from Toffoli & 37 & \textbf{A} & Why the oracle alone reveals nothing \\
28 & \textbf{A} & Global vs.\ relative phase & 38 & \textbf{C} & Relative phases are observable \\
29 & \textbf{B} & Subspace measurement & 39 & \textbf{D} & Toffoli as a phase oracle \\
30 & \textbf{B} & Unitary vs.\ projector & 40 & \textbf{B} & Ancilla stays unentangled \\
\bottomrule
\end{tabularx}}

\needspace{14\baselineskip}
\subsection*{Section C --- Grover's Algorithm Fundamentals}
{\small\renewcommand{\arraystretch}{1.15}
\begin{tabularx}{\linewidth}{@{}r c L @{\hspace{1em}} r c L@{}}
\toprule
\textbf{Q.} & \textbf{Ans.} & \textbf{Topic tested} & \textbf{Q.} & \textbf{Ans.} & \textbf{Topic tested}\\ \midrule
41 & \textbf{D} & $N = 4$ is exact in one step & 51 & \textbf{D} & $M = 3N/4$ \\
42 & \textbf{B} & One iteration, $N = 8$ & 52 & \textbf{A} & Over-rotation at $2k_{\rm opt}$ \\
43 & \textbf{D} & $k_{\rm opt}$ and $P$ for $n = 4$ & 53 & \textbf{A} & Size of the speed-up \\
44 & \textbf{B} & Matrix of $D$ for $N = 4$ & 54 & \textbf{B} & Shot statistics \\
45 & \textbf{C} & Building the oracle & 55 & \textbf{A} & $k_{\rm opt}$ against $n$ \\
46 & \textbf{A} & Gate count per iteration & 56 & \textbf{A} & Blocks of the circuit \\
47 & \textbf{B} & Diffusion circuit $= -D$ & 57 & \textbf{B} & Independence of the marked item \\
48 & \textbf{C} & Unmarked amplitude & 58 & \textbf{D} & Oscillation and $\sqrt N$ scaling \\
49 & \textbf{C} & $M = N/4$ & 59 & \textbf{C} & Amplitude trace, first iteration \\
50 & \textbf{D} & $M = N/2$ & 60 & \textbf{C} & Amplitude trace, second iteration \\
\bottomrule
\end{tabularx}}

\needspace{14\baselineskip}
\subsection*{Section D --- Advanced Grover and Implementation}
{\small\renewcommand{\arraystretch}{1.15}
\begin{tabularx}{\linewidth}{@{}r c L @{\hspace{1em}} r c L@{}}
\toprule
\textbf{Q.} & \textbf{Ans.} & \textbf{Topic tested} & \textbf{Q.} & \textbf{Ans.} & \textbf{Topic tested}\\ \midrule
61 & \textbf{B} & Several solutions: $k_{\rm opt}$ & 71 & \textbf{B} & Simulation cost \\
62 & \textbf{A} & BBHT threshold for unknown $M$ & 72 & \textbf{C} & Reading the sweep code \\
63 & \textbf{D} & Exact Grover search & 73 & \textbf{D} & Interpreting $2000/2000$ shots \\
64 & \textbf{D} & Optimality (BBBV, Zalka) & 74 & \textbf{C} & Repeat-until-success cost \\
65 & \textbf{C} & Reading gate statistics & 75 & \textbf{A} & Query complexities \\
66 & \textbf{B} & Amplitude amplification cost & 76 & \textbf{B} & Variants of Grover \\
67 & \textbf{D} & Notebook memory estimate & 77 & \textbf{D} & Grover and NP-completeness \\
68 & \textbf{A} & Key search & 78 & \textbf{A} & Unknown $M$; flat curve at $M = N/2$ \\
69 & \textbf{C} & Noise and depth & 79 & \textbf{B} & Quantum counting outcomes \\
70 & \textbf{A} & Toffoli count & 80 & \textbf{C} & Quantum counting estimate \\
\bottomrule
\end{tabularx}}

\subsection*{Quick-marking grid}
\begin{center}\renewcommand{\arraystretch}{1.3}
\begin{tabular}{|c|c|c|c|c|c|c|c|c|c|}\hline
1\,\textbf{B} & 2\,\textbf{C} & 3\,\textbf{D} & 4\,\textbf{D} & 5\,\textbf{A} & 6\,\textbf{D} & 7\,\textbf{A} & 8\,\textbf{C} & 9\,\textbf{A} & 10\,\textbf{B} \\
\hline
11\,\textbf{A} & 12\,\textbf{B} & 13\,\textbf{B} & 14\,\textbf{C} & 15\,\textbf{C} & 16\,\textbf{D} & 17\,\textbf{C} & 18\,\textbf{A} & 19\,\textbf{D} & 20\,\textbf{B} \\
\hline
21\,\textbf{B} & 22\,\textbf{A} & 23\,\textbf{C} & 24\,\textbf{C} & 25\,\textbf{D} & 26\,\textbf{C} & 27\,\textbf{D} & 28\,\textbf{A} & 29\,\textbf{B} & 30\,\textbf{B} \\
\hline
31\,\textbf{A} & 32\,\textbf{D} & 33\,\textbf{A} & 34\,\textbf{D} & 35\,\textbf{C} & 36\,\textbf{B} & 37\,\textbf{A} & 38\,\textbf{C} & 39\,\textbf{D} & 40\,\textbf{B} \\
\hline
41\,\textbf{D} & 42\,\textbf{B} & 43\,\textbf{D} & 44\,\textbf{B} & 45\,\textbf{C} & 46\,\textbf{A} & 47\,\textbf{B} & 48\,\textbf{C} & 49\,\textbf{C} & 50\,\textbf{D} \\
\hline
51\,\textbf{D} & 52\,\textbf{A} & 53\,\textbf{A} & 54\,\textbf{B} & 55\,\textbf{A} & 56\,\textbf{A} & 57\,\textbf{B} & 58\,\textbf{D} & 59\,\textbf{C} & 60\,\textbf{C} \\
\hline
61\,\textbf{B} & 62\,\textbf{A} & 63\,\textbf{D} & 64\,\textbf{D} & 65\,\textbf{C} & 66\,\textbf{B} & 67\,\textbf{D} & 68\,\textbf{A} & 69\,\textbf{C} & 70\,\textbf{A} \\
\hline
71\,\textbf{B} & 72\,\textbf{C} & 73\,\textbf{D} & 74\,\textbf{C} & 75\,\textbf{A} & 76\,\textbf{B} & 77\,\textbf{D} & 78\,\textbf{A} & 79\,\textbf{B} & 80\,\textbf{C} \\
\hline\end{tabular}\end{center}
\newpage
\section*{Answer Bank --- questions with detailed solutions}\label{bank}

\begin{sectionbanner}SECTION A \enspace$\cdot$\enspace Mathematics and Search Foundations \hfill Q.1--20\end{sectionbanner}

\needspace{13\baselineskip}
\begin{qitem}{1}
\qmeta{Single-correct MCQ}{Triple-angle identity}
If $\sin\theta = \tfrac14$, the value of $\sin^2 3\theta$ is
\par\smallskip\noindent\begin{tabularx}{\linewidth}{@{}*{4}{L}@{}}
\optlabel{A}\,$\dfrac{1}{16}$ & \correctlabel{B}\,$\dfrac{121}{256}$ & \optlabel{C}\,$\dfrac{9}{16}$ & \optlabel{D}\,$\dfrac{11}{16}$
\end{tabularx}
\end{qitem}
\begin{solbox}{B}
\textbf{\color{brand}Concepts required.} Triple-angle identity $\sin3\theta = 3\sin\theta - 4\sin^3\theta$. In Grover's algorithm $P_1 = \sin^2 3\theta$.

\textbf{\color{brand}Relevant formulae.} $\sin3\theta = 3\sin\theta - 4\sin^3\theta$, where $\theta$ is any angle.

\textbf{\color{brand}Step-by-step explanation.} Use the triple-angle identity $\sin 3\theta = 3\sin\theta - 4\sin^3\theta$:
\[ \sin 3\theta = \frac34 - 4\cdot\frac{1}{64} = \frac{12}{16} - \frac{1}{16} = \frac{11}{16},
   \qquad \sin^2 3\theta = \frac{121}{256}. \]
This is exactly the success probability of one Grover iteration on $N = 16$ items.
The value $\frac{9}{16} = (3\sin\theta)^2$ is the small-angle estimate; it ignores the cubic correction.

\textbf{\color{brand}Why the other options are incorrect.}
\begin{itemize}[leftmargin=2.4em, itemsep=1pt, topsep=2pt]
\item[\textbf{(A)}] $\frac{1}{16} = \sin^2\theta$ is the probability before any iteration.
\item[\textbf{(C)}] $\frac{9}{16} = (3\sin\theta)^2$ uses the small-angle estimate $\sin3\theta \approx 3\sin\theta$ and drops $-4\sin^3\theta$.
\item[\textbf{(D)}] $\frac{11}{16}$ is $\sin3\theta$ itself: the final squaring was forgotten.
\end{itemize}
\textbf{\color{brand}Key takeaway.} Keep the exact identity for small $N$; the small-angle estimate is only reliable for large $N$.
\end{solbox}
\qsep

\needspace{13\baselineskip}
\begin{qitem}{2}
\qmeta{Single-correct MCQ}{Reflection matrix $2uu^{T} - I$}
Let $u = \big(\tfrac35, \tfrac45\big)^{T}$. The matrix $2uu^{T} - I$ equals
\par\smallskip\noindent\begin{tabularx}{\linewidth}{@{}*{4}{L}@{}}
\optlabel{A}\,$\dfrac{1}{25}\begin{pmatrix}18 & 24\\ 24 & 32\end{pmatrix}$ & \optlabel{B}\,$\dfrac{1}{25}\begin{pmatrix}-7 & 12\\ 12 & 7\end{pmatrix}$ & \correctlabel{C}\,$\dfrac{1}{25}\begin{pmatrix}-7 & 24\\ 24 & 7\end{pmatrix}$ & \optlabel{D}\,$\dfrac{1}{25}\begin{pmatrix}7 & -24\\ -24 & -7\end{pmatrix}$
\end{tabularx}
\end{qitem}
\begin{solbox}{C}
\textbf{\color{brand}Concepts required.} $uu^{T}$ is the matrix with entries $u_iu_j$. $2uu^{T} - I$ reflects vectors in the line along $u$: it fixes $u$ and negates the perpendicular direction.

\textbf{\color{brand}Relevant formulae.} $(uu^{T})_{jk} = u_ju_k$ for a real unit vector $u$; $I$ is the identity matrix. $R = 2uu^{T} - I$ satisfies $Ru = u$.

\textbf{\color{brand}Step-by-step explanation.} $uu^{T} = \frac{1}{25}\begin{psmallmatrix}9 & 12\\ 12 & 16\end{psmallmatrix}$, so
$2uu^{T} = \frac{1}{25}\begin{psmallmatrix}18 & 24\\ 24 & 32\end{psmallmatrix}$. Subtracting $I = \frac{1}{25}\begin{psmallmatrix}25 & 0\\ 0 & 25\end{psmallmatrix}$ gives
$\frac{1}{25}\begin{psmallmatrix}-7 & 24\\ 24 & 7\end{psmallmatrix}$.
\emph{Check:} it fixes $u$, since $\frac{1}{25}\cdot\frac15(-21 + 96,\ 72 + 28) = \frac{1}{125}(75, 100) = (\frac35, \frac45)$. Its trace is $0$ and its
determinant is $-1$, as for every reflection in the plane.

\textbf{\color{brand}Why the other options are incorrect.}
\begin{itemize}[leftmargin=2.4em, itemsep=1pt, topsep=2pt]
\item[\textbf{(A)}] This is $2uu^{T}$ alone: the $I$ was not subtracted, and the matrix is not orthogonal.
\item[\textbf{(B)}] The off-diagonal entries are those of $uu^{T}$ and were not doubled; the columns are then not unit vectors.
\item[\textbf{(D)}] This is $I - 2uu^{T}$, the negative of the required matrix; it sends $u$ to $-u$.
\end{itemize}
\textbf{\color{brand}Key takeaway.} Check a reflection by testing that it fixes $u$, and that its trace is $0$ and determinant $-1$ in the plane.
\end{solbox}
\qsep

\needspace{13\baselineskip}
\begin{qitem}{3}
\qmeta{Single-correct MCQ}{Two reflections make a rotation}
A vector in the plane is reflected first in the line through the origin at angle $10^\circ$ to the $x$-axis, and then in the
line through the origin at angle $35^\circ$. The combined map is
\par\smallskip\noindent\begin{tabularx}{\linewidth}{@{}*{2}{L}@{}}
\optlabel{A}\,a reflection in the line at angle $45^\circ$ & \optlabel{B}\,a rotation by $+25^\circ$ \\[3pt]
\optlabel{C}\,a rotation by $-50^\circ$ & \correctlabel{D}\,a rotation by $+50^\circ$
\end{tabularx}
\end{qitem}
\begin{solbox}{D}
\textbf{\color{brand}Concepts required.} Reflection in the line at angle $\alpha$: $\phi \mapsto 2\alpha - \phi$. Two reflections compose to a rotation by twice the angle between the lines.

\textbf{\color{brand}Relevant formulae.} A reflection in the line at angle $\alpha$ maps angle $\phi$ to $2\alpha - \phi$. Reflecting at $\alpha$ and then at $\beta$ rotates by $2(\beta - \alpha)$.

\textbf{\color{brand}Step-by-step explanation.} A reflection in the line at angle $\alpha$ sends a vector at angle $\phi$ to angle $2\alpha - \phi$. Applying $\alpha = 10^\circ$ and then
$\beta = 35^\circ$:
\[ \phi \;\to\; 20^\circ - \phi \;\to\; 70^\circ - (20^\circ - \phi) = \phi + 50^\circ . \]
The result is a rotation by $2(\beta - \alpha) = +50^\circ$: \textbf{twice} the angle between the lines, in the direction from the
first line to the second. This is the geometric heart of Grover's algorithm.

\textbf{\color{brand}Why the other options are incorrect.}
\begin{itemize}[leftmargin=2.4em, itemsep=1pt, topsep=2pt]
\item[\textbf{(A)}] Two reflections have determinant $(-1)(-1) = +1$, so the product is a rotation and cannot be a reflection.
\item[\textbf{(B)}] $25^\circ$ is the angle between the lines; the rotation angle is twice this.
\item[\textbf{(C)}] $-50^\circ$ is the result for the opposite order (first $35^\circ$, then $10^\circ$).
\end{itemize}
\textbf{\color{brand}Key takeaway.} Two reflections make a rotation by twice the angle between their axes: this is why $G$ rotates by $2\theta$.
\end{solbox}
\qsep

\needspace{13\baselineskip}
\begin{qitem}{4}
\qmeta{Single-correct MCQ}{Eigenvalues of a rotation}
Let $\sin\theta = \tfrac14$ with $0 < \theta < \tfrac\pi2$, and let
$G = \begin{psmallmatrix}\cos 2\theta & -\sin 2\theta\\ \sin 2\theta & \cos 2\theta\end{psmallmatrix}$. The eigenvalues of $G$ are
\par\smallskip\noindent\begin{tabularx}{\linewidth}{@{}*{4}{L}@{}}
\optlabel{A}\,$\dfrac{7 \pm \sqrt{15}}{8}$ & \optlabel{B}\,$\dfrac{1 \pm i\sqrt{15}}{4}$ & \optlabel{C}\,$\pm 1$ & \correctlabel{D}\,$\dfrac{7 \pm i\sqrt{15}}{8}$
\end{tabularx}
\end{qitem}
\begin{solbox}{D}
\textbf{\color{brand}Concepts required.} A plane rotation by $\alpha$ has eigenvalues $e^{\pm i\alpha}$, and $\cos2\theta = 1 - 2\sin^2\theta$.

\textbf{\color{brand}Relevant formulae.} A rotation by $\alpha$ has eigenvalues $\lambda_\pm = e^{\pm i\alpha}$; $\;\cos2\theta = 1 - 2\sin^2\theta$; $\;\sin2\theta = 2\sin\theta\cos\theta$.

\textbf{\color{brand}Step-by-step explanation.} A rotation by angle $\alpha$ has eigenvalues $e^{\pm i\alpha} = \cos\alpha \pm i\sin\alpha$ (trace $2\cos\alpha$, determinant $1$).
Here $\alpha = 2\theta$, with
\[ \cos 2\theta = 1 - 2\sin^2\theta = 1 - \frac{2}{16} = \frac78, \qquad
   \sin 2\theta = 2\cdot\frac14\cdot\frac{\sqrt{15}}{4} = \frac{\sqrt{15}}{8}. \]
Therefore the eigenvalues are $\frac{7 \pm i\sqrt{15}}{8}$; their modulus is $\sqrt{(49 + 15)/64} = 1$, as required.

\textbf{\color{brand}Why the other options are incorrect.}
\begin{itemize}[leftmargin=2.4em, itemsep=1pt, topsep=2pt]
\item[\textbf{(A)}] These real numbers are not of modulus $1$; the factor $i$ was lost.
\item[\textbf{(B)}] These are $e^{\pm i\theta}$, using $\theta$ instead of the rotation angle $2\theta$.
\item[\textbf{(C)}] $\pm1$ are the eigenvalues of a reflection, or of a rotation by $0$ or $\pi$ only.
\end{itemize}
\textbf{\color{brand}Key takeaway.} Eigenvalues of a rotation lie on the unit circle at the rotation angle: quantum counting reads $\theta$ from them.
\end{solbox}
\qsep

\needspace{13\baselineskip}
\begin{qitem}{5}
\qmeta{Single-correct MCQ}{Powers of a complex phase}
Let $z = \dfrac{1 + i}{\sqrt2}$. The value of $z^6$ is
\par\smallskip\noindent\begin{tabularx}{\linewidth}{@{}*{4}{L}@{}}
\correctlabel{A}\,$-i$ & \optlabel{B}\,$-1$ & \optlabel{C}\,$1$ & \optlabel{D}\,$i$
\end{tabularx}
\end{qitem}
\begin{solbox}{A}
\textbf{\color{brand}Concepts required.} Polar form $z = |z|e^{i\varphi}$, and $e^{i\alpha}e^{i\beta} = e^{i(\alpha+\beta)}$.

\textbf{\color{brand}Relevant formulae.} $z = |z|e^{i\varphi}$ with modulus $|z|$ and angle $\varphi$; $\;z^m = |z|^me^{im\varphi}$; $\;e^{i\varphi} = \cos\varphi + i\sin\varphi$.

\textbf{\color{brand}Step-by-step explanation.} Write $z$ in polar form: $|z| = \sqrt{\frac12 + \frac12} = 1$ and its angle is $45^\circ$, so $z = e^{i\pi/4}$. Powers multiply the angle:
\[ z^6 = e^{6i\pi/4} = e^{3i\pi/2} = \cos 270^\circ + i\sin 270^\circ = -i. \]
Since $|z| = 1$, multiplying an amplitude by $z$ changes only its phase and never its probability.

\textbf{\color{brand}Why the other options are incorrect.}
\begin{itemize}[leftmargin=2.4em, itemsep=1pt, topsep=2pt]
\item[\textbf{(B)}] $-1 = e^{i\pi}$ is $z^4$.
\item[\textbf{(C)}] $1$ would need a whole number of turns; $z^8 = 1$, but $z^6$ has turned only $270^\circ$.
\item[\textbf{(D)}] $i = e^{i\pi/2}$ corresponds to $z^2$ or a sign slip in the angle $270^\circ$.
\end{itemize}
\textbf{\color{brand}Key takeaway.} Convert to polar form first: powers of a phase simply add angles.
\end{solbox}
\qsep

\needspace{13\baselineskip}
\begin{qitem}{6}
\qmeta{Single-correct MCQ}{Optimal integer $k$; why $-\frac12$ matters}
Let $\sin\theta = \tfrac{1}{10}$, $0 < \theta < \tfrac\pi2$; numerically $\theta = 0.10017$. Among the integers $0 \le k \le 10$, the one that makes
$\sin^2\big((2k+1)\theta\big)$ largest is
\par\smallskip\noindent\begin{tabularx}{\linewidth}{@{}*{4}{L}@{}}
\optlabel{A}\,$8$ & \optlabel{B}\,$10$ & \optlabel{C}\,$5$ & \correctlabel{D}\,$7$
\end{tabularx}
\end{qitem}
\begin{solbox}{D}
\textbf{\color{brand}Concepts required.} $\sin^2((2k+1)\theta)$ is largest when $(2k+1)\theta$ is closest to $\frac\pi2$, that is for $k$ nearest to $\frac{\pi}{4\theta} - \frac12$.

\textbf{\color{brand}Relevant formulae.} $\sin^2\alpha$ is largest at $\alpha = \frac\pi2$; so $(2k+1)\theta = \frac\pi2$ gives the real value $k = \frac{\pi}{4\theta} - \frac12$, to be rounded to an integer.

\textbf{\color{brand}Step-by-step explanation.} We want $(2k+1)\theta$ as close as possible to $\frac\pi2$, i.e.\ $k$ closest to $\frac{\pi}{4\theta} - \frac12$. Here
\[ \frac{\pi}{4\theta} - \frac12 = \frac{3.1416}{0.40067} - 0.5 = 7.841 - 0.5 = 7.341, \]
so $k = 7$: $15\theta = 1.5025$ rad is $0.068$ below $\frac\pi2$, while $17\theta = 1.7028$ rad is $0.132$ above it. Then
$\sin^2 15\theta = 0.9953$ and $\sin^2 17\theta = 0.9827$.
The tempting estimate $\frac\pi4\cdot 10 = 7.85$ rounds to $8$, which is wrong: it omits the $-\frac12$.

\textbf{\color{brand}Why the other options are incorrect.}
\begin{itemize}[leftmargin=2.4em, itemsep=1pt, topsep=2pt]
\item[\textbf{(A)}] $k = 8$ comes from rounding $\frac\pi4\sqrt N = 7.85$; it overshoots, and $\sin^2 17\theta = 0.983 < 0.995$.
\item[\textbf{(B)}] $k = 10 = \sqrt N$ omits the factor $\frac\pi4$, giving $21\theta = 2.10$ rad and $\sin^2 \approx 0.74$.
\item[\textbf{(C)}] $k = 5$ stops too soon: $11\theta = 1.10$ rad and $\sin^2 \approx 0.79$.
\end{itemize}
\textbf{\color{brand}Key takeaway.} Use $\frac{\pi}{4\theta} - \frac12$, not $\frac\pi4\sqrt N$, when the exact integer matters.
\end{solbox}
\qsep

\needspace{13\baselineskip}
\begin{qitem}{7}
\qmeta{Single-correct MCQ}{Inversion about the mean}
The vector $v \in \mathbb{R}^{8}$ has seven entries equal to $+1$ and one entry equal to $-1$. Each entry $v_j$ is replaced by $2\bar v - v_j$,
where $\bar v$ is the mean of the entries. The new value of the entry that was $-1$ is
\par\smallskip\noindent\begin{tabularx}{\linewidth}{@{}*{4}{L}@{}}
\correctlabel{A}\,$\dfrac52$ & \optlabel{B}\,$\dfrac32$ & \optlabel{C}\,$\dfrac74$ & \optlabel{D}\,$\dfrac12$
\end{tabularx}
\end{qitem}
\begin{solbox}{A}
\textbf{\color{brand}Concepts required.} Inversion about the mean: $v_j \mapsto 2\bar v - v_j$. It preserves the length of $v$ and leaves the mean unchanged.

\textbf{\color{brand}Relevant formulae.} Mean $\bar v = \frac1N\sum_jv_j$ of the $N$ entries; inversion about the mean: $v_j \mapsto 2\bar v - v_j$.

\textbf{\color{brand}Step-by-step explanation.} The mean is $\bar v = \frac{7 - 1}{8} = \frac34$. The special entry becomes $2\cdot\frac34 - (-1) = \frac52$; each other entry becomes
$\frac32 - 1 = \frac12$. \emph{Check:} the length is preserved, $7\cdot\frac14 + \frac{25}{4} = 8 = 7 + 1$.
Inversion about the mean reflects every entry through the mean, so an entry far below the mean ends up far above it.

\textbf{\color{brand}Why the other options are incorrect.}
\begin{itemize}[leftmargin=2.4em, itemsep=1pt, topsep=2pt]
\item[\textbf{(B)}] $\frac32 = 2\bar v$ forgets the term $-v_j = +1$.
\item[\textbf{(C)}] $\frac74$ comes from using the mean $\frac78$ of a wrong sum and is not consistent with length preservation.
\item[\textbf{(D)}] $\frac12$ is the new value of the \emph{other} entries.
\end{itemize}
\textbf{\color{brand}Key takeaway.} An entry a distance $d$ below the mean ends a distance $d$ above it. Always check that the length is preserved.
\end{solbox}
\qsep

\needspace{13\baselineskip}
\begin{qitem}{8}
\qmeta{Single-correct MCQ}{Classical search without replacement}
A box holds $31$ cards of which exactly $3$ are marked. Cards are drawn one at a time at random without replacement until the first
marked card appears. The expected number of draws is
\par\smallskip\noindent\begin{tabularx}{\linewidth}{@{}*{4}{L}@{}}
\optlabel{A}\,$\dfrac{31}{4}$ & \optlabel{B}\,$\dfrac{31}{3}$ & \correctlabel{C}\,$8$ & \optlabel{D}\,$16$
\end{tabularx}
\end{qitem}
\begin{solbox}{C}
\textbf{\color{brand}Concepts required.} Without replacement, the expected position of the first of $M$ marked cards among $N$ is $\frac{N+1}{M+1}$ (symmetry of the gaps).

\textbf{\color{brand}Relevant formulae.} Expected draws without replacement: $\frac{N+1}{M+1}$, with $N$ items of which $M$ are marked. With replacement: $\frac NM$.

\textbf{\color{brand}Step-by-step explanation.} The $M = 3$ marked cards cut the random order of the $N - M = 28$ unmarked cards into $M + 1 = 4$ gaps, and by symmetry each gap
holds on average $\frac{28}{4} = 7$ unmarked cards. The first marked card therefore comes after $7$ unmarked ones on average, on draw number $8$.
In general the expected number of draws is $\frac{N+1}{M+1} = \frac{32}{4} = 8$. The value $\frac{31}{3}$ would be right
for drawing \emph{with} replacement.

\textbf{\color{brand}Why the other options are incorrect.}
\begin{itemize}[leftmargin=2.4em, itemsep=1pt, topsep=2pt]
\item[\textbf{(A)}] $\frac{31}{4} = \frac{N}{M+1}$ uses $N$ where $N + 1$ is needed.
\item[\textbf{(B)}] $\frac{31}{3} = N/M$ is the expectation \emph{with} replacement (geometric distribution).
\item[\textbf{(D)}] $16 = \frac{N+1}{2}$ is the formula for one marked card.
\end{itemize}
\textbf{\color{brand}Key takeaway.} With $M = 1$ the formula gives the classical average $\frac{N+1}{2}$ that Grover is compared with.
\end{solbox}
\qsep

\needspace{13\baselineskip}
\begin{qitem}{9}
\qmeta{Single-correct MCQ}{Binary vs. linear search}
A key, which may be absent, is looked up in a list of $1023$ distinct entries using classical comparisons. The worst-case numbers of comparisons are, for binary search on the \emph{sorted} list and for linear search on the \emph{unsorted} list, respectively,
\par\smallskip\noindent\begin{tabularx}{\linewidth}{@{}*{4}{L}@{}}
\correctlabel{A}\,$10$ and $1023$ & \optlabel{B}\,$1023$ and $1023$ & \optlabel{C}\,$10$ and $512$ & \optlabel{D}\,$32$ and $1023$
\end{tabularx}
\end{qitem}
\begin{solbox}{A}
\textbf{\color{brand}Concepts required.} Binary search halves a sorted range at each step, so it needs about $\log_2 N$ comparisons. Without order, nothing can be excluded without looking at it.

\textbf{\color{brand}Relevant formulae.} Binary search on $N = 2^b - 1$ sorted entries: $b$ comparisons at most. Linear search on $N$ unsorted entries: $N$ comparisons at most.

\textbf{\color{brand}Step-by-step explanation.} Each comparison in binary search halves the remaining range. $1023 = 2^{10} - 1$ entries are exhausted after exactly $10$ comparisons. An unsorted list has no structure to exploit: if the key is absent or in the last place, all $1023$ entries must be checked.
Grover's algorithm addresses the second, \emph{unstructured}, case. ($512$ is about the \emph{average} cost of a linear search when the key is present.)

\textbf{\color{brand}Why the other options are incorrect.}
\begin{itemize}[leftmargin=2.4em, itemsep=1pt, topsep=2pt]
\item[\textbf{(B)}] $1023$ for the sorted list ignores the order, which is exactly what binary search exploits.
\item[\textbf{(C)}] $512$ is about the average cost of a linear search when the key is present, not the worst case.
\item[\textbf{(D)}] $32 \approx \sqrt{1023}$ is a square-root cost, which no classical comparison search on a sorted list needs.
\end{itemize}
\textbf{\color{brand}Key takeaway.} Structure gives $\log_2 N$. Without structure the classical cost is $N$; this is the setting of Grover's algorithm.
\end{solbox}
\qsep

\needspace{13\baselineskip}
\begin{qitem}{10}
\qmeta{Single-correct MCQ}{Counting solutions of a Boolean function}
Let $f(x_1, x_2, x_3, x_4) = (x_1 \lor x_2) \land (x_3 \oplus x_4)$ on $\{0, 1\}^4$. The fraction of inputs with $f = 1$ is
\par\smallskip\noindent\begin{tabularx}{\linewidth}{@{}*{4}{L}@{}}
\optlabel{A}\,$\dfrac12$ & \correctlabel{B}\,$\dfrac38$ & \optlabel{C}\,$\dfrac14$ & \optlabel{D}\,$\dfrac58$
\end{tabularx}
\end{qitem}
\begin{solbox}{B}
\textbf{\color{brand}Concepts required.} Count the solutions of a Boolean function. If factors share no variables, the counts multiply. $\sin^2\theta = M/N$.

\textbf{\color{brand}Relevant formulae.} Fraction of solutions $= M/N$, with $M$ inputs satisfying $f$ out of $N = 2^n$. Conditions on disjoint variables: counts multiply.

\textbf{\color{brand}Step-by-step explanation.} The two factors involve different variables, so the counts multiply. $x_1 \lor x_2 = 1$ for $3$ of the $4$ pairs $(x_1, x_2)$;
$x_3 \oplus x_4 = 1$ for $2$ of the $4$ pairs $(x_3, x_4)$. So $M = 3\times2 = 6$ of $N = 16$ inputs, a fraction $\frac{6}{16} = \frac38$.
This fraction is the quantity $\sin^2\theta$ in a Grover search for the solutions of $f$.

\textbf{\color{brand}Why the other options are incorrect.}
\begin{itemize}[leftmargin=2.4em, itemsep=1pt, topsep=2pt]
\item[\textbf{(A)}] $\frac12$ is the fraction for $x_3 \oplus x_4$ alone.
\item[\textbf{(C)}] $\frac14 = \frac{4}{16}$ counts only two solutions of $x_1 \lor x_2$ instead of three.
\item[\textbf{(D)}] $\frac58$ averages the fractions $\frac34$ and $\frac12$ instead of multiplying them.
\end{itemize}
\textbf{\color{brand}Key takeaway.} Find $M$ by factoring the Boolean expression into independent parts.
\end{solbox}
\qsep

\needspace{13\baselineskip}
\begin{qitem}{11}
\qmeta{Single-correct MCQ}{The signed sum $N - 2M$}
A Boolean function $f$ on $4$ bits takes the value $1$ on exactly $3$ inputs. The value of $\displaystyle\sum_{x\in\{0,1\}^4}(-1)^{f(x)}$ is
\par\smallskip\noindent\begin{tabularx}{\linewidth}{@{}*{4}{L}@{}}
\correctlabel{A}\,$10$ & \optlabel{B}\,$7$ & \optlabel{C}\,$6$ & \optlabel{D}\,$13$
\end{tabularx}
\end{qitem}
\begin{solbox}{A}
\textbf{\color{brand}Concepts required.} $(-1)^{f(x)}$ is $+1$ on the $N - M$ unmarked inputs and $-1$ on the $M$ marked ones, so the sum is $N - 2M$.

\textbf{\color{brand}Relevant formulae.} $\sum_x(-1)^{f(x)} = (N - M) - M = N - 2M$, with $M$ inputs having $f(x) = 1$.

\textbf{\color{brand}Step-by-step explanation.} Each of the $N - M = 13$ inputs with $f = 0$ contributes $+1$, and each of the $M = 3$ inputs with $f = 1$ contributes $-1$:
\[ \sum_x (-1)^{f(x)} = (N - M) - M = N - 2M = 16 - 6 = 10. \]
Divided by $N$ this is the factor $1 - \frac{2M}{N}$ that appears in the mean amplitude after the oracle.

\textbf{\color{brand}Why the other options are incorrect.}
\begin{itemize}[leftmargin=2.4em, itemsep=1pt, topsep=2pt]
\item[\textbf{(B)}] $7$ is $N/2 - 1$ and does not follow from the counts.
\item[\textbf{(C)}] $6 = 2M$ is the amount lost from $N$, not the sum.
\item[\textbf{(D)}] $13 = N - M$ counts the unmarked inputs and ignores the $-1$ from the marked ones.
\end{itemize}
\textbf{\color{brand}Key takeaway.} The mean amplitude after the oracle is proportional to $N - 2M$.
\end{solbox}
\qsep

\needspace{13\baselineskip}
\begin{qitem}{12}
\qmeta{Single-correct MCQ}{Normalising a state vector}
The vector $\ket\psi = c\,\big(\ket{00} + i\ket{01} + (1 + i)\ket{10} + 2\ket{11}\big)$ with $c > 0$ is normalised. The value of $c$ is
\par\smallskip\noindent\begin{tabularx}{\linewidth}{@{}*{4}{L}@{}}
\optlabel{A}\,$\dfrac15$ & \correctlabel{B}\,$\dfrac{1}{2\sqrt2}$ & \optlabel{C}\,$\dfrac18$ & \optlabel{D}\,$\dfrac{1}{\sqrt7}$
\end{tabularx}
\end{qitem}
\begin{solbox}{B}
\textbf{\color{brand}Concepts required.} Normalisation: $\sum_x|c_x|^2 = 1$, with $|a + ib|^2 = a^2 + b^2$.

\textbf{\color{brand}Relevant formulae.} Normalisation: $\sum_x|c_x|^2 = 1$ for amplitudes $c_x$; $\;|a + ib|^2 = a^2 + b^2$ for real $a, b$.

\textbf{\color{brand}Step-by-step explanation.} Normalisation means $\sum_x|c_x|^2 = 1$, where each term is a squared \emph{modulus}: $|1|^2 = 1$, $|i|^2 = 1$, $|1 + i|^2 = 1^2 + 1^2 = 2$, $|2|^2 = 4$. So
\[ c^2(1 + 1 + 2 + 4) = 8c^2 = 1 \;\Rightarrow\; c = \frac{1}{\sqrt8} = \frac{1}{2\sqrt2}. \]

\textbf{\color{brand}Why the other options are incorrect.}
\begin{itemize}[leftmargin=2.4em, itemsep=1pt, topsep=2pt]
\item[\textbf{(A)}] $\frac15$ adds non-squared sizes such as $1 + 1 + 1 + 2$ instead of squared moduli.
\item[\textbf{(C)}] $\frac18$ forgets the square root in $c^2 = \frac18$.
\item[\textbf{(D)}] $\frac{1}{\sqrt7}$ takes $|1 + i|^2 = 1$ instead of $2$.
\end{itemize}
\textbf{\color{brand}Key takeaway.} Use squared moduli, never squares of the complex numbers themselves.
\end{solbox}
\qsep

\needspace{13\baselineskip}
\begin{qitem}{13}
\qmeta{Single-correct MCQ}{Square-root scaling}
Let $g(n) = \tfrac\pi4\sqrt{2^n}$. Then $g(n + 6)/g(n)$ equals
\par\smallskip\noindent\begin{tabularx}{\linewidth}{@{}*{4}{L}@{}}
\optlabel{A}\,$6$ & \correctlabel{B}\,$8$ & \optlabel{C}\,$\sqrt6$ & \optlabel{D}\,$64$
\end{tabularx}
\end{qitem}
\begin{solbox}{B}
\textbf{\color{brand}Concepts required.} $\sqrt{2^{n}} = 2^{n/2}$, so adding $a$ bits multiplies $\sqrt N$ by $2^{a/2}$.

\textbf{\color{brand}Relevant formulae.} $\sqrt{2^{n}} = 2^{n/2}$, so $\frac{\sqrt{2^{n+a}}}{\sqrt{2^n}} = 2^{a/2}$ when $a$ bits are added.

\textbf{\color{brand}Step-by-step explanation.} $\frac{g(n+6)}{g(n)} = \sqrt{\frac{2^{n+6}}{2^n}} = \sqrt{2^6} = 2^3 = 8$. Six extra bits make the search space $64$ times larger but the
number of Grover iterations only $8$ times larger.

\textbf{\color{brand}Why the other options are incorrect.}
\begin{itemize}[leftmargin=2.4em, itemsep=1pt, topsep=2pt]
\item[\textbf{(A)}] $6$ treats the cost as linear in the number of bits.
\item[\textbf{(C)}] $\sqrt6$ takes the root of the \emph{number of added bits} rather than of $2^6$.
\item[\textbf{(D)}] $64 = 2^6$ is the growth of $N$, which is the classical scaling.
\end{itemize}
\textbf{\color{brand}Key takeaway.} Each extra qubit multiplies the iteration count by $\sqrt2$.
\end{solbox}
\qsep

\needspace{13\baselineskip}
\begin{qitem}{14}
\qmeta{Single-correct MCQ}{Complex inner product}
Let $u = \frac{1}{\sqrt2}(1, i, 0)^{T}$ and $v = \frac{1}{\sqrt3}(1, 1, 1)^{T}$. The inner product $\braket{u|v}$ is
\par\smallskip\noindent\begin{tabularx}{\linewidth}{@{}*{4}{L}@{}}
\optlabel{A}\,$\dfrac{1 + i}{\sqrt6}$ & \optlabel{B}\,$\dfrac{2}{\sqrt6}$ & \correctlabel{C}\,$\dfrac{1 - i}{\sqrt6}$ & \optlabel{D}\,$0$
\end{tabularx}
\end{qitem}
\begin{solbox}{C}
\textbf{\color{brand}Concepts required.} $\braket{u|v} = \sum_k\overline{u_k}v_k$: conjugate the first vector.

\textbf{\color{brand}Relevant formulae.} $\braket{u|v} = \sum_k\overline{u_k}v_k$, where the bar means the complex conjugate.

\textbf{\color{brand}Step-by-step explanation.} $\braket{u|v} = \sum_k \overline{u_k}v_k$: the \emph{first} vector is conjugated. Thus
\[ \braket{u|v} = \frac{1}{\sqrt2}\cdot\frac{1}{\sqrt3}\big(\bar1\cdot1 + \bar i\cdot1 + 0\big) = \frac{1 - i}{\sqrt6}. \]
Forgetting the conjugate gives $\frac{1 + i}{\sqrt6}$, which is $\braket{v|u}$. The probability $|\braket{u|v}|^2 = \frac{2}{6} = \frac13$ is the same either way.

\textbf{\color{brand}Why the other options are incorrect.}
\begin{itemize}[leftmargin=2.4em, itemsep=1pt, topsep=2pt]
\item[\textbf{(A)}] $\frac{1 + i}{\sqrt6}$ forgets the conjugate on $u$; it is $\braket{v|u}$.
\item[\textbf{(B)}] $\frac{2}{\sqrt6}$ uses the moduli of the components and loses the phase $i$.
\item[\textbf{(D)}] $0$ would mean orthogonal vectors; the first component alone gives a non-cancelling term.
\end{itemize}
\textbf{\color{brand}Key takeaway.} Swapping the two vectors conjugates the inner product: $\braket{v|u} = \overline{\braket{u|v}}$.
\end{solbox}
\qsep

\needspace{13\baselineskip}
\begin{qitem}{15}
\qmeta{Matching lists}{$\sin^2 3\theta$ as a polynomial}
Match the value of $\sin^2\theta$ (List-I) with the value of $\sin^2 3\theta$ (List-II).
\par\smallskip\noindent{\renewcommand{\arraystretch}{1.35}%
\begin{tabularx}{\linewidth}{@{}c>{\raggedright\arraybackslash}p{0.33\linewidth}@{\hspace{1.2em}}cL@{}}
\toprule
\multicolumn{2}{@{}l}{\textbf{List-I}} & \multicolumn{2}{l@{}}{\textbf{List-II}} \\ \midrule
\textbf{(P)} & $\sin^2\theta = \frac18$ & \textbf{(1)} & $1$ \\
\textbf{(Q)} & $\sin^2\theta = \frac14$ & \textbf{(2)} & $0$ \\
\textbf{(R)} & $\sin^2\theta = \frac12$ & \textbf{(3)} & $\frac{25}{32}$ \\
\textbf{(S)} & $\sin^2\theta = \frac34$ & \textbf{(4)} & $\frac12$ \\
\bottomrule
\end{tabularx}}
\par\smallskip\noindent\begin{tabularx}{\linewidth}{@{}*{2}{L}@{}}
\optlabel{A}\,P$\to$4, Q$\to$1, R$\to$3, S$\to$2 & \optlabel{B}\,P$\to$3, Q$\to$1, R$\to$2, S$\to$4 \\[3pt]
\correctlabel{C}\,P$\to$3, Q$\to$1, R$\to$4, S$\to$2 & \optlabel{D}\,P$\to$1, Q$\to$3, R$\to$4, S$\to$2
\end{tabularx}
\end{qitem}
\begin{solbox}{C}
\textbf{\color{brand}Concepts required.} $\sin^2 3\theta = s(3 - 4s)^2$ with $s = \sin^2\theta$.

\textbf{\color{brand}Relevant formulae.} $\sin^2 3\theta = s(3 - 4s)^2$ with $s = \sin^2\theta$.

\textbf{\color{brand}Step-by-step explanation.} Since $\sin 3\theta = \sin\theta\,(3 - 4\sin^2\theta)$, with $s = \sin^2\theta$ we get $\sin^2 3\theta = s(3 - 4s)^2$.
$s = \frac18$: $\frac18\cdot\frac{25}{4} = \frac{25}{32}$ (3). $\;s = \frac14$: $\frac14\cdot 4 = 1$ (1). $\;s = \frac12$: $\frac12\cdot 1 = \frac12$ (4).
$\;s = \frac34$: $\frac34\cdot 0 = 0$ (2). Hence P$\to$3, Q$\to$1, R$\to$4, S$\to$2.

\textbf{\color{brand}Why the other options are incorrect.}
\begin{itemize}[leftmargin=2.4em, itemsep=1pt, topsep=2pt]
\item[\textbf{(A)}] This gives $s = \frac18$ the value $\frac12$, but $\frac18\cdot\frac{25}{4} = \frac{25}{32}$.
\item[\textbf{(B)}] This exchanges R and S; for $s = \frac34$ we have $3 - 4s = 0$, so the value is $0$.
\item[\textbf{(D)}] This exchanges P and Q; for $s = \frac14$ the value is $1$, not $\frac{25}{32}$.
\end{itemize}
\textbf{\color{brand}Key takeaway.} $s = \frac14$ is perfect, $s = \frac12$ is flat and $s = \frac34$ is the worst case for one iteration.
\end{solbox}
\qsep

\needspace{13\baselineskip}
\begin{qitem}{16}
\qmeta{Matching lists}{Traces of projectors and reflections}
Let $\ket{u}$ be a unit vector in $\mathbb{C}^{N}$, with $N \ge 3$. Match each matrix in List-I with its trace in List-II.
\par\smallskip\noindent{\renewcommand{\arraystretch}{1.35}%
\begin{tabularx}{\linewidth}{@{}c>{\raggedright\arraybackslash}p{0.33\linewidth}@{\hspace{1.2em}}cL@{}}
\toprule
\multicolumn{2}{@{}l}{\textbf{List-I}} & \multicolumn{2}{l@{}}{\textbf{List-II}} \\ \midrule
\textbf{(P)} & $2\ket{u}\!\bra{u} - I$ & \textbf{(1)} & $N$ \\
\textbf{(Q)} & $I - 2\ket{u}\!\bra{u}$ & \textbf{(2)} & $N - 2$ \\
\textbf{(R)} & $\ket{u}\!\bra{u}$ & \textbf{(3)} & $1$ \\
\textbf{(S)} & $I$ & \textbf{(4)} & $2 - N$ \\
\bottomrule
\end{tabularx}}
\par\smallskip\noindent\begin{tabularx}{\linewidth}{@{}*{2}{L}@{}}
\optlabel{A}\,P$\to$3, Q$\to$2, R$\to$4, S$\to$1 & \optlabel{B}\,P$\to$4, Q$\to$2, R$\to$1, S$\to$3 \\[3pt]
\optlabel{C}\,P$\to$2, Q$\to$4, R$\to$3, S$\to$1 & \correctlabel{D}\,P$\to$4, Q$\to$2, R$\to$3, S$\to$1
\end{tabularx}
\end{qitem}
\begin{solbox}{D}
\textbf{\color{brand}Concepts required.} $\operatorname{tr}\ket{u}\!\bra{u} = \braket{u|u} = 1$ and trace is linear.

\textbf{\color{brand}Relevant formulae.} $\operatorname{tr}\ket{u}\!\bra{u} = \braket{u|u} = 1$; $\;\operatorname{tr}I = N$ in dimension $N$; trace is linear.

\textbf{\color{brand}Step-by-step explanation.} The key fact is $\operatorname{tr}\ket{u}\!\bra{u} = \braket{u|u} = 1$, and $\operatorname{tr} I = N$. So
$\operatorname{tr}(2\ket{u}\!\bra{u} - I) = 2 - N$ (4); $\;\operatorname{tr}(I - 2\ket{u}\!\bra{u}) = N - 2$ (2); $\;\operatorname{tr}\ket{u}\!\bra{u} = 1$ (3);
$\;\operatorname{tr} I = N$ (1). Hence P$\to$4, Q$\to$2, R$\to$3, S$\to$1.

\textbf{\color{brand}Why the other options are incorrect.}
\begin{itemize}[leftmargin=2.4em, itemsep=1pt, topsep=2pt]
\item[\textbf{(A)}] This gives $2\ket{u}\!\bra{u} - I$ trace $1$ and the projector trace $2 - N$, which ignores linearity.
\item[\textbf{(B)}] This swaps R and S: the projector has trace $1$, the identity has trace $N$.
\item[\textbf{(C)}] This swaps P and Q: the sign in front of the identity decides which of $N - 2$, $2 - N$ applies.
\end{itemize}
\textbf{\color{brand}Key takeaway.} Trace identifies eigenvalue multiplicities: $D$ has one $+1$ and $N - 1$ eigenvalues $-1$.
\end{solbox}
\qsep

\needspace{13\baselineskip}
\begin{qitem}{17}
\qmeta{Assertion--Reason}{Rounding rule vs.\ $\arcsin x > x$}
\textbf{Assertion (A):} Rounding $\frac\pi4\sqrt N$ to the nearest integer does not always give the integer $k$ for which
$(2k+1)\theta$ is closest to $\frac\pi2$, where $\sin\theta = 1/\sqrt N$.\\[1mm]
\textbf{Reason (R):} For $0 < x < 1$ we have $\arcsin x < x$.
\begin{itemize}[leftmargin=2.6em, labelsep=0.4em, itemsep=2pt, topsep=3pt]
\item[\optlabel{A}] Both Assertion (A) and Reason (R) are true, and R is the correct explanation of A
\item[\optlabel{B}] Both A and R are true, but R is \emph{not} the correct explanation of A
\item[\correctlabel{C}] A is true, but R is false
\item[\optlabel{D}] A is false, but R is true
\end{itemize}
\end{qitem}
\begin{solbox}{C}
\textbf{\color{brand}Concepts required.} Decide A and R independently. $\arcsin x = x + x^3/6 + \dots$, so it exceeds $x$ for $x > 0$.

\textbf{\color{brand}Relevant formulae.} $k^\ast = \frac{\pi}{4\theta} - \frac12$; $\;\arcsin x = x + \frac{x^3}{6} + \dots$ for $|x| < 1$.

\textbf{\color{brand}Step-by-step explanation.} \textbf{A is true.} For $N = 100$, $\frac\pi4\sqrt N = 7.85$ rounds to $8$, but the true optimum is $k = 7$ (see Q.6).
\textbf{R is false:} the inequality goes the other way. $\arcsin x = x + \frac{x^3}{6} + \dots > x$, for example $\arcsin 0.1 = 0.10017$.
The actual reason for A is the omitted $-\frac12$ in $k = \frac{\pi}{4\theta} - \frac12$. Answer: (C).

\textbf{\color{brand}Why the other options are incorrect.}
\begin{itemize}[leftmargin=2.4em, itemsep=1pt, topsep=2pt]
\item[\textbf{(A)}] R is false, so it cannot explain A.
\item[\textbf{(B)}] This needs R to be true, but the inequality in R points the wrong way.
\item[\textbf{(D)}] A is true: $N = 100$ is a counterexample to the rounding rule, and R is false.
\end{itemize}
\textbf{\color{brand}Key takeaway.} An Assertion can be true while the Reason offered for it is false: test both with a number.
\end{solbox}
\qsep

\needspace{13\baselineskip}
\begin{qitem}{18}
\qmeta{Statement I / II}{Spectrum of a reflection}
Let $\ket{u}$ be a unit vector in $\mathbb{C}^{N}$, with $N \ge 2$, and let $D = 2\ket{u}\!\bra{u} - I$.\\[1mm]
\textbf{Statement I:} $D^2 = I$.\\[1mm]
\textbf{Statement II:} $D$ has the eigenvalue $+1$ with multiplicity $1$ and the eigenvalue $-1$ with multiplicity $N - 1$.
\begin{itemize}[leftmargin=2.6em, labelsep=0.4em, itemsep=2pt, topsep=3pt]
\item[\correctlabel{A}] Both Statement I and Statement II are true
\item[\optlabel{B}] Both Statement I and Statement II are false
\item[\optlabel{C}] Statement I is true, but Statement II is false
\item[\optlabel{D}] Statement I is false, but Statement II is true
\end{itemize}
\end{qitem}
\begin{solbox}{A}
\textbf{\color{brand}Concepts required.} $D$ is a reflection: $D^2 = I$, it fixes $\ket{u}$, and it negates the orthogonal complement.

\textbf{\color{brand}Relevant formulae.} $D = 2\ket{u}\!\bra{u} - I$ with $\braket{u|u} = 1$; $\;D\ket{u} = \ket{u}$; $\;D\ket{w} = -\ket{w}$ if $\braket{u|w} = 0$.

\textbf{\color{brand}Step-by-step explanation.} \textbf{I is true:} $D^2 = 4\ket{u}\braket{u|u}\bra{u} - 4\ket{u}\!\bra{u} + I = I$, because $\braket{u|u} = 1$.
\textbf{II is true:} $D\ket{u} = 2\ket{u} - \ket{u} = \ket{u}$, and for any $\ket{w}$ orthogonal to $\ket{u}$ we have $D\ket{w} = -\ket{w}$; such vectors span
an $(N-1)$-dimensional space. \emph{Check:} the trace is $1 - (N-1) = 2 - N$, in agreement with Q.16. Answer: (A).

\textbf{\color{brand}Why the other options are incorrect.}
\begin{itemize}[leftmargin=2.4em, itemsep=1pt, topsep=2pt]
\item[\textbf{(B)}] Both are true: $\braket{u|u} = 1$ gives I, and acting on $\ket{u}$ and on its complement gives II.
\item[\textbf{(C)}] II is also true, as the trace check $1 - (N-1) = 2 - N$ confirms.
\item[\textbf{(D)}] I is true by direct expansion.
\end{itemize}
\textbf{\color{brand}Key takeaway.} Every reflection is its own inverse and has only eigenvalues $\pm1$.
\end{solbox}
\qsep

\needspace{20\baselineskip}
\begin{passagebox}{Paragraph for Questions 19 and 20}
Define $p(k) = \sin^2\big((2k+1)\theta\big)$ for integers $k \ge 0$, with $\theta = \pi/18 = 10^\circ$. (This is the
shape of the Grover success curve, but the two questions below need only trigonometry.) Note that $(2k+1)\theta$ takes the
values $10^\circ, 30^\circ, 50^\circ, \dots$
\end{passagebox}
\needspace{13\baselineskip}
\begin{qitem}{19}
\qmeta{Paragraph based}{First maximum and period of $p(k)$}
The smallest $k$ with $p(k) = 1$, and the smallest positive period of the sequence $p(k)$, are
\begin{itemize}[leftmargin=2.6em, labelsep=0.4em, itemsep=2pt, topsep=3pt]
\item[\optlabel{A}] first maximum at $k = 4$; smallest period $18$
\item[\optlabel{B}] first maximum at $k = 5$; smallest period $9$
\item[\optlabel{C}] first maximum at $k = 9$; smallest period $18$
\item[\correctlabel{D}] first maximum at $k = 4$; smallest period $9$
\end{itemize}
\end{qitem}
\begin{solbox}{D}
\textbf{\color{brand}Concepts required.} $p(k) = 1$ requires an odd multiple of $90^\circ$. The period comes from $2\theta$ steps compared with $180^\circ$.

\textbf{\color{brand}Relevant formulae.} $p(k) = \sin^2((2k+1)\theta)$; $\sin^2$ has period $180^\circ$; moving $k$ by $q$ adds $2q\theta$ to the angle.

\textbf{\color{brand}Step-by-step explanation.} $p(k) = 1$ requires $(2k+1)\cdot 10^\circ = 90^\circ$ (the first odd multiple of $90^\circ$), i.e.\ $2k + 1 = 9$, so $k = 4$.
Moving from $k$ to $k + q$ adds $20^\circ q$ to the angle. $\sin^2$ has period $180^\circ$, so we need $20^\circ q = 180^\circ m$; the smallest
solution is $q = 9$. Hence period $9$ ($q = 18$ corresponds to the period $360^\circ$ of $\sin$ and is not the smallest).

\textbf{\color{brand}Why the other options are incorrect.}
\begin{itemize}[leftmargin=2.4em, itemsep=1pt, topsep=2pt]
\item[\textbf{(A)}] $18$ is a period but not the smallest: it corresponds to $360^\circ$.
\item[\textbf{(B)}] $k = 5$ gives $110^\circ$, already past the maximum.
\item[\textbf{(C)}] $k = 9$ gives $190^\circ$, which is near a minimum, and $18$ is not the smallest period.
\end{itemize}
\textbf{\color{brand}Key takeaway.} Convert to degrees and look at the list of angles: the sequence becomes easy to read.
\end{solbox}
\qsep

\needspace{13\baselineskip}
\begin{qitem}{20}
\qmeta{Paragraph based}{Tolerance window around the peak}
The number of integers $k$ with $0 \le k \le 17$ for which $p(k) \ge \tfrac34$ is
\par\smallskip\noindent\begin{tabularx}{\linewidth}{@{}*{4}{L}@{}}
\optlabel{A}\,$3$ & \correctlabel{B}\,$6$ & \optlabel{C}\,$9$ & \optlabel{D}\,$4$
\end{tabularx}
\end{qitem}
\begin{solbox}{B}
\textbf{\color{brand}Concepts required.} $\sin^2\alpha \ge \frac34 \iff |\sin\alpha| \ge \frac{\sqrt3}{2}$; $\alpha$ is within $30^\circ$ of an odd multiple of $90^\circ$.

\textbf{\color{brand}Relevant formulae.} $\sin^2\alpha \ge \frac34 \iff |\sin\alpha| \ge \frac{\sqrt3}{2}$, i.e.\ $\alpha$ is within $30^\circ$ of an odd multiple of $90^\circ$.

\textbf{\color{brand}Step-by-step explanation.} $\sin^2\alpha \ge \frac34$ means $|\sin\alpha| \ge \frac{\sqrt3}{2}$, i.e.\ $\alpha$ lies within $30^\circ$ of $90^\circ$ or of $270^\circ$. For $k = 0, \dots, 17$ the angle
$(2k+1)\cdot10^\circ$ runs over $10^\circ, 30^\circ, \dots, 350^\circ$. The values in $[60^\circ, 120^\circ]$ are $70^\circ, 90^\circ, 110^\circ$ ($k = 3, 4, 5$);
the values in $[240^\circ, 300^\circ]$ are $250^\circ, 270^\circ, 290^\circ$ ($k = 12, 13, 14$). That makes $6$ values of $k$.

\textbf{\color{brand}Why the other options are incorrect.}
\begin{itemize}[leftmargin=2.4em, itemsep=1pt, topsep=2pt]
\item[\textbf{(A)}] $3$ counts only one of the two peaks in the range $0 \le k \le 17$.
\item[\textbf{(C)}] $9$ is the period and not a count of angles in the window.
\item[\textbf{(D)}] $4$ counts two values at each peak; each peak has three: $70^\circ, 90^\circ, 110^\circ$ and $250^\circ, 270^\circ, 290^\circ$.
\end{itemize}
\textbf{\color{brand}Key takeaway.} A window of $\pm30^\circ$ around the peak corresponds to $P \ge \frac34$: there is some tolerance around $k_{\rm opt}$.
\end{solbox}
\qsep

\newpage

\begin{sectionbanner}SECTION B \enspace$\cdot$\enspace Quantum Computing Foundations \hfill Q.21--40\end{sectionbanner}

\needspace{13\baselineskip}
\begin{qitem}{21}
\qmeta{Single-correct MCQ}{Uniform state, partial outcome}
The state $H^{\otimes3}\ket{000}$ is measured in the computational basis. The probability that qubit 1 gives $1$ and qubit 3 gives $0$ is
\par\smallskip\noindent\begin{tabularx}{\linewidth}{@{}*{4}{L}@{}}
\optlabel{A}\,$\dfrac18$ & \correctlabel{B}\,$\dfrac14$ & \optlabel{C}\,$\dfrac{1}{\sqrt8}$ & \optlabel{D}\,$\dfrac12$
\end{tabularx}
\end{qitem}
\begin{solbox}{B}
\textbf{\color{brand}Concepts required.} $H^{\otimes n}\ket{0\dots0}$ is uniform and a product state, so its qubits are independent.

\textbf{\color{brand}Relevant formulae.} $H^{\otimes n}\ket{0\dots0} = \frac{1}{\sqrt N}\sum_x\ket{x}$ with $N = 2^n$; Born rule: $P(x) = |c_x|^2$; the probability of an event is the sum over its outcomes.

\textbf{\color{brand}Step-by-step explanation.} All $8$ outcomes have probability $\left|\frac{1}{\sqrt8}\right|^2 = \frac18$. The event ``$x_1 = 1$, $x_3 = 0$'' contains the two outcomes
$\ket{100}$ and $\ket{110}$, so $P = \frac28 = \frac14$. (Equivalently the qubits are independent: $\frac12\cdot\frac12$.)

\textbf{\color{brand}Why the other options are incorrect.}
\begin{itemize}[leftmargin=2.4em, itemsep=1pt, topsep=2pt]
\item[\textbf{(A)}] $\frac18$ is the probability of one specific full string; here qubit 2 is free.
\item[\textbf{(C)}] $\frac{1}{\sqrt8}$ is an amplitude, not a probability.
\item[\textbf{(D)}] $\frac12$ is the probability for one qubit only.
\end{itemize}
\textbf{\color{brand}Key takeaway.} Sum over all full strings consistent with a partial outcome.
\end{solbox}
\qsep

\needspace{13\baselineskip}
\begin{qitem}{22}
\qmeta{Single-correct MCQ}{Tensor-product amplitude}
Let $\ket\psi = \frac{1}{\sqrt2}\big(\ket0 + i\ket1\big) \otimes \frac{1}{\sqrt2}\big(\ket0 - \ket1\big)$. The amplitude of $\ket{10}$ in $\ket\psi$ is
\par\smallskip\noindent\begin{tabularx}{\linewidth}{@{}*{4}{L}@{}}
\correctlabel{A}\,$\dfrac{i}{2}$ & \optlabel{B}\,$-\dfrac{i}{2}$ & \optlabel{C}\,$\dfrac12$ & \optlabel{D}\,$-\dfrac12$
\end{tabularx}
\end{qitem}
\begin{solbox}{A}
\textbf{\color{brand}Concepts required.} Tensor product: the amplitude of $\ket{jk}$ is $a_jb_k$, with qubit 1 written first.

\textbf{\color{brand}Relevant formulae.} $(a_0\ket0 + a_1\ket1)\otimes(b_0\ket0 + b_1\ket1) = \sum_{j,k}a_jb_k\ket{jk}$.

\textbf{\color{brand}Step-by-step explanation.} Multiply out: $(a_0\ket0 + a_1\ket1)\otimes(b_0\ket0 + b_1\ket1)$ has the amplitude $a_jb_k$ on $\ket{jk}$. For $\ket{10}$: $a_1b_0 = \frac{i}{\sqrt2}\cdot\frac{1}{\sqrt2} = \frac{i}{2}$.
The whole state is $\frac12(\ket{00} - \ket{01} + i\ket{10} - i\ket{11})$.

\textbf{\color{brand}Why the other options are incorrect.}
\begin{itemize}[leftmargin=2.4em, itemsep=1pt, topsep=2pt]
\item[\textbf{(B)}] $-\frac i2$ is the amplitude of $\ket{11}$.
\item[\textbf{(C)}] $\frac12$ is the amplitude of $\ket{00}$.
\item[\textbf{(D)}] $-\frac12$ is the amplitude of $\ket{01}$.
\end{itemize}
\textbf{\color{brand}Key takeaway.} Read the first bit from the first factor and the second bit from the second factor.
\end{solbox}
\qsep

\needspace{13\baselineskip}
\begin{qitem}{23}
\qmeta{Single-correct MCQ}{Hadamard layer signs}
$H \otimes H$ is applied to $\frac12\big(\ket{00} + \ket{01} + \ket{10} - \ket{11}\big)$. The amplitude of $\ket{11}$ in the result is
\par\smallskip\noindent\begin{tabularx}{\linewidth}{@{}*{4}{L}@{}}
\optlabel{A}\,$\dfrac12$ & \optlabel{B}\,$0$ & \correctlabel{C}\,$-\dfrac12$ & \optlabel{D}\,$1$
\end{tabularx}
\end{qitem}
\begin{solbox}{C}
\textbf{\color{brand}Concepts required.} $H^{\otimes2}$ has entries $\frac12(-1)^{x\cdot y}$; each output amplitude is a signed sum of the input amplitudes.

\textbf{\color{brand}Relevant formulae.} $H^{\otimes n}\ket{x} = \frac{1}{\sqrt N}\sum_y(-1)^{x\cdot y}\ket{y}$, with $x\cdot y = \sum_jx_jy_j$.

\textbf{\color{brand}Step-by-step explanation.} Use $H^{\otimes2}\ket{x} = \frac12\sum_y(-1)^{x\cdot y}\ket{y}$. For $y = 11$ the signs $(-1)^{x\cdot y}$ for $x = 00, 01, 10, 11$ are $+, -, -, +$, so the amplitude is
\[ \frac12\cdot\frac12\big[(+1)(1) + (-1)(1) + (-1)(1) + (+1)(-1)\big] = \frac14(1 - 1 - 1 - 1) = -\frac12. \]
The full result is $\frac12(\ket{00} + \ket{01} + \ket{10} - \ket{11})$: this state is unchanged by $H \otimes H$.

\textbf{\color{brand}Why the other options are incorrect.}
\begin{itemize}[leftmargin=2.4em, itemsep=1pt, topsep=2pt]
\item[\textbf{(A)}] $\frac12$ has the wrong sign, from treating the entry for $x = 11$, $y = 11$ as negative.
\item[\textbf{(B)}] $0$ would need the signed sum to cancel; here it is $-2$ before the factor $\frac14$.
\item[\textbf{(D)}] $1$ is the magnitude when all four terms align, which they do not.
\end{itemize}
\textbf{\color{brand}Key takeaway.} Compute one output amplitude at a time with $(-1)^{x\cdot y}$.
\end{solbox}
\qsep

\needspace{13\baselineskip}
\begin{qitem}{24}
\qmeta{Single-correct MCQ}{Phase kickback with $\ket{-}$}
Let $U_f\ket{x}\ket{y} = \ket{x}\ket{y \oplus f(x)}$ and $\ket{-} = \frac{1}{\sqrt2}(\ket0 - \ket1)$. Then $U_f\ket{x}\ket{-}$ equals
\par\smallskip\noindent\begin{tabularx}{\linewidth}{@{}*{4}{L}@{}}
\optlabel{A}\,$\ket{x}\ket{f(x)}$ & \optlabel{B}\,$\ket{x}\ket{-}$ & \correctlabel{C}\,$(-1)^{f(x)}\ket{x}\ket{-}$ & \optlabel{D}\,$\ket{f(x)}\ket{-}$
\end{tabularx}
\end{qitem}
\begin{solbox}{C}
\textbf{\color{brand}Concepts required.} Phase kickback: an ancilla in $\ket{-}$ is an eigenvector of $X$ with eigenvalue $-1$.

\textbf{\color{brand}Relevant formulae.} $U_f\ket{x}\ket{y} = \ket{x}\ket{y\oplus f(x)}$, with $\oplus$ the XOR; $\;\ket{-} = \frac{1}{\sqrt2}(\ket0 - \ket1)$; $\;X\ket{-} = -\ket{-}$.

\textbf{\color{brand}Step-by-step explanation.} Expand the ancilla: $U_f\ket{x}\ket{-} = \frac{1}{\sqrt2}\ket{x}\big(\ket{f(x)} - \ket{1 \oplus f(x)}\big)$.
If $f(x) = 0$ the bracket is $\ket0 - \ket1$; if $f(x) = 1$ it is $\ket1 - \ket0 = -(\ket0 - \ket1)$. So the result is $(-1)^{f(x)}\ket{x}\ket{-}$. This is
\emph{phase kickback}: the ancilla is unchanged and the sign lands on the input register.

\textbf{\color{brand}Why the other options are incorrect.}
\begin{itemize}[leftmargin=2.4em, itemsep=1pt, topsep=2pt]
\item[\textbf{(A)}] The ancilla does not store $f(x)$ here; that would be the case for an ancilla in $\ket0$.
\item[\textbf{(B)}] This drops the sign, which is the entire point of the construction.
\item[\textbf{(D)}] The input register is not overwritten; it remains $\ket{x}$.
\end{itemize}
\textbf{\color{brand}Key takeaway.} Ancilla in $\ket{-}$ turns the bit oracle into the phase oracle.
\end{solbox}
\qsep

\needspace{13\baselineskip}
\begin{qitem}{25}
\qmeta{Single-correct MCQ}{No kickback with $\ket{+}$}
With $U_f$ as in the previous question and $\ket{+} = \frac{1}{\sqrt2}(\ket0 + \ket1)$, the state $U_f\ket{x}\ket{+}$ equals
\par\smallskip\noindent\begin{tabularx}{\linewidth}{@{}*{4}{L}@{}}
\optlabel{A}\,$\ket{x}\ket{-}$ & \optlabel{B}\,$(-1)^{f(x)}\ket{x}\ket{-}$ & \optlabel{C}\,$(-1)^{f(x)}\ket{x}\ket{+}$ & \correctlabel{D}\,$\ket{x}\ket{+}$
\end{tabularx}
\end{qitem}
\begin{solbox}{D}
\textbf{\color{brand}Concepts required.} $\ket{+}$ is an eigenvector of $X$ with eigenvalue $+1$, so a controlled bit flip leaves it unchanged.

\textbf{\color{brand}Relevant formulae.} $\ket{+} = \frac{1}{\sqrt2}(\ket0 + \ket1)$; $\;X\ket{+} = +\ket{+}$.

\textbf{\color{brand}Step-by-step explanation.} $U_f\ket{x}\ket{+} = \frac{1}{\sqrt2}\ket{x}\big(\ket{f(x)} + \ket{1 \oplus f(x)}\big) = \ket{x}\ket{+}$ in both cases, because $\ket0 + \ket1$ is symmetric under the bit flip
($X\ket{+} = +\ket{+}$). No phase is kicked back, so the ancilla must be prepared in $\ket{-}$, not $\ket{+}$.

\textbf{\color{brand}Why the other options are incorrect.}
\begin{itemize}[leftmargin=2.4em, itemsep=1pt, topsep=2pt]
\item[\textbf{(A)}] $U_f$ cannot change the ancilla from $\ket{+}$ to $\ket{-}$: flipping $\ket0 + \ket1$ gives the same state.
\item[\textbf{(B)}] This is the result for an ancilla starting in $\ket{-}$.
\item[\textbf{(C)}] The sign requires the eigenvalue $-1$, that is the ancilla $\ket{-}$.
\end{itemize}
\textbf{\color{brand}Key takeaway.} The ancilla's eigenvalue under $X$ decides the kicked-back sign.
\end{solbox}
\qsep

\needspace{13\baselineskip}
\begin{qitem}{26}
\qmeta{Single-correct MCQ}{$X$ conjugation selects the item}
Let CCZ be the three-qubit gate that multiplies only $\ket{111}$ by $-1$. Which basis state has its sign flipped by $X_2\,\mathrm{CCZ}\,X_2$, where $X_2$ is $X$ on qubit 2?
\par\smallskip\noindent\begin{tabularx}{\linewidth}{@{}*{4}{L}@{}}
\optlabel{A}\,$\ket{111}$ & \optlabel{B}\,$\ket{010}$ & \correctlabel{C}\,$\ket{101}$ & \optlabel{D}\,$\ket{000}$
\end{tabularx}
\end{qitem}
\begin{solbox}{C}
\textbf{\color{brand}Concepts required.} Conjugating by $X$ moves the affected basis state: $X_2\,\mathrm{CCZ}\,X_2 = I - 2\ket{101}\!\bra{101}$.

\textbf{\color{brand}Relevant formulae.} $\mathrm{CCZ} = I - 2\ket{111}\!\bra{111}$; $\;U(I - 2\ket{a}\!\bra{a})U^\dagger = I - 2U\ket{a}\!\bra{a}U^\dagger$; $X_2$ flips qubit 2.

\textbf{\color{brand}Step-by-step explanation.} The rightmost $X_2$ acts first. A state picks up the sign exactly when $X_2$ maps it to $\ket{111}$, which happens for $\ket{101}$. The last $X_2$ restores the
bit. So the sign flip sits on $\ket{101}$. This is the standard oracle recipe: apply $X$ wherever the marked string has a $0$.

\textbf{\color{brand}Why the other options are incorrect.}
\begin{itemize}[leftmargin=2.4em, itemsep=1pt, topsep=2pt]
\item[\textbf{(A)}] $\ket{111}$ is the state for CCZ alone; after $X_2$ it is mapped to $\ket{101}$, which receives no sign.
\item[\textbf{(B)}] $\ket{010}$ is the bit-complement of the right answer.
\item[\textbf{(D)}] $\ket{000}$ would need $X$ on all three qubits.
\end{itemize}
\textbf{\color{brand}Key takeaway.} Put an $X$ on every qubit where the marked string has a $0$.
\end{solbox}
\qsep

\needspace{13\baselineskip}
\begin{qitem}{27}
\qmeta{Single-correct MCQ}{CCZ from Toffoli}
Let CCX be the Toffoli gate with target qubit 3. The product $(I \otimes I \otimes H)\,\mathrm{CCX}\,(I \otimes I \otimes H)$ equals
\par\smallskip\noindent\begin{tabularx}{\linewidth}{@{}*{4}{L}@{}}
\optlabel{A}\,CCX & \optlabel{B}\,$I \otimes I \otimes I$ & \optlabel{C}\,$Z \otimes Z \otimes Z$ & \correctlabel{D}\,CCZ
\end{tabularx}
\end{qitem}
\begin{solbox}{D}
\textbf{\color{brand}Concepts required.} $HXH = Z$ and $HZH = X$. Conjugating the target of a controlled gate conjugates the gate itself.

\textbf{\color{brand}Relevant formulae.} $HXH = Z$, $\;HIH = I$; a controlled gate acts on its target only in the branch where all the controls are $1$.

\textbf{\color{brand}Step-by-step explanation.} Use $HXH = Z$ and $HIH = I$. In the branch where both controls are $1$, the target sees $HXH = Z$; in all other branches it sees $HIH = I$. So the product is
the doubly controlled $Z$, i.e.\ CCZ $= \operatorname{diag}(1, \dots, 1, -1)$.

\textbf{\color{brand}Why the other options are incorrect.}
\begin{itemize}[leftmargin=2.4em, itemsep=1pt, topsep=2pt]
\item[\textbf{(A)}] CCX is the gate before conjugation; $HXH \ne X$.
\item[\textbf{(B)}] The product cannot be $I$ because in the branch $\ket{11}$ the target sees $Z$.
\item[\textbf{(C)}] $Z \otimes Z \otimes Z$ flips the sign for every odd-weight string, not only for $\ket{111}$.
\end{itemize}
\textbf{\color{brand}Key takeaway.} A multi-controlled $Z$ is a multi-controlled $X$ with Hadamards around the target.
\end{solbox}
\qsep

\needspace{13\baselineskip}
\begin{qitem}{28}
\qmeta{Single-correct MCQ}{Global vs.\ relative phase}
Which pair of states are physically \emph{indistinguishable} (they give the same statistics for every possible measurement)?
\begin{itemize}[leftmargin=2.6em, labelsep=0.4em, itemsep=2pt, topsep=3pt]
\item[\correctlabel{A}] $\frac{1}{\sqrt2}(\ket0 + \ket1)$ and $-\frac{1}{\sqrt2}(\ket0 + \ket1)$
\item[\optlabel{B}] $\frac{1}{\sqrt2}(\ket{00} + \ket{11})$ and $\frac{1}{\sqrt2}(\ket{00} - \ket{11})$
\item[\optlabel{C}] $\frac{1}{\sqrt2}(\ket0 + \ket1)$ and $\frac{1}{\sqrt2}(\ket0 - \ket1)$
\item[\optlabel{D}] $\frac{1}{\sqrt2}(\ket0 + \ket1)$ and $\frac{1}{\sqrt2}(\ket0 + i\ket1)$
\end{itemize}
\end{qitem}
\begin{solbox}{A}
\textbf{\color{brand}Concepts required.} Global phase: $e^{i\gamma}\ket\psi$ is the same physical state. A relative phase is observable.

\textbf{\color{brand}Relevant formulae.} Same physical state: $\ket{\psi'} = e^{i\gamma}\ket\psi$ for a real $\gamma$. $\;|\braket{\phi|\psi'}|^2 = |\braket{\phi|\psi}|^2$ for all $\ket\phi$.

\textbf{\color{brand}Step-by-step explanation.} Two state vectors describe the same physical state iff they differ by one overall factor $e^{i\gamma}$. Only the first pair does: the factor is $-1$. In each
of the other pairs the \emph{relative} phase differs, and a suitable measurement separates them. For example $H$ maps $\ket{+}$ and $\ket{-}$ to $\ket0$ and $\ket1$.
The minus sign in the diffusion circuit (see Q.47) is of the harmless global kind.

\textbf{\color{brand}Why the other options are incorrect.}
\begin{itemize}[leftmargin=2.4em, itemsep=1pt, topsep=2pt]
\item[\textbf{(B)}] These two Bell states are orthogonal, so a suitable measurement separates them.
\item[\textbf{(C)}] $\ket{+}$ and $\ket{-}$ are orthogonal, and an $X$-basis measurement separates them perfectly.
\item[\textbf{(D)}] These differ by a relative phase $i$ and have different $\langle Y\rangle$.
\end{itemize}
\textbf{\color{brand}Key takeaway.} Test with one common factor across all components: if none exists, the states differ.
\end{solbox}
\qsep

\needspace{13\baselineskip}
\begin{qitem}{29}
\qmeta{Single-correct MCQ}{Subspace measurement}
Let $\ket\psi = \frac{1}{\sqrt{10}}\big(\ket{00} + 2\ket{01} - 2\ket{10} + i\ket{11}\big)$. A measurement asks only whether the two qubits have \emph{different} values
(outcome ``different'' on the span of $\ket{01}, \ket{10}$). The probability of ``different'' and the state afterwards are
\par\smallskip\noindent\begin{tabularx}{\linewidth}{@{}*{2}{L}@{}}
\optlabel{A}\,$\frac45$ and $\ket{01}$ & \correctlabel{B}\,$\frac45$ and $\frac{1}{\sqrt2}(\ket{01} - \ket{10})$ \\[3pt]
\optlabel{C}\,$\frac45$ and $\frac{1}{\sqrt{10}}(2\ket{01} - 2\ket{10})$ & \optlabel{D}\,$\frac25$ and $\frac{1}{\sqrt2}(\ket{01} - \ket{10})$
\end{tabularx}
\end{qitem}
\begin{solbox}{B}
\textbf{\color{brand}Concepts required.} Projective measurement: sum the probabilities in the subspace, project, and re-normalise.

\textbf{\color{brand}Relevant formulae.} $P = \sum_{x\in S}|c_x|^2$ for the subspace spanned by the set $S$; the new state is $\frac{1}{\sqrt P}\sum_{x\in S}c_x\ket{x}$.

\textbf{\color{brand}Step-by-step explanation.} Sum the probabilities in the subspace: $P = \frac{4 + 4}{10} = \frac45$. The post-measurement state is the projected vector, re-normalised:
$\frac{(2\ket{01} - 2\ket{10})/\sqrt{10}}{\sqrt{4/5}} = \frac{2\ket{01} - 2\ket{10}}{\sqrt8} = \frac{1}{\sqrt2}(\ket{01} - \ket{10})$.
Relative phases inside the subspace survive, and the measurement does not pick one of the basis states.

\textbf{\color{brand}Why the other options are incorrect.}
\begin{itemize}[leftmargin=2.4em, itemsep=1pt, topsep=2pt]
\item[\textbf{(A)}] The measurement does not distinguish $\ket{01}$ from $\ket{10}$, so it does not collapse to one of them.
\item[\textbf{(C)}] This vector was not re-normalised; its length is $\sqrt{4/5} \ne 1$.
\item[\textbf{(D)}] $\frac25$ uses the amplitudes instead of squares or counts only one of the two terms.
\end{itemize}
\textbf{\color{brand}Key takeaway.} A coarse measurement keeps superpositions inside the measured subspace.
\end{solbox}
\qsep

\needspace{13\baselineskip}
\begin{qitem}{30}
\qmeta{Single-correct MCQ}{Unitary vs.\ projector}
Let $\ket{u}$ be a unit vector in a space of dimension $N \ge 2$. Which of these operators is \emph{not} unitary?
\par\smallskip\noindent\begin{tabularx}{\linewidth}{@{}*{2}{L}@{}}
\optlabel{A}\,$H^{\otimes n}$ (for $N = 2^n$) & \correctlabel{B}\,$\ket{u}\!\bra{u}$ \\[3pt]
\optlabel{C}\,$I - 2\ket{u}\!\bra{u}$ & \optlabel{D}\,$2\ket{u}\!\bra{u} - I$
\end{tabularx}
\end{qitem}
\begin{solbox}{B}
\textbf{\color{brand}Concepts required.} Unitary means $U^\dagger U = I$, and hence invertible. Projectors have eigenvalue $0$.

\textbf{\color{brand}Relevant formulae.} Unitary: $U^\dagger U = I$. Projector: $\Pi = \ket{u}\!\bra{u}$, $\;\Pi^2 = \Pi$. Reflection: $R = R^\dagger$, $\;R^2 = I$.

\textbf{\color{brand}Step-by-step explanation.} $\ket{u}\!\bra{u}$ is a projector: it has eigenvalue $0$, so it is not invertible and cannot be unitary. A reflection $R = \pm(I - 2\ket{u}\!\bra{u})$ is
Hermitian with $R^2 = I$, so $R^\dagger R = I$. $H^{\otimes n}$ is a tensor product of unitaries.

\textbf{\color{brand}Why the other options are incorrect.}
\begin{itemize}[leftmargin=2.4em, itemsep=1pt, topsep=2pt]
\item[\textbf{(A)}] A tensor product of unitaries is unitary.
\item[\textbf{(C)}] This is a reflection: Hermitian and squares to $I$, so it is unitary.
\item[\textbf{(D)}] This is also a reflection, the negative of the previous one.
\end{itemize}
\textbf{\color{brand}Key takeaway.} Projectors are not gates; reflections are.
\end{solbox}
\qsep

\needspace{13\baselineskip}
\begin{qitem}{31}
\qmeta{Single-correct MCQ}{Overlap $\cos\theta$}
For $N = 16$ let $\ket{s}$ be the uniform superposition of all $16$ basis states and $\ket{r}$ the uniform superposition of $15$ of them (all except one). The value of $\braket{r|s}$ is
\par\smallskip\noindent\begin{tabularx}{\linewidth}{@{}*{4}{L}@{}}
\correctlabel{A}\,$\dfrac{\sqrt{15}}{4}$ & \optlabel{B}\,$\dfrac{15}{16}$ & \optlabel{C}\,$\dfrac14$ & \optlabel{D}\,$\dfrac{1}{\sqrt{15}}$
\end{tabularx}
\end{qitem}
\begin{solbox}{A}
\textbf{\color{brand}Concepts required.} Overlap is the sum of products of matching amplitudes. Here $\cos\theta = \braket{r|s}$ and $\sin\theta = \braket{m|s}$.

\textbf{\color{brand}Relevant formulae.} $\braket{r|s} = \sum_x\overline{r_x}s_x$; with $\ket{s}$ uniform over $N$ states and $\ket{r}$ uniform over $N - M$: $\;\braket{r|s} = \sqrt{(N-M)/N} = \cos\theta$.

\textbf{\color{brand}Step-by-step explanation.} $\braket{r|s} = \sum_{x}\frac{1}{\sqrt{15}}\cdot\frac{1}{\sqrt{16}}$ over the $15$ common basis states, which is $\frac{15}{\sqrt{15}\cdot 4} = \frac{\sqrt{15}}{4}$.
It equals $\cos\theta$ with $\sin\theta = \frac14$. $\frac{15}{16}$ is $\cos^2\theta$: the probability, not the overlap.

\textbf{\color{brand}Why the other options are incorrect.}
\begin{itemize}[leftmargin=2.4em, itemsep=1pt, topsep=2pt]
\item[\textbf{(B)}] $\frac{15}{16}$ is $|\braket{r|s}|^2$, the probability.
\item[\textbf{(C)}] $\frac14 = \sin\theta$ is the overlap of $\ket{s}$ with the \emph{marked} state.
\item[\textbf{(D)}] $\frac{1}{\sqrt{15}}$ is a single amplitude of $\ket{r}$.
\end{itemize}
\textbf{\color{brand}Key takeaway.} The overlaps of $\ket{s}$ with the two axes define $\theta$.
\end{solbox}
\qsep

\needspace{13\baselineskip}
\begin{qitem}{32}
\qmeta{Single-correct MCQ}{State-vector memory}
A state vector of $n = 33$ qubits is stored in double precision ($16$ bytes per complex amplitude). The memory needed is
\par\smallskip\noindent\begin{tabularx}{\linewidth}{@{}*{4}{L}@{}}
\optlabel{A}\,$256$ GiB & \optlabel{B}\,$16$ GiB & \optlabel{C}\,$64$ GiB & \correctlabel{D}\,$128$ GiB
\end{tabularx}
\end{qitem}
\begin{solbox}{D}
\textbf{\color{brand}Concepts required.} Memory $= 2^n \times$ bytes per amplitude, and $2^{30}$ bytes $= 1$ GiB.

\textbf{\color{brand}Relevant formulae.} Memory $= 2^n\times$ (bytes per amplitude); double precision uses $16$ bytes; $1$ GiB $= 2^{30}$ bytes.

\textbf{\color{brand}Step-by-step explanation.} $2^{33}\times 16 = 2^{33}\times 2^{4} = 2^{37}$ bytes $= 2^{7}\times 2^{30}$ bytes $= 128$ GiB. Every extra qubit doubles it. ($16$ GiB is the value for $n = 30$.)

\textbf{\color{brand}Why the other options are incorrect.}
\begin{itemize}[leftmargin=2.4em, itemsep=1pt, topsep=2pt]
\item[\textbf{(A)}] $256$ GiB is the value for $n = 34$.
\item[\textbf{(B)}] $16$ GiB is the value for $n = 30$.
\item[\textbf{(C)}] $64$ GiB is the value for $n = 32$, or for $n = 33$ in single precision.
\end{itemize}
\textbf{\color{brand}Key takeaway.} Memory doubles with every qubit: $n = 30$ needs $16$ GiB in double precision.
\end{solbox}
\qsep

\needspace{13\baselineskip}
\begin{qitem}{33}
\qmeta{Single-correct MCQ}{Diffusion for one qubit}
For a single qubit ($n = 1$), the operator $2\ket{+}\!\bra{+} - I$ equals
\par\smallskip\noindent\begin{tabularx}{\linewidth}{@{}*{4}{L}@{}}
\correctlabel{A}\,$X$ & \optlabel{B}\,$H$ & \optlabel{C}\,$Z$ & \optlabel{D}\,$-I$
\end{tabularx}
\end{qitem}
\begin{solbox}{A}
\textbf{\color{brand}Concepts required.} $\ket{+}\!\bra{+}$ has all entries $\frac12$; the reflection about $\ket{+}$ fixes $\ket{+}$ and negates $\ket{-}$.

\textbf{\color{brand}Relevant formulae.} $\ket{+}\!\bra{+} = \frac12\begin{psmallmatrix}1&1\\1&1\end{psmallmatrix}$; $\;X = \begin{psmallmatrix}0&1\\1&0\end{psmallmatrix}$.

\textbf{\color{brand}Step-by-step explanation.} $\ket{+}\!\bra{+} = \frac12\begin{psmallmatrix}1&1\\1&1\end{psmallmatrix}$, so $2\ket{+}\!\bra{+} - I = \begin{psmallmatrix}1&1\\1&1\end{psmallmatrix} - \begin{psmallmatrix}1&0\\0&1\end{psmallmatrix} = \begin{psmallmatrix}0&1\\1&0\end{psmallmatrix} = X$.
\emph{Check:} $X\ket{+} = \ket{+}$ and $X\ket{-} = -\ket{-}$, which are the eigenvalues required of this reflection.

\textbf{\color{brand}Why the other options are incorrect.}
\begin{itemize}[leftmargin=2.4em, itemsep=1pt, topsep=2pt]
\item[\textbf{(B)}] $H$ does not fix $\ket{+}$: it maps $\ket{+}$ to $\ket0$.
\item[\textbf{(C)}] $Z$ fixes $\ket0$ and negates $\ket1$: that is the reflection about $\ket0$.
\item[\textbf{(D)}] $-I$ would negate $\ket{+}$ too, not fix it.
\end{itemize}
\textbf{\color{brand}Key takeaway.} Diffusion on one qubit is simply $X$.
\end{solbox}
\qsep

\needspace{13\baselineskip}
\begin{qitem}{34}
\qmeta{Single-correct MCQ}{$\ket{s}$ is fixed by $D$}
For $n$ qubits let $\ket{s} = H^{\otimes n}\ket{0\dots0}$ and $D = 2\ket{s}\!\bra{s} - I$. If the oracle marks nothing (it acts as the identity), then after $k$ iterations the state is
\par\smallskip\noindent\begin{tabularx}{\linewidth}{@{}*{2}{L}@{}}
\optlabel{A}\,$(-1)^k\ket{s}$ & \optlabel{B}\,$\ket{0\dots0}$ for odd $k$ \\[3pt]
\optlabel{C}\,a state whose amplitudes depend on $k$ & \correctlabel{D}\,$\ket{s}$ for every $k$
\end{tabularx}
\end{qitem}
\begin{solbox}{D}
\textbf{\color{brand}Concepts required.} $\ket{s}$ is the $+1$ eigenvector of $D$.

\textbf{\color{brand}Relevant formulae.} $D\ket{s} = 2\ket{s}\braket{s|s} - \ket{s} = \ket{s}$, with $\braket{s|s} = 1$.

\textbf{\color{brand}Step-by-step explanation.} $D\ket{s} = 2\ket{s}\braket{s|s} - \ket{s} = \ket{s}$, so $\ket{s}$ is an eigenvector of $D$ with eigenvalue $+1$. With $U_\omega = I$, each iteration is just $D$ and leaves $\ket{s}$ unchanged.
Measuring gives a uniformly random string, which matches $\sin^2\theta = 0$: nothing can be amplified when no item is marked.

\textbf{\color{brand}Why the other options are incorrect.}
\begin{itemize}[leftmargin=2.4em, itemsep=1pt, topsep=2pt]
\item[\textbf{(A)}] The eigenvalue is $+1$, so no factor $(-1)^k$ arises. (The circuit form $-D$ would give such a factor, but it is a global phase.)
\item[\textbf{(B)}] $D$ does not map $\ket{s}$ to $\ket{0\dots0}$; that is what $H^{\otimes n}$ does.
\item[\textbf{(C)}] The amplitudes stay at $1/\sqrt N$ because $\ket{s}$ is an eigenvector.
\end{itemize}
\textbf{\color{brand}Key takeaway.} Without a marked item the algorithm does nothing, as $\theta = 0$ predicts.
\end{solbox}
\qsep

\needspace{13\baselineskip}
\begin{qitem}{35}
\qmeta{Matching lists}{Single-qubit gate actions}
Match each expression in List-I with its value in List-II.
\par\smallskip\noindent{\renewcommand{\arraystretch}{1.35}%
\begin{tabularx}{\linewidth}{@{}c>{\raggedright\arraybackslash}p{0.33\linewidth}@{\hspace{1.2em}}cL@{}}
\toprule
\multicolumn{2}{@{}l}{\textbf{List-I}} & \multicolumn{2}{l@{}}{\textbf{List-II}} \\ \midrule
\textbf{(P)} & $H\ket1$ & \textbf{(1)} & $-\ket{-}$ \\
\textbf{(Q)} & $Z\ket{-}$ & \textbf{(2)} & $\ket0$ \\
\textbf{(R)} & $X\ket{-}$ & \textbf{(3)} & $\ket{+}$ \\
\textbf{(S)} & $H\ket{+}$ & \textbf{(4)} & $\ket{-}$ \\
\bottomrule
\end{tabularx}}
\par\smallskip\noindent\begin{tabularx}{\linewidth}{@{}*{2}{L}@{}}
\optlabel{A}\,P$\to$4, Q$\to$1, R$\to$3, S$\to$2 & \optlabel{B}\,P$\to$4, Q$\to$3, R$\to$2, S$\to$1 \\[3pt]
\correctlabel{C}\,P$\to$4, Q$\to$3, R$\to$1, S$\to$2 & \optlabel{D}\,P$\to$3, Q$\to$4, R$\to$1, S$\to$2
\end{tabularx}
\end{qitem}
\begin{solbox}{C}
\textbf{\color{brand}Concepts required.} $H\ket0 = \ket{+}$, $H\ket1 = \ket{-}$, $Z$ exchanges $\ket{\pm}$, and $X\ket{\pm} = \pm\ket{\pm}$.

\textbf{\color{brand}Relevant formulae.} $H\ket0 = \ket{+}$, $H\ket1 = \ket{-}$, $H^2 = I$; $\;Z\ket\pm = \ket\mp$; $\;X\ket\pm = \pm\ket\pm$.

\textbf{\color{brand}Step-by-step explanation.} $H\ket1 = \ket{-}$ (4). $Z\ket{-} = \frac{1}{\sqrt2}(\ket0 + \ket1) = \ket{+}$ (3). $X\ket{-} = \frac{1}{\sqrt2}(\ket1 - \ket0) = -\ket{-}$ (1).
$H\ket{+} = \ket0$ because $H^2 = I$ (2). Hence P$\to$4, Q$\to$3, R$\to$1, S$\to$2.

\textbf{\color{brand}Why the other options are incorrect.}
\begin{itemize}[leftmargin=2.4em, itemsep=1pt, topsep=2pt]
\item[\textbf{(A)}] This gives $Z\ket{-} = -\ket{-}$, which is the action of $X$; $Z$ exchanges $\ket{\pm}$.
\item[\textbf{(B)}] This exchanges R and S: $X\ket{-}$ cannot be $\ket0$, and $H\ket{+}$ is $\ket0$.
\item[\textbf{(D)}] This exchanges P and Q: $H\ket1$ is $\ket{-}$, and $Z\ket{-}$ is $\ket{+}$.
\end{itemize}
\textbf{\color{brand}Key takeaway.} Keep the table of $H, X, Z$ on $\ket0, \ket1, \ket\pm$ at your fingertips.
\end{solbox}
\qsep

\needspace{13\baselineskip}
\begin{qitem}{36}
\qmeta{Matching lists}{Product vs.\ entangled states}
Match each two-qubit state in List-I with its description in List-II.
\par\smallskip\noindent{\renewcommand{\arraystretch}{1.35}%
\begin{tabularx}{\linewidth}{@{}c>{\raggedright\arraybackslash}p{0.33\linewidth}@{\hspace{1.2em}}cL@{}}
\toprule
\multicolumn{2}{@{}l}{\textbf{List-I}} & \multicolumn{2}{l@{}}{\textbf{List-II}} \\ \midrule
\textbf{(P)} & $\frac12(\ket{00} + \ket{01} + \ket{10} + \ket{11})$ & \textbf{(1)} & entangled, with four amplitudes of equal size \\
\textbf{(Q)} & $\frac12(\ket{00} - \ket{01} - \ket{10} + \ket{11})$ & \textbf{(2)} & the product $\ket{-}\ket{-}$ \\
\textbf{(R)} & $\frac12(\ket{00} + \ket{01} + \ket{10} - \ket{11})$ & \textbf{(3)} & entangled, with only two non-zero amplitudes \\
\textbf{(S)} & $\frac{1}{\sqrt2}(\ket{00} + \ket{11})$ & \textbf{(4)} & the product $\ket{+}\ket{+}$ \\
\bottomrule
\end{tabularx}}
\par\smallskip\noindent\begin{tabularx}{\linewidth}{@{}*{2}{L}@{}}
\optlabel{A}\,P$\to$4, Q$\to$2, R$\to$3, S$\to$1 & \correctlabel{B}\,P$\to$4, Q$\to$2, R$\to$1, S$\to$3 \\[3pt]
\optlabel{C}\,P$\to$4, Q$\to$1, R$\to$2, S$\to$3 & \optlabel{D}\,P$\to$2, Q$\to$4, R$\to$1, S$\to$3
\end{tabularx}
\end{qitem}
\begin{solbox}{B}
\textbf{\color{brand}Concepts required.} A two-qubit state is a product iff $c_{00}c_{11} = c_{01}c_{10}$.

\textbf{\color{brand}Relevant formulae.} $\sum_{j,k}c_{jk}\ket{jk}$ is a product state $\iff c_{00}c_{11} = c_{01}c_{10}$.

\textbf{\color{brand}Step-by-step explanation.} Test $c_{00}c_{11}$ against $c_{01}c_{10}$; they are equal exactly for a product state.
P: $\frac14 = \frac14$, product, equal to $\ket{+}\ket{+}$ (4). Q: $\frac14 = \frac14$, product $\ket{-}\ket{-}$ (2).
R: $-\frac14 \ne \frac14$, entangled with four equal magnitudes (1); this is the state after the oracle for $n = 2$. S: $\frac12 \ne 0$, entangled with two non-zero amplitudes (3).

\textbf{\color{brand}Why the other options are incorrect.}
\begin{itemize}[leftmargin=2.4em, itemsep=1pt, topsep=2pt]
\item[\textbf{(A)}] This exchanges R and S: state R has four non-zero amplitudes, and state S has two.
\item[\textbf{(C)}] This calls Q entangled, but it passes the product test and equals $\ket{-}\ket{-}$.
\item[\textbf{(D)}] This exchanges P and Q: all plus signs give $\ket{+}\ket{+}$.
\end{itemize}
\textbf{\color{brand}Key takeaway.} One sign flip on $\ket{s}$ with $n = 2$ gives an entangled state, while Hadamards alone never do.
\end{solbox}
\qsep

\needspace{13\baselineskip}
\begin{qitem}{37}
\qmeta{Assertion--Reason}{Why the oracle alone reveals nothing}
\textbf{Assertion (A):} Measuring in the computational basis immediately after a phase oracle acts on $\ket{s}$ gives a uniformly random string; the oracle step alone reveals nothing.\\[1mm]
\textbf{Reason (R):} The oracle multiplies each amplitude by $\pm1$, and the Born probabilities depend only on the magnitudes of the amplitudes.
\begin{itemize}[leftmargin=2.6em, labelsep=0.4em, itemsep=2pt, topsep=3pt]
\item[\correctlabel{A}] Both Assertion (A) and Reason (R) are true, and R is the correct explanation of A
\item[\optlabel{B}] Both A and R are true, but R is \emph{not} the correct explanation of A
\item[\optlabel{C}] A is true, but R is false
\item[\optlabel{D}] A is false, but R is true
\end{itemize}
\end{qitem}
\begin{solbox}{A}
\textbf{\color{brand}Concepts required.} Probabilities depend only on $|c_x|^2$, while the oracle changes only signs.

\textbf{\color{brand}Relevant formulae.} $U_\omega\ket{x} = (-1)^{f(x)}\ket{x}$; Born rule $P(x) = |c_x|^2$; $\;|\pm a|^2 = |a|^2$.

\textbf{\color{brand}Step-by-step explanation.} \textbf{R is true:} $U_\omega\ket{x} = (-1)^{f(x)}\ket{x}$ and $|\pm a|^2 = |a|^2$. \textbf{A is true} because of R: the magnitudes stay $1/\sqrt N$, so the distribution stays uniform.
R is the correct explanation. Answer: (A).

\textbf{\color{brand}Why the other options are incorrect.}
\begin{itemize}[leftmargin=2.4em, itemsep=1pt, topsep=2pt]
\item[\textbf{(B)}] R does explain A: it gives the exact reason that the distribution stays uniform.
\item[\textbf{(C)}] R is true: $|\pm a|^2 = |a|^2$.
\item[\textbf{(D)}] A is true: the output is uniformly random whichever item is marked.
\end{itemize}
\textbf{\color{brand}Key takeaway.} The oracle's effect is invisible until an interfering step such as diffusion follows.
\end{solbox}
\qsep

\needspace{13\baselineskip}
\begin{qitem}{38}
\qmeta{Statement I / II}{Relative phases are observable}
\textbf{Statement I:} The states $\ket\psi$ and $-\ket\psi$ give identical statistics for every measurement.\\[1mm]
\textbf{Statement II:} Multiplying only \emph{one} basis-state amplitude of a superposition by $-1$ can never change the statistics of any measurement.
\begin{itemize}[leftmargin=2.6em, labelsep=0.4em, itemsep=2pt, topsep=3pt]
\item[\optlabel{A}] Both Statement I and Statement II are true
\item[\optlabel{B}] Both Statement I and Statement II are false
\item[\correctlabel{C}] Statement I is true, but Statement II is false
\item[\optlabel{D}] Statement I is false, but Statement II is true
\end{itemize}
\end{qitem}
\begin{solbox}{C}
\textbf{\color{brand}Concepts required.} Global phases are unobservable; relative phases can be observed after interference.

\textbf{\color{brand}Relevant formulae.} Global phase: $e^{i\gamma}\ket\psi$. Relative phase: a phase on only some terms, as in $\ket0 - \ket1$ against $\ket0 + \ket1$.

\textbf{\color{brand}Step-by-step explanation.} \textbf{I is true:} a global phase cancels in $|\braket{\phi|\psi}|^2$ for every $\ket\phi$.
\textbf{II is false:} that is a \emph{relative} phase. Counterexample: $\ket{+}$ and $\ket{-}$ differ by the sign of $\ket1$ only, yet a measurement in the $\{\ket{+}, \ket{-}\}$ basis separates them perfectly.
Grover's algorithm exists because the sign is detectable after interference. Answer: (C).

\textbf{\color{brand}Why the other options are incorrect.}
\begin{itemize}[leftmargin=2.4em, itemsep=1pt, topsep=2pt]
\item[\textbf{(A)}] II is false: $\ket{+}$ and $\ket{-}$ differ only by the sign of one amplitude and yet are distinguishable.
\item[\textbf{(B)}] I is true: a global phase cancels in every probability.
\item[\textbf{(D)}] I is true, so this option fails.
\end{itemize}
\textbf{\color{brand}Key takeaway.} A sign on one amplitude is a relative phase: it is exactly the signal that Grover amplifies.
\end{solbox}
\qsep

\needspace{20\baselineskip}
\begin{passagebox}{Paragraph for Questions 39 and 40}
Three qubits start in $\ket{0}\ket{0}\ket{1}$ (input qubits 1 and 2, ancilla qubit 3). A Hadamard gate is applied to all three. Then the
Toffoli gate is applied with qubits 1 and 2 as controls and qubit 3 as target: it implements
$U_f\ket{x}\ket{y} = \ket{x}\ket{y \oplus f(x)}$ with $f(x_1x_2) = x_1 \land x_2$.
\end{passagebox}
\needspace{13\baselineskip}
\begin{qitem}{39}
\qmeta{Paragraph based}{Toffoli as a phase oracle}
The state of the three qubits after the Toffoli gate is
\par\smallskip\noindent\begin{tabularx}{\linewidth}{@{}*{2}{L}@{}}
\optlabel{A}\,$\frac12\big(\ket{00} - \ket{01} - \ket{10} - \ket{11}\big)\ket{-}$ & \optlabel{B}\,$\frac12\big(\ket{00} + \ket{01} + \ket{10} - \ket{11}\big)\ket{+}$ \\[3pt]
\optlabel{C}\,$\frac12\big(\ket{00} + \ket{01} + \ket{10} + \ket{11}\big)\ket{-}$ & \correctlabel{D}\,$\frac12\big(\ket{00} + \ket{01} + \ket{10} - \ket{11}\big)\ket{-}$
\end{tabularx}
\end{qitem}
\begin{solbox}{D}
\textbf{\color{brand}Concepts required.} $H\ket1 = \ket{-}$; kickback puts the sign $(-1)^{f(x)}$ on $\ket{x}$.

\textbf{\color{brand}Relevant formulae.} $H\ket1 = \ket{-}$; $\;U_f\ket{x}\ket{-} = (-1)^{f(x)}\ket{x}\ket{-}$; here $f(x_1x_2) = x_1\land x_2$.

\textbf{\color{brand}Step-by-step explanation.} After the Hadamards: $\frac12\sum_x\ket{x}\ket{-}$, since $H\ket1 = \ket{-}$. By kickback each term gets $(-1)^{f(x)}$, and $f = 1$ only for $x = 11$.
So the state is $\frac12(\ket{00} + \ket{01} + \ket{10} - \ket{11})\ket{-}$. The ancilla stays in $\ket{-}$, unentangled with the inputs.

\textbf{\color{brand}Why the other options are incorrect.}
\begin{itemize}[leftmargin=2.4em, itemsep=1pt, topsep=2pt]
\item[\textbf{(A)}] Only $x = 11$ has $f = 1$, so only one sign flips.
\item[\textbf{(B)}] The ancilla remains in $\ket{-}$; it does not change to $\ket{+}$.
\item[\textbf{(C)}] This forgets the kicked-back sign on $\ket{11}$.
\end{itemize}
\textbf{\color{brand}Key takeaway.} Toffoli plus ancilla $\ket{-}$ is the phase oracle for $x_1 \land x_2$.
\end{solbox}
\qsep

\needspace{13\baselineskip}
\begin{qitem}{40}
\qmeta{Paragraph based}{Ancilla stays unentangled}
Now the ancilla alone is measured in the computational basis. The probability of the outcome $1$, and the state of the input qubits afterwards, are
\begin{itemize}[leftmargin=2.6em, labelsep=0.4em, itemsep=2pt, topsep=3pt]
\item[\optlabel{A}] $\frac14$; the input qubits collapse to $\ket{11}$
\item[\correctlabel{B}] $\frac12$; the input qubits remain in $\frac12\big(\ket{00} + \ket{01} + \ket{10} - \ket{11}\big)$
\item[\optlabel{C}] $\frac12$; the input qubits collapse to $\ket{11}$
\item[\optlabel{D}] $1$; the input qubits remain in $\frac12\big(\ket{00} + \ket{01} + \ket{10} - \ket{11}\big)$
\end{itemize}
\end{qitem}
\begin{solbox}{B}
\textbf{\color{brand}Concepts required.} A measurement on one factor of a product state does not disturb the other factor.

\textbf{\color{brand}Relevant formulae.} For a product state $\ket\phi\ket{-}$, measuring the second factor leaves $\ket\phi$ unchanged; $\;P(1) = |\braket{1|-}|^2 = \frac12$.

\textbf{\color{brand}Step-by-step explanation.} The state after the gate is a product of the input state and $\ket{-}$ (see Q.39). Measuring $\ket{-}$ gives $1$ with probability $\frac12$, and because the ancilla is not entangled
with the inputs, the input qubits are not disturbed.

\textbf{\color{brand}Why the other options are incorrect.}
\begin{itemize}[leftmargin=2.4em, itemsep=1pt, topsep=2pt]
\item[\textbf{(A)}] This treats the ancilla as if it were in $\ket{f(x)}$, i.e.\ started in $\ket0$.
\item[\textbf{(C)}] The probability is right, but the inputs do not collapse because there is no entanglement.
\item[\textbf{(D)}] $\ket{-}$ gives outcome $1$ with probability $\frac12$, not $1$.
\end{itemize}
\textbf{\color{brand}Key takeaway.} With an ancilla in $\ket{-}$ the ancilla stays unentangled and can be re-used.
\end{solbox}
\qsep

\newpage

\begin{sectionbanner}SECTION C \enspace$\cdot$\enspace Grover's Algorithm Fundamentals \hfill Q.41--60\end{sectionbanner}

\needspace{13\baselineskip}
\begin{qitem}{41}
\qmeta{Single-correct MCQ}{$N = 4$ is exact in one step}
Grover's algorithm is run on $n = 2$ qubits ($N = 4$) with one marked item. After exactly one iteration, the probability of measuring the marked item is
\par\smallskip\noindent\begin{tabularx}{\linewidth}{@{}*{4}{L}@{}}
\optlabel{A}\,$\dfrac12$ & \optlabel{B}\,$\dfrac14$ & \optlabel{C}\,$\dfrac34$ & \correctlabel{D}\,$1$
\end{tabularx}
\end{qitem}
\begin{solbox}{D}
\textbf{\color{brand}Concepts required.} $\sin\theta = 1/\sqrt N$ and $P_k = \sin^2((2k+1)\theta)$. For $N = 4$, $\theta = 30^\circ$.

\textbf{\color{brand}Relevant formulae.} $\sin\theta = \sqrt{M/N}$, with $N$ items and $M$ marked; $\;P_k = \sin^2((2k+1)\theta)$ after $k$ iterations.

\textbf{\color{brand}Step-by-step explanation.} $\sin\theta = \frac{1}{\sqrt4} = \frac12$, so $\theta = 30^\circ$. After $k = 1$ iteration $P = \sin^2(3\theta) = \sin^2 90^\circ = 1$.
In amplitudes: after the oracle the entries are $(\frac12, \frac12, \frac12, -\frac12)$ with mean $\frac14$; inversion about the mean gives $(0, 0, 0, 1)$.

\textbf{\color{brand}Why the other options are incorrect.}
\begin{itemize}[leftmargin=2.4em, itemsep=1pt, topsep=2pt]
\item[\textbf{(A)}] $\frac12$ is the result for $M = N/2$, not for this case.
\item[\textbf{(B)}] $\frac14$ is the probability before any iteration.
\item[\textbf{(C)}] $\frac34$ is the total probability of the unmarked items at the start.
\end{itemize}
\textbf{\color{brand}Key takeaway.} $N = 4$ is the unique single-item case where one iteration is exact.
\end{solbox}
\qsep

\needspace{13\baselineskip}
\begin{qitem}{42}
\qmeta{Single-correct MCQ}{One iteration, $N = 8$}
For $N = 8$ with one marked item, the probability of measuring the marked item after one iteration is
\par\smallskip\noindent\begin{tabularx}{\linewidth}{@{}*{4}{L}@{}}
\optlabel{A}\,$1$ & \correctlabel{B}\,$\dfrac{25}{32}$ & \optlabel{C}\,$\dfrac58$ & \optlabel{D}\,$\dfrac{9}{16}$
\end{tabularx}
\end{qitem}
\begin{solbox}{B}
\textbf{\color{brand}Concepts required.} Trace the amplitudes with inversion about the mean, or use $P_1 = s(3 - 4s)^2$ with $s = 1/N$.

\textbf{\color{brand}Relevant formulae.} Oracle: marked $c \mapsto -c$. Diffusion: $c_x \mapsto 2\bar c - c_x$ with $\bar c = \frac1N\sum_xc_x$. Also $P_1 = s(3 - 4s)^2$, with $s = M/N$.

\textbf{\color{brand}Step-by-step explanation.} Write $a = 1/\sqrt8$. After the oracle: seven entries $a$, one entry $-a$. Mean $= \frac{6a}{8} = \frac34a$. The marked entry becomes $2\cdot\frac34a + a = \frac52a$, so
$P = \frac{25}{4}a^2 = \frac{25}{4}\cdot\frac18 = \frac{25}{32}$. \emph{Check:} $\sin^2 3\theta = s(3 - 4s)^2$ with $s = \frac18$ gives $\frac18\cdot\frac{25}{4} = \frac{25}{32}$.

\textbf{\color{brand}Why the other options are incorrect.}
\begin{itemize}[leftmargin=2.4em, itemsep=1pt, topsep=2pt]
\item[\textbf{(A)}] $P = 1$ happens for $N = 4$ only; here $3\theta \ne 90^\circ$.
\item[\textbf{(C)}] $\frac58$ uses the amplitude factor $\frac52$ divided by $4$ and forgets to square.
\item[\textbf{(D)}] $\frac{9}{16}$ uses the mean $\frac34$ squared instead of the new amplitude.
\end{itemize}
\textbf{\color{brand}Key takeaway.} For small $N$, inversion about the mean is the fastest method.
\end{solbox}
\qsep

\needspace{13\baselineskip}
\begin{qitem}{43}
\qmeta{Single-correct MCQ}{$k_{\rm opt}$ and $P$ for $n = 4$}
For $n = 4$ qubits with one marked item ($\theta \approx 0.2527$ rad), the optimal number of iterations $k_{\rm opt}$ and the resulting success probability are
\par\smallskip\noindent\begin{tabularx}{\linewidth}{@{}*{2}{L}@{}}
\optlabel{A}\,$k_{\rm opt} = 2$ and $P \approx 0.908$ & \optlabel{B}\,$k_{\rm opt} = 4$ and $P \approx 0.582$ \\[3pt]
\optlabel{C}\,$k_{\rm opt} = 3$ and $P = 1$ & \correctlabel{D}\,$k_{\rm opt} = 3$ and $P \approx 0.961$
\end{tabularx}
\end{qitem}
\begin{solbox}{D}
\textbf{\color{brand}Concepts required.} $k_{\rm opt} = \operatorname{round}\big(\frac{\pi}{4\theta} - \frac12\big)$ and $P = \sin^2((2k+1)\theta)$.

\textbf{\color{brand}Relevant formulae.} $k_{\rm opt} = \operatorname{round}\big(\frac{\pi}{4\theta} - \frac12\big)$; $\;P_k = \sin^2((2k+1)\theta)$; $\;\sin\theta = 1/\sqrt N$ for one marked item.

\textbf{\color{brand}Step-by-step explanation.} $\frac{\pi}{4\theta} - \frac12 = \frac{0.7854}{0.2527} - 0.5 = 3.108 - 0.5 = 2.608$, which rounds to $k_{\rm opt} = 3$. Then
$P = \sin^2(7\theta) = \sin^2(1.769) \approx 0.961$. The angle $7\theta$ does not hit $\frac\pi2$ exactly, so $P < 1$.
For comparison, $k = 2$ gives $\sin^2 5\theta \approx 0.908$ and $k = 4$ gives $\sin^2 9\theta \approx 0.582$.

\textbf{\color{brand}Why the other options are incorrect.}
\begin{itemize}[leftmargin=2.4em, itemsep=1pt, topsep=2pt]
\item[\textbf{(A)}] $k = 2$ is below the optimum: $5\theta$ is further from $\frac\pi2$ than $7\theta$.
\item[\textbf{(B)}] $k = 4$ over-rotates, and $P$ falls to $0.58$.
\item[\textbf{(C)}] $k = 3$ is right but $P < 1$: $7\theta = 1.769$ rad, not $\frac\pi2$.
\end{itemize}
\textbf{\color{brand}Key takeaway.} Optimal does not mean perfect: $P$ is close to $1$ but usually not equal to $1$.
\end{solbox}
\qsep

\needspace{13\baselineskip}
\begin{qitem}{44}
\qmeta{Single-correct MCQ}{Matrix of $D$ for $N = 4$}
For $n = 2$ ($N = 4$), the matrix of the diffusion operator $D = 2\ket{s}\!\bra{s} - I$ has
\begin{itemize}[leftmargin=2.6em, labelsep=0.4em, itemsep=2pt, topsep=3pt]
\item[\optlabel{A}] diagonal entries $+\frac12$ and off-diagonal entries $-\frac12$
\item[\correctlabel{B}] diagonal entries $-\frac12$ and off-diagonal entries $+\frac12$
\item[\optlabel{C}] diagonal entries $-1$ and off-diagonal entries $+1$
\item[\optlabel{D}] all entries equal to $\frac12$
\end{itemize}
\end{qitem}
\begin{solbox}{B}
\textbf{\color{brand}Concepts required.} $\ket{s}\!\bra{s}$ has every entry $\frac1N$, so $D$ has diagonal $\frac2N - 1$ and off-diagonal $\frac2N$.

\textbf{\color{brand}Relevant formulae.} $D = 2\ket{s}\!\bra{s} - I$; $\;(\ket{s}\!\bra{s})_{xy} = \frac1N$; so $D_{xx} = \frac2N - 1$ and $D_{xy} = \frac2N$ for $x \ne y$.

\textbf{\color{brand}Step-by-step explanation.} Every entry of $\ket{s}\!\bra{s}$ is $\frac{1}{\sqrt4}\cdot\frac{1}{\sqrt4} = \frac14$, so every entry of $2\ket{s}\!\bra{s}$ is $\frac12$. Subtracting $I$ changes only the diagonal: $\frac12 - 1 = -\frac12$. In general the diagonal is $\frac2N - 1$ and the off-diagonal is $\frac2N$.
\emph{Check:} acting on the amplitudes $(\frac12, \frac12, \frac12, -\frac12)$, the last row gives $\frac12(\frac12 + \frac12 + \frac12) - \frac12(-\frac12) = \frac34 + \frac14 = 1$, the same as inversion about the mean.

\textbf{\color{brand}Why the other options are incorrect.}
\begin{itemize}[leftmargin=2.4em, itemsep=1pt, topsep=2pt]
\item[\textbf{(A)}] This is $-D = I - 2\ket{s}\!\bra{s}$, which is what the circuit implements; it differs from $D$ by a global phase.
\item[\textbf{(C)}] Forgets the factor $\frac14$ from the normalisation of $\ket{s}$; this matrix is not unitary.
\item[\textbf{(D)}] This is $2\ket{s}\!\bra{s}$, before $I$ has been subtracted.
\end{itemize}
\textbf{\color{brand}Key takeaway.} Row by row, $D$ computes $2\bar c - c_x$, which is inversion about the mean.
\end{solbox}
\qsep

\needspace{13\baselineskip}
\begin{qitem}{45}
\qmeta{Single-correct MCQ}{Building the oracle}
An oracle for the marked string $\ket{0101}$ is built as: $X$ gates, then the multi-controlled $Z$ on all four qubits, then the same $X$ gates. The $X$ gates act on
\par\smallskip\noindent\begin{tabularx}{\linewidth}{@{}*{4}{L}@{}}
\optlabel{A}\,qubits 2 and 4 & \optlabel{B}\,no qubit & \correctlabel{C}\,qubits 1 and 3 & \optlabel{D}\,all four qubits
\end{tabularx}
\end{qitem}
\begin{solbox}{C}
\textbf{\color{brand}Concepts required.} The multi-controlled $Z$ flips $\ket{1\dots1}$ only, so the $X$ gates must map the marked string to all ones.

\textbf{\color{brand}Relevant formulae.} $\mathrm{C}^{n-1}\mathrm{Z} = I - 2\ket{1\dots1}\!\bra{1\dots1}$; $\;X\ket0 = \ket1$; so $X$ on the $0$ bits of $x^\ast$ maps $\ket{x^\ast}$ to $\ket{1\dots1}$.

\textbf{\color{brand}Step-by-step explanation.} The multi-controlled $Z$ flips the sign of $\ket{1111}$ only, so the $X$ layer must map the marked string to $\ket{1111}$. That means flipping exactly the bits where
$x^\ast$ has a $0$. In $0101$ these are qubits 1 and 3. The second $X$ layer undoes the first, so the sign flip ends up on $\ket{0101}$.

\textbf{\color{brand}Why the other options are incorrect.}
\begin{itemize}[leftmargin=2.4em, itemsep=1pt, topsep=2pt]
\item[\textbf{(A)}] Qubits 2 and 4 hold $1$ already: flipping them maps $\ket{1010}$ to all ones, the wrong item.
\item[\textbf{(B)}] With no $X$ gates the oracle marks $\ket{1111}$.
\item[\textbf{(D)}] Flipping all qubits marks the complement $\ket{1010}$.
\end{itemize}
\textbf{\color{brand}Key takeaway.} Rule: $X$ on the zeros of $x^\ast$.
\end{solbox}
\qsep

\needspace{13\baselineskip}
\begin{qitem}{46}
\qmeta{Single-correct MCQ}{Gate count per iteration}
Use the circuit in the conventions box and count each multi-controlled $Z$ as one gate. The number of gates in \emph{one} Grover iteration for $n = 6$, $x^\ast = 010011$ is
\par\smallskip\noindent\begin{tabularx}{\linewidth}{@{}*{4}{L}@{}}
\correctlabel{A}\,$32$ & \optlabel{B}\,$30$ & \optlabel{C}\,$26$ & \optlabel{D}\,$38$
\end{tabularx}
\end{qitem}
\begin{solbox}{A}
\textbf{\color{brand}Concepts required.} Per iteration: oracle $2z + 1$ gates and diffusion $4n + 1$ gates, so $4n + 2z + 2$ in total.

\textbf{\color{brand}Relevant formulae.} Gates per iteration $= (2z + 1) + (4n + 1) = 4n + 2z + 2$, with $n$ qubits and $z$ zeros in $x^\ast$.

\textbf{\color{brand}Step-by-step explanation.} The oracle has $2z$ $X$ gates, where $z$ is the number of $0$s in $x^\ast$ (here $z = 3$), and one multi-controlled $Z$: $6 + 1 = 7$ gates.
The diffusion has $2n$ $H$, $2n$ $X$ and one multi-controlled $Z$: $12 + 12 + 1 = 25$ gates. Total $= 4n + 2z + 2 = 24 + 6 + 2 = 32$.

\textbf{\color{brand}Why the other options are incorrect.}
\begin{itemize}[leftmargin=2.4em, itemsep=1pt, topsep=2pt]
\item[\textbf{(B)}] $30$ omits the two multi-controlled $Z$ gates.
\item[\textbf{(C)}] $26$ omits the $2z$ oracle $X$ gates.
\item[\textbf{(D)}] $38$ counts $2n$ oracle $X$ gates, as if every bit were $0$.
\end{itemize}
\textbf{\color{brand}Key takeaway.} The formula $4n + 2z + 2$ is the one the notebook uses in its time estimate.
\end{solbox}
\qsep

\needspace{13\baselineskip}
\begin{qitem}{47}
\qmeta{Single-correct MCQ}{Diffusion circuit $= -D$}
Let $C$ be the multi-controlled $Z$ on $n$ qubits. The circuit $H^{\otimes n}X^{\otimes n}\,C\,X^{\otimes n}H^{\otimes n}$ implements
\par\smallskip\noindent\begin{tabularx}{\linewidth}{@{}*{2}{L}@{}}
\optlabel{A}\,$2\ket{s}\!\bra{s} - I = D$ & \correctlabel{B}\,$I - 2\ket{s}\!\bra{s} = -D$ \\[3pt]
\optlabel{C}\,$I - 2\ket{0\dots0}\!\bra{0\dots0}$ & \optlabel{D}\,$I - 2\ket{1\dots1}\!\bra{1\dots1}$
\end{tabularx}
\end{qitem}
\begin{solbox}{B}
\textbf{\color{brand}Concepts required.} $U(I - 2\ket{a}\!\bra{a})U^\dagger = I - 2U\ket{a}\!\bra{a}U^\dagger$; a reflection conjugated by a unitary is a reflection about the image.

\textbf{\color{brand}Relevant formulae.} $\mathrm{C}^{n-1}\mathrm{Z} = I - 2\ket{1\dots1}\!\bra{1\dots1}$; $\;U(I - 2\ket{a}\!\bra{a})U^\dagger = I - 2U\ket{a}\!\bra{a}U^\dagger$; $\;D = 2\ket{s}\!\bra{s} - I$.

\textbf{\color{brand}Step-by-step explanation.} $C = I - 2\ket{1\dots1}\!\bra{1\dots1}$. Conjugating by $X^{\otimes n}$ gives $I - 2\ket{0\dots0}\!\bra{0\dots0}$, and conjugating by $H^{\otimes n}$ gives $I - 2\ket{s}\!\bra{s}$, since $H^{\otimes n}\ket{0\dots0} = \ket{s}$.
That is $-D$: the circuit equals the diffusion up to the global phase $-1$, which has no observable effect.

\textbf{\color{brand}Why the other options are incorrect.}
\begin{itemize}[leftmargin=2.4em, itemsep=1pt, topsep=2pt]
\item[\textbf{(A)}] This is the intended operator, but the circuit gives its negative. Only a global phase separates them, but the exact identity is $-D$.
\item[\textbf{(C)}] This is the operator at an intermediate step, before the $H$ layers are included.
\item[\textbf{(D)}] This is the multi-controlled $Z$ itself.
\end{itemize}
\textbf{\color{brand}Key takeaway.} Diffusion circuit $= -D$, and the minus sign is a harmless global phase.
\end{solbox}
\qsep

\needspace{13\baselineskip}
\begin{qitem}{48}
\qmeta{Single-correct MCQ}{Unmarked amplitude}
For $N = 16$ with one marked item, the amplitude of each \emph{unmarked} item after one iteration is
\par\smallskip\noindent\begin{tabularx}{\linewidth}{@{}*{4}{L}@{}}
\optlabel{A}\,$\dfrac14$ & \optlabel{B}\,$\dfrac{11}{16}$ & \correctlabel{C}\,$\dfrac{3}{16}$ & \optlabel{D}\,$\dfrac{7}{32}$
\end{tabularx}
\end{qitem}
\begin{solbox}{C}
\textbf{\color{brand}Concepts required.} Inversion about the mean: new amplitude $= 2\times\text{mean} - \text{old}$, with the mean taken after the oracle.

\textbf{\color{brand}Relevant formulae.} $\bar c = \frac1N\sum_xc_x$ after the oracle; new amplitude $= 2\bar c - c_x$.

\textbf{\color{brand}Step-by-step explanation.} Start with amplitudes $\frac14$. After the oracle the mean is $\frac{15 - 1}{16}\cdot\frac14 = \frac{7}{32}$. An unmarked entry becomes $2\cdot\frac{7}{32} - \frac14 = \frac{7}{16} - \frac{4}{16} = \frac{3}{16}$,
and the marked entry becomes $\frac{7}{16} + \frac{4}{16} = \frac{11}{16}$. \emph{Check:} $15\cdot\frac{9}{256} + \frac{121}{256} = 1$.

\textbf{\color{brand}Why the other options are incorrect.}
\begin{itemize}[leftmargin=2.4em, itemsep=1pt, topsep=2pt]
\item[\textbf{(A)}] $\frac14$ is the initial amplitude.
\item[\textbf{(B)}] $\frac{11}{16}$ is the amplitude of the marked item.
\item[\textbf{(D)}] $\frac{7}{32}$ is the mean, not the new amplitude.
\end{itemize}
\textbf{\color{brand}Key takeaway.} Unmarked amplitudes shrink a little and the marked one grows a lot.
\end{solbox}
\qsep

\needspace{13\baselineskip}
\begin{qitem}{49}
\qmeta{Single-correct MCQ}{$M = N/4$}
Let $N = 16$ with $M = 4$ marked items. After one iteration the probability of measuring a marked item is
\par\smallskip\noindent\begin{tabularx}{\linewidth}{@{}*{4}{L}@{}}
\optlabel{A}\,$\dfrac34$ & \optlabel{B}\,$\dfrac{25}{32}$ & \correctlabel{C}\,$1$ & \optlabel{D}\,$\dfrac14$
\end{tabularx}
\end{qitem}
\begin{solbox}{C}
\textbf{\color{brand}Concepts required.} Several solutions: $\sin^2\theta = M/N$. $M = N/4$ gives $\theta = 30^\circ$.

\textbf{\color{brand}Relevant formulae.} $\sin^2\theta = M/N$; $\;P_1 = \sin^2 3\theta$.

\textbf{\color{brand}Step-by-step explanation.} $\sin^2\theta = \frac{M}{N} = \frac14$, so $\theta = 30^\circ$ and $P_1 = \sin^2 3\theta = \sin^2 90^\circ = 1$. With $M = N/4$ a single iteration succeeds with certainty, for every $N$ for which $M = N/4$ is an integer.
The final state is the uniform superposition of the four marked items.

\textbf{\color{brand}Why the other options are incorrect.}
\begin{itemize}[leftmargin=2.4em, itemsep=1pt, topsep=2pt]
\item[\textbf{(A)}] $\frac34$ is the probability of the unmarked set before any iteration.
\item[\textbf{(B)}] $\frac{25}{32}$ is the value for $s = \frac18$ and not for $s = \frac14$.
\item[\textbf{(D)}] $\frac14$ is the probability before any iteration.
\end{itemize}
\textbf{\color{brand}Key takeaway.} $M = N/4$ is the special ratio for a one-step exact search.
\end{solbox}
\qsep

\needspace{13\baselineskip}
\begin{qitem}{50}
\qmeta{Single-correct MCQ}{$M = N/2$}
Let $N = 8$ with $M = 4$ marked items. The probability of a marked outcome after $k$ iterations is
\par\smallskip\noindent\begin{tabularx}{\linewidth}{@{}*{2}{L}@{}}
\optlabel{A}\,$1$ for even $k$ and $0$ for odd $k$ & \optlabel{B}\,$0$ for $k = 1$ \\[3pt]
\optlabel{C}\,$1$ for $k = 1$ & \correctlabel{D}\,$\frac12$ for every $k$
\end{tabularx}
\end{qitem}
\begin{solbox}{D}
\textbf{\color{brand}Concepts required.} $M = N/2$ gives $\theta = 45^\circ$, and $2\theta = 90^\circ$ steps skip between the diagonals.

\textbf{\color{brand}Relevant formulae.} $\sin^2\theta = M/N$; $\;P_k = \sin^2((2k+1)\theta)$; $\;\sin^2 45^\circ = \sin^2 135^\circ = \frac12$.

\textbf{\color{brand}Step-by-step explanation.} $\sin^2\theta = \frac12$ gives $\theta = 45^\circ$, and $P_k = \sin^2\big((2k+1)45^\circ\big)$. The angles $45^\circ, 135^\circ, 225^\circ, \dots$ all have $\sin^2 = \frac12$.
The iteration rotates by $2\theta = 90^\circ$, which exchanges the two diagonals and never brings the state closer to the marked axis. Grover gives no gain at $M = N/2$.

\textbf{\color{brand}Why the other options are incorrect.}
\begin{itemize}[leftmargin=2.4em, itemsep=1pt, topsep=2pt]
\item[\textbf{(A)}] The angles $45^\circ, 135^\circ, \dots$ all have $\sin^2 = \frac12$, so there is no oscillation between $0$ and $1$.
\item[\textbf{(B)}] That would need $3\theta = 180^\circ$, i.e.\ $\theta = 60^\circ$.
\item[\textbf{(C)}] That would need $3\theta = 90^\circ$, i.e.\ $\theta = 30^\circ$.
\end{itemize}
\textbf{\color{brand}Key takeaway.} At $M = N/2$, classical guessing is already as good as Grover.
\end{solbox}
\qsep

\needspace{13\baselineskip}
\begin{qitem}{51}
\qmeta{Single-correct MCQ}{$M = 3N/4$}
Let $N = 8$ with $M = 6$ marked items. The probability of a marked outcome after one iteration is
\par\smallskip\noindent\begin{tabularx}{\linewidth}{@{}*{4}{L}@{}}
\optlabel{A}\,$\dfrac34$ & \optlabel{B}\,$\dfrac14$ & \optlabel{C}\,$1$ & \correctlabel{D}\,$0$
\end{tabularx}
\end{qitem}
\begin{solbox}{D}
\textbf{\color{brand}Concepts required.} $M > N/2$ gives $\theta > 45^\circ$, and iterations can rotate past the marked axis at once.

\textbf{\color{brand}Relevant formulae.} $\sin^2\theta = M/N$; $\;P_1 = \sin^2 3\theta$; $\;\sin 180^\circ = 0$.

\textbf{\color{brand}Step-by-step explanation.} $\sin^2\theta = \frac34$ gives $\theta = 60^\circ$, so $P_1 = \sin^2 180^\circ = 0$. Amplitudes confirm this: with $a = 1/\sqrt8$ the oracle gives six entries $-a$ and two entries $+a$; the mean is $-\frac12a$;
the marked entries become $-a + a = 0$. Measuring \emph{without} any iteration ($P_0 = \frac34$) is far better.

\textbf{\color{brand}Why the other options are incorrect.}
\begin{itemize}[leftmargin=2.4em, itemsep=1pt, topsep=2pt]
\item[\textbf{(A)}] $\frac34$ is $P_0$, the probability with no iteration.
\item[\textbf{(B)}] $\frac14$ is the initial probability of the unmarked items.
\item[\textbf{(C)}] $1$ would need $3\theta = 90^\circ$; here $3\theta = 180^\circ$.
\end{itemize}
\textbf{\color{brand}Key takeaway.} If most items are marked, do not iterate: just measure.
\end{solbox}
\qsep

\needspace{13\baselineskip}
\begin{qitem}{52}
\qmeta{Single-correct MCQ}{Over-rotation at $2k_{\rm opt}$}
On $2^{15}$ items with one marked item, $k_{\rm opt} = 142$. If $k = 284$ iterations are used instead, the probability of the marked item is about
\par\smallskip\noindent\begin{tabularx}{\linewidth}{@{}*{4}{L}@{}}
\correctlabel{A}\,$3\times10^{-6}$ & \optlabel{B}\,$\dfrac12$ & \optlabel{C}\,$\dfrac{1}{32768}$ & \optlabel{D}\,$1$
\end{tabularx}
\end{qitem}
\begin{solbox}{A}
\textbf{\color{brand}Concepts required.} $P_k$ is periodic: at $k \approx 2k_{\rm opt}$ the state has rotated by nearly $\pi$ and returns to near the unmarked axis.

\textbf{\color{brand}Relevant formulae.} $P_k = \sin^2((2k+1)\theta)$; $\;\sin^2(\pi + \delta) \approx \delta^2$ for a small $\delta$.

\textbf{\color{brand}Step-by-step explanation.} The angle is $(2k+1)\theta = 569\times0.0055243 = 3.1433$ rad, which is $0.0017$ past $\pi$. So $P = \sin^2(569\theta) \approx (0.0017)^2 \approx 3\times10^{-6}$.
This is over-rotation: twice the optimal number of iterations brings the state back near the unmarked axis, and the marked item is then \emph{less} likely than a random guess ($1/32768 \approx 3\times10^{-5}$).

\textbf{\color{brand}Why the other options are incorrect.}
\begin{itemize}[leftmargin=2.4em, itemsep=1pt, topsep=2pt]
\item[\textbf{(B)}] $\frac12$ is the value around $k \approx k_{\rm opt}/2$, or $k \approx \frac32 k_{\rm opt}$.
\item[\textbf{(C)}] $\frac{1}{32768} \approx 3\times10^{-5}$ is the starting value at $k = 0$, about ten times larger than the value at $k = 284$.
\item[\textbf{(D)}] $1$ would be right near $k_{\rm opt}$, not near $2k_{\rm opt}$.
\end{itemize}
\textbf{\color{brand}Key takeaway.} More iterations are not always better: Grover is a rotation, not a monotonic filter.
\end{solbox}
\qsep

\needspace{13\baselineskip}
\begin{qitem}{53}
\qmeta{Single-correct MCQ}{Size of the speed-up}
For $n = 20$ qubits with one marked item, Grover needs $804$ iterations. A classical random search needs on average about $2^{19}$ queries. The classical cost divided by the quantum cost is about
\par\smallskip\noindent\begin{tabularx}{\linewidth}{@{}*{4}{L}@{}}
\correctlabel{A}\,$650$ & \optlabel{B}\,$2$ & \optlabel{C}\,$1000$ & \optlabel{D}\,$20$
\end{tabularx}
\end{qitem}
\begin{solbox}{A}
\textbf{\color{brand}Concepts required.} Ratio $= \frac{N/2}{(\pi/4)\sqrt N} = \frac2\pi\sqrt N$.

\textbf{\color{brand}Relevant formulae.} Classical average $\approx N/2$; Grover $k \approx \frac\pi4\sqrt N$; ratio $= \frac{N/2}{(\pi/4)\sqrt N} = \frac2\pi\sqrt N$.

\textbf{\color{brand}Step-by-step explanation.} $2^{19} = 524\,288$ and $524\,288/804 \approx 652$. In general the ratio is $\frac{N/2}{(\pi/4)\sqrt N} = \frac{2}{\pi}\sqrt N$, and here $\frac2\pi\times1024 \approx 652$. The speed-up is quadratic, not exponential.

\textbf{\color{brand}Why the other options are incorrect.}
\begin{itemize}[leftmargin=2.4em, itemsep=1pt, topsep=2pt]
\item[\textbf{(B)}] $2$ would be a constant speed-up.
\item[\textbf{(C)}] $1000$ rounds $\sqrt N = 1024$ and forgets the factor $\frac2\pi$.
\item[\textbf{(D)}] $20$ is $n$; the ratio is not polynomial in $n$.
\end{itemize}
\textbf{\color{brand}Key takeaway.} Quadratic speed-up: the gain is $\frac2\pi\sqrt N$.
\end{solbox}
\qsep

\needspace{13\baselineskip}
\begin{qitem}{54}
\qmeta{Single-correct MCQ}{Shot statistics}
A Grover circuit with success probability $P = 0.9613$ is run for $2000$ shots. The expected marked count and its binomial standard deviation are about
\par\smallskip\noindent\begin{tabularx}{\linewidth}{@{}*{4}{L}@{}}
\optlabel{A}\,$1923 \pm 44$ & \correctlabel{B}\,$1923 \pm 9$ & \optlabel{C}\,$2000 \pm 0$ & \optlabel{D}\,$1923 \pm 77$
\end{tabularx}
\end{qitem}
\begin{solbox}{B}
\textbf{\color{brand}Concepts required.} Binomial: mean $SP$, standard deviation $\sqrt{SP(1-P)}$.

\textbf{\color{brand}Relevant formulae.} Binomial counts from $S$ shots: mean $= SP$, standard deviation $= \sqrt{SP(1 - P)}$.

\textbf{\color{brand}Step-by-step explanation.} Mean $= SP = 2000\times0.9613 \approx 1923$. Standard deviation $= \sqrt{SP(1-P)} = \sqrt{2000\times0.9613\times0.0387} = \sqrt{74.4} \approx 8.6$.
$\pm44 = \sqrt{1923}$ is the Poisson estimate, which is wrong here because $P$ is close to $1$.

\textbf{\color{brand}Why the other options are incorrect.}
\begin{itemize}[leftmargin=2.4em, itemsep=1pt, topsep=2pt]
\item[\textbf{(A)}] $\pm44 = \sqrt{1923}$ uses a Poisson estimate and ignores the factor $1 - P$.
\item[\textbf{(C)}] $P < 1$, so not every shot succeeds.
\item[\textbf{(D)}] $\pm77 = S(1-P)$ is the expected number of failures, not a spread.
\end{itemize}
\textbf{\color{brand}Key takeaway.} When $P$ is close to $1$, the spread is small: compare measurement against $\sqrt{S P(1-P)}$.
\end{solbox}
\qsep

\needspace{13\baselineskip}
\begin{qitem}{55}
\qmeta{Matching lists}{$k_{\rm opt}$ against $n$}
Match the number of qubits (one marked item) in List-I with $k_{\rm opt}$ in List-II. Use $k_{\rm opt} = \operatorname{round}\big(\frac{\pi}{4\theta} - \frac12\big)$.
\par\smallskip\noindent{\renewcommand{\arraystretch}{1.35}%
\begin{tabularx}{\linewidth}{@{}c>{\raggedright\arraybackslash}p{0.33\linewidth}@{\hspace{1.2em}}cL@{}}
\toprule
\multicolumn{2}{@{}l}{\textbf{List-I}} & \multicolumn{2}{l@{}}{\textbf{List-II}} \\ \midrule
\textbf{(P)} & $n = 2$ & \textbf{(1)} & $12$ \\
\textbf{(Q)} & $n = 6$ & \textbf{(2)} & $25$ \\
\textbf{(R)} & $n = 8$ & \textbf{(3)} & $1$ \\
\textbf{(S)} & $n = 10$ & \textbf{(4)} & $6$ \\
\bottomrule
\end{tabularx}}
\par\smallskip\noindent\begin{tabularx}{\linewidth}{@{}*{2}{L}@{}}
\correctlabel{A}\,P$\to$3, Q$\to$4, R$\to$1, S$\to$2 & \optlabel{B}\,P$\to$3, Q$\to$1, R$\to$4, S$\to$2 \\[3pt]
\optlabel{C}\,P$\to$3, Q$\to$4, R$\to$2, S$\to$1 & \optlabel{D}\,P$\to$4, Q$\to$3, R$\to$1, S$\to$2
\end{tabularx}
\end{qitem}
\begin{solbox}{A}
\textbf{\color{brand}Concepts required.} $k_{\rm opt} \approx \frac\pi4 2^{n/2}$, with exact rounding from $\frac{\pi}{4\theta} - \frac12$.

\textbf{\color{brand}Relevant formulae.} $\sin\theta = 2^{-n/2}$; $\;k_{\rm opt} = \operatorname{round}\big(\frac{\pi}{4\theta} - \frac12\big)$.

\textbf{\color{brand}Step-by-step explanation.} $n = 2$: $\theta = \frac\pi6$, $\frac{\pi}{4\theta} - \frac12 = 1$ (3). $\;n = 6$: $\theta \approx 0.1253$, giving $5.77 \to 6$ (4). $\;n = 8$: $\theta \approx 0.0625$, giving $12.06 \to 12$ (1).
$\;n = 10$: $\theta \approx 0.03126$, giving $24.63 \to 25$ (2). Hence P$\to$3, Q$\to$4, R$\to$1, S$\to$2.

\textbf{\color{brand}Why the other options are incorrect.}
\begin{itemize}[leftmargin=2.4em, itemsep=1pt, topsep=2pt]
\item[\textbf{(B)}] This exchanges Q and R: the iterations grow with $n$, so $n = 6$ cannot need more than $n = 8$.
\item[\textbf{(C)}] This exchanges R and S: $n = 8$ gives $12$ and $n = 10$ gives $25$.
\item[\textbf{(D)}] This exchanges P and Q: $n = 2$ gives $k = 1$ and $n = 6$ gives $k = 6$.
\end{itemize}
\textbf{\color{brand}Key takeaway.} Adding $2$ qubits roughly doubles $k_{\rm opt}$: $6, 12, 25$.
\end{solbox}
\qsep

\needspace{13\baselineskip}
\begin{qitem}{56}
\qmeta{Matching lists}{Blocks of the circuit}
Match each part of the Grover circuit in List-I with its effect in List-II.
\par\smallskip\noindent{\renewcommand{\arraystretch}{1.35}%
\begin{tabularx}{\linewidth}{@{}c>{\raggedright\arraybackslash}p{0.33\linewidth}@{\hspace{1.2em}}cL@{}}
\toprule
\multicolumn{2}{@{}l}{\textbf{List-I}} & \multicolumn{2}{l@{}}{\textbf{List-II}} \\ \midrule
\textbf{(P)} & $H^{\otimes n}$ on $\ket{0\dots0}$ & \textbf{(1)} & returns one bit string with its Born probability \\
\textbf{(Q)} & $X$ gates on the $0$ bits of $x^\ast$, multi-controlled $Z$, the same $X$ gates & \textbf{(2)} & creates the uniform superposition $\ket{s}$ \\
\textbf{(R)} & $H^{\otimes n}X^{\otimes n}$, multi-controlled $Z$, $X^{\otimes n}H^{\otimes n}$ & \textbf{(3)} & reflects about $\ket{s}$, up to a global phase \\
\textbf{(S)} & measuring all qubits & \textbf{(4)} & flips the sign of $\ket{x^\ast}$ only \\
\bottomrule
\end{tabularx}}
\par\smallskip\noindent\begin{tabularx}{\linewidth}{@{}*{2}{L}@{}}
\correctlabel{A}\,P$\to$2, Q$\to$4, R$\to$3, S$\to$1 & \optlabel{B}\,P$\to$2, Q$\to$4, R$\to$1, S$\to$3 \\[3pt]
\optlabel{C}\,P$\to$4, Q$\to$2, R$\to$3, S$\to$1 & \optlabel{D}\,P$\to$2, Q$\to$3, R$\to$4, S$\to$1
\end{tabularx}
\end{qitem}
\begin{solbox}{A}
\textbf{\color{brand}Concepts required.} The circuit has four distinct stages: preparation, oracle, diffusion and measurement.

\textbf{\color{brand}Relevant formulae.} $\ket{s} = H^{\otimes n}\ket{0\dots0}$; $\;U_\omega = I - 2\ket{x^\ast}\!\bra{x^\ast}$; $\;H^{\otimes n}X^{\otimes n}\,\mathrm{C}^{n-1}\mathrm{Z}\,X^{\otimes n}H^{\otimes n} = -D$.

\textbf{\color{brand}Step-by-step explanation.} P: $H^{\otimes n}\ket{0\dots0} = \ket{s}$ (2). Q: the oracle recipe of Q.45 flips only $\ket{x^\ast}$ (4). R: this is the diffusion circuit of Q.47 (3). S: the measurement samples one string from $|c_x|^2$ (1).
Hence P$\to$2, Q$\to$4, R$\to$3, S$\to$1.

\textbf{\color{brand}Why the other options are incorrect.}
\begin{itemize}[leftmargin=2.4em, itemsep=1pt, topsep=2pt]
\item[\textbf{(B)}] This exchanges R and S: the diffusion is unitary and measurement is not.
\item[\textbf{(C)}] This exchanges P and Q: the Hadamard layer prepares $\ket{s}$, and the $X$-MCZ-$X$ block marks.
\item[\textbf{(D)}] This exchanges Q and R: the block with $H$ gates is the diffusion, not the oracle.
\end{itemize}
\textbf{\color{brand}Key takeaway.} Oracle: reflect about the marked axis. Diffusion: reflect about $\ket{s}$.
\end{solbox}
\qsep

\needspace{13\baselineskip}
\begin{qitem}{57}
\qmeta{Assertion--Reason}{Independence of the marked item}
\textbf{Assertion (A):} The success probability of Grover's algorithm does not depend on which single item is marked.\\[1mm]
\textbf{Reason (R):} The oracle is a diagonal matrix in the computational basis.
\begin{itemize}[leftmargin=2.6em, labelsep=0.4em, itemsep=2pt, topsep=3pt]
\item[\optlabel{A}] Both Assertion (A) and Reason (R) are true, and R is the correct explanation of A
\item[\correctlabel{B}] Both A and R are true, but R is \emph{not} the correct explanation of A
\item[\optlabel{C}] A is true, but R is false
\item[\optlabel{D}] A is false, but R is true
\end{itemize}
\end{qitem}
\begin{solbox}{B}
\textbf{\color{brand}Concepts required.} Decide A and R first, then ask whether R causes A. The cause here is the symmetry of $\ket{s}$.

\textbf{\color{brand}Relevant formulae.} $\braket{x^\ast|s} = 1/\sqrt N$ for every basis state $\ket{x^\ast}$; $\;U_\omega = \operatorname{diag}\big((-1)^{f(x)}\big)$.

\textbf{\color{brand}Step-by-step explanation.} \textbf{A is true:} only $\braket{x^\ast|s} = 1/\sqrt N$ enters the analysis, and this is the same for every $x^\ast$.
\textbf{R is true:} $U_\omega = \operatorname{diag}\big((-1)^{f(x)}\big)$. However, being diagonal does not explain A: an arbitrary diagonal gate need not treat all items alike. The explanation is the symmetry of $\ket{s}$, which has the same overlap with every basis state. Answer: (B).

\textbf{\color{brand}Why the other options are incorrect.}
\begin{itemize}[leftmargin=2.4em, itemsep=1pt, topsep=2pt]
\item[\textbf{(A)}] R is true but is not the explanation: not every diagonal gate treats all items equally.
\item[\textbf{(C)}] R is true: $U_\omega$ is diagonal with entries $\pm1$.
\item[\textbf{(D)}] A is true: $\braket{x^\ast|s} = 1/\sqrt N$ for every $x^\ast$.
\end{itemize}
\textbf{\color{brand}Key takeaway.} True does not imply explanatory: in Assertion--Reason questions, check for a causal link.
\end{solbox}
\qsep

\needspace{13\baselineskip}
\begin{qitem}{58}
\qmeta{Statement I / II}{Oscillation and $\sqrt N$ scaling}
\textbf{Statement I:} In Grover's algorithm, applying more iterations never lowers the success probability.\\[1mm]
\textbf{Statement II:} For one marked item, the optimal number of iterations grows as $\sqrt N$.
\begin{itemize}[leftmargin=2.6em, labelsep=0.4em, itemsep=2pt, topsep=3pt]
\item[\optlabel{A}] Both Statement I and Statement II are true
\item[\optlabel{B}] Both Statement I and Statement II are false
\item[\optlabel{C}] Statement I is true, but Statement II is false
\item[\correctlabel{D}] Statement I is false, but Statement II is true
\end{itemize}
\end{qitem}
\begin{solbox}{D}
\textbf{\color{brand}Concepts required.} $P_k$ oscillates; $k_{\rm opt} \propto \sqrt N$.

\textbf{\color{brand}Relevant formulae.} $P_k = \sin^2((2k+1)\theta)$; $\;k_{\rm opt} \approx \frac{\pi}{4\theta} \approx \frac\pi4\sqrt N$.

\textbf{\color{brand}Step-by-step explanation.} \textbf{I is false:} $P_k = \sin^2((2k+1)\theta)$ oscillates. For $N = 4$ we have $P_1 = 1$ but $P_2 = \sin^2 150^\circ = \frac14$.
\textbf{II is true:} $k_{\rm opt} \approx \frac{\pi}{4\theta} \approx \frac\pi4\sqrt N$. Answer: (D).

\textbf{\color{brand}Why the other options are incorrect.}
\begin{itemize}[leftmargin=2.4em, itemsep=1pt, topsep=2pt]
\item[\textbf{(A)}] I is false: over-rotation is real, for instance for $N = 4$, $k = 2$.
\item[\textbf{(B)}] II is true: $k_{\rm opt} \approx \frac\pi4\sqrt N$.
\item[\textbf{(C)}] I is false and II is true, which is the opposite of this option.
\end{itemize}
\textbf{\color{brand}Key takeaway.} One counterexample is enough to reject a ``never'' statement.
\end{solbox}
\qsep

\needspace{20\baselineskip}
\begin{passagebox}{Paragraph for Questions 59 and 60}
Grover's algorithm runs on $n = 3$ qubits with the single marked item $x^\ast = \ket{101}$. It starts from the uniform state $\ket{s}$, in which every
amplitude is $a = 1/\sqrt8$. One iteration is the oracle $U_\omega$ followed by the diffusion $D$. Trace the amplitudes by hand with inversion about the mean.
\end{passagebox}
\needspace{13\baselineskip}
\begin{qitem}{59}
\qmeta{Paragraph based}{Amplitude trace, first iteration}
After the first iteration the amplitudes are
\par\smallskip\noindent\begin{tabularx}{\linewidth}{@{}*{2}{L}@{}}
\optlabel{A}\,marked $\frac32a$; each unmarked $\frac12a$ & \optlabel{B}\,marked $\frac74a$; each unmarked $\frac34a$ \\[3pt]
\correctlabel{C}\,marked $\frac52a$; each unmarked $\frac12a$ & \optlabel{D}\,marked $\frac52a$; each unmarked $\frac34a$
\end{tabularx}
\end{qitem}
\begin{solbox}{C}
\textbf{\color{brand}Concepts required.} Inversion about the mean with the mean computed after the oracle.

\textbf{\color{brand}Relevant formulae.} $\bar c = \frac1N\sum_xc_x$; marked: $2\bar c + a$; unmarked: $2\bar c - a$, with $a = 1/\sqrt8$ the initial amplitude.

\textbf{\color{brand}Step-by-step explanation.} After $U_\omega$: seven entries $a$ and one entry $-a$; the mean is $\frac{7a - a}{8} = \frac34a$. Inversion about the mean gives: unmarked $2\cdot\frac34a - a = \frac12a$; marked $2\cdot\frac34a + a = \frac52a$.
\emph{Check:} $(7\cdot\frac14 + \frac{25}{4})a^2 = 8a^2 = 1$.

\textbf{\color{brand}Why the other options are incorrect.}
\begin{itemize}[leftmargin=2.4em, itemsep=1pt, topsep=2pt]
\item[\textbf{(A)}] $\frac32a = 2\times$mean omits the term $+a$ for the marked item: $2\cdot\frac34a - (-a) = \frac52a$.
\item[\textbf{(B)}] $\frac74a$ and $\frac34a$ use the mean as the amplitude and do not preserve the length.
\item[\textbf{(D)}] The unmarked value is wrong: $\frac34a$ is the mean itself.
\end{itemize}
\textbf{\color{brand}Key takeaway.} Marked: $2\text{mean} + a$. Unmarked: $2\text{mean} - a$.
\end{solbox}
\qsep

\needspace{13\baselineskip}
\begin{qitem}{60}
\qmeta{Paragraph based}{Amplitude trace, second iteration}
After the second iteration the probability of measuring $\ket{101}$ is
\par\smallskip\noindent\begin{tabularx}{\linewidth}{@{}*{4}{L}@{}}
\optlabel{A}\,$\dfrac{25}{32}$ & \optlabel{B}\,$\dfrac{11}{16}$ & \correctlabel{C}\,$\dfrac{121}{128}$ & \optlabel{D}\,$1$
\end{tabularx}
\end{qitem}
\begin{solbox}{C}
\textbf{\color{brand}Concepts required.} Iterate the same two steps; the second mean is smaller because the marked amplitude is now large.

\textbf{\color{brand}Relevant formulae.} Repeat the oracle ($c \mapsto -c$ on the marked item) and diffusion ($c_x \mapsto 2\bar c - c_x$); $\;P = |c_{\rm marked}|^2$; check: $P_2 = \sin^2 5\theta$.

\textbf{\color{brand}Step-by-step explanation.} Start from the result of Q.59: unmarked $\frac12a$, marked $\frac52a$. After $U_\omega$ the marked entry is $-\frac52a$, and the mean is $\frac{7\cdot\frac12a - \frac52a}{8} = \frac{a}{8}$.
The marked entry becomes $2\cdot\frac a8 + \frac52a = \frac{11}{4}a$, so $P = \frac{121}{16}a^2 = \frac{121}{16\cdot8} = \frac{121}{128}$. (Each unmarked entry becomes $\frac a4 - \frac a2 = -\frac a4$.)
\emph{Check:} $P_2 = \sin^2 5\theta \approx 0.945$ with $\sin^2\theta = \frac18$. A third iteration would over-rotate.

\textbf{\color{brand}Why the other options are incorrect.}
\begin{itemize}[leftmargin=2.4em, itemsep=1pt, topsep=2pt]
\item[\textbf{(A)}] $\frac{25}{32}$ is the value after the \emph{first} iteration.
\item[\textbf{(B)}] $\frac{11}{16}$ uses $\frac{11}{4}$ without squaring correctly with $a^2 = \frac18$.
\item[\textbf{(D)}] $P = 1$ would need $5\theta = 90^\circ$ exactly; here it is $103.5^\circ$.
\end{itemize}
\textbf{\color{brand}Key takeaway.} $k = 2$ is optimal for $N = 8$, with $P = \frac{121}{128} \approx 0.945$.
\end{solbox}
\qsep

\newpage

\begin{sectionbanner}SECTION D \enspace$\cdot$\enspace Advanced Grover and Implementation \hfill Q.61--80\end{sectionbanner}

\needspace{13\baselineskip}
\begin{qitem}{61}
\qmeta{Single-correct MCQ}{Several solutions: $k_{\rm opt}$}
For $N = 1024$ with $M = 16$ marked items, the optimal number of iterations is
\par\smallskip\noindent\begin{tabularx}{\linewidth}{@{}*{4}{L}@{}}
\optlabel{A}\,$25$ & \correctlabel{B}\,$6$ & \optlabel{C}\,$5$ & \optlabel{D}\,$8$
\end{tabularx}
\end{qitem}
\begin{solbox}{B}
\textbf{\color{brand}Concepts required.} With $M$ solutions $\sin\theta = \sqrt{M/N}$, and $k_{\rm opt} \approx \frac\pi4\sqrt{N/M}$.

\textbf{\color{brand}Relevant formulae.} $\sin\theta = \sqrt{M/N}$; $\;k_{\rm opt} = \operatorname{round}\big(\frac{\pi}{4\theta} - \frac12\big) \approx \frac\pi4\sqrt{N/M}$.

\textbf{\color{brand}Step-by-step explanation.} $\sin\theta = \sqrt{16/1024} = \frac18$, so $\theta = 0.12533$ rad. Then $\frac{\pi}{4\theta} - \frac12 = 6.267 - 0.5 = 5.767$, which rounds to $6$. The success probability is $\sin^2 13\theta \approx 0.9966$.
($25$ is the value for one marked item; $8 = \sqrt{N/M}$ forgets $\frac\pi4$.)

\textbf{\color{brand}Why the other options are incorrect.}
\begin{itemize}[leftmargin=2.4em, itemsep=1pt, topsep=2pt]
\item[\textbf{(A)}] $25$ is the value for $M = 1$ and $N = 1024$.
\item[\textbf{(C)}] $5$ rounds down instead of to the nearest integer: $5.77$ rounds to $6$.
\item[\textbf{(D)}] $8 = \sqrt{N/M}$ forgets $\frac\pi4$.
\end{itemize}
\textbf{\color{brand}Key takeaway.} More solutions mean fewer iterations: divide $N$ by $M$ before taking the root.
\end{solbox}
\qsep

\needspace{13\baselineskip}
\begin{qitem}{62}
\qmeta{Single-correct MCQ}{BBHT threshold for unknown $M$}
When $M$ is unknown, the algorithm of Boyer, Brassard, H{\o}yer and Tapp picks $k$ uniformly from $\{0, \dots, m-1\}$. The average success probability is at least $\frac14$ once $m \ge 1/\sin2\theta$.
For $N = 2^{20}$ and $M = 1$ this threshold on $m$ is about
\par\smallskip\noindent\begin{tabularx}{\linewidth}{@{}*{4}{L}@{}}
\correctlabel{A}\,$512$ & \optlabel{B}\,$1024$ & \optlabel{C}\,$804$ & \optlabel{D}\,$2^{20}$
\end{tabularx}
\end{qitem}
\begin{solbox}{A}
\textbf{\color{brand}Concepts required.} The BBHT lemma: the average of $P_k$ over $k < m$ is $\frac12 - \frac{\sin4m\theta}{4m\sin2\theta}$.

\textbf{\color{brand}Relevant formulae.} $\frac1m\sum_{k=0}^{m-1}\sin^2((2k+1)\theta) = \frac12 - \frac{\sin4m\theta}{4m\sin2\theta}$, for $k$ uniform on $\{0, \dots, m-1\}$; for small $\theta$, $\sin2\theta \approx 2\sin\theta$.

\textbf{\color{brand}Step-by-step explanation.} $\sin\theta = 2^{-10}$ is small, so $\sin 2\theta \approx 2\sin\theta = 2^{-9}$ and $\frac{1}{\sin2\theta} \approx 2^{9} = 512$. In general the threshold is about $\frac12\sqrt{N/M}$, which comes from the identity
$\frac1m\sum_{k=0}^{m-1}\sin^2((2k+1)\theta) = \frac12 - \frac{\sin 4m\theta}{4m\sin2\theta}$; the last term is at most $\frac14$ in size once $m \ge 1/\sin2\theta$.

\textbf{\color{brand}Why the other options are incorrect.}
\begin{itemize}[leftmargin=2.4em, itemsep=1pt, topsep=2pt]
\item[\textbf{(B)}] $1024 = \sqrt N$ forgets $\sin2\theta \approx 2\sin\theta$.
\item[\textbf{(C)}] $804 \approx \frac\pi4\sqrt N$ is $k_{\rm opt}$, not the threshold.
\item[\textbf{(D)}] $2^{20} = N$ would destroy the speed-up.
\end{itemize}
\textbf{\color{brand}Key takeaway.} Random $k$ below about $\frac12\sqrt{N/M}$ already gives average success $\ge \frac14$.
\end{solbox}
\qsep

\needspace{13\baselineskip}
\begin{qitem}{63}
\qmeta{Single-correct MCQ}{Exact Grover search}
To obtain success probability exactly $1$, the angle is reduced from $\theta$ to $\theta'$ so that $(2k+1)\theta' = \pi/2$ with $k = \lceil\frac{\pi}{4\theta} - \frac12\rceil$. For $N = 16$, $M = 1$ this requires
\par\smallskip\noindent\begin{tabularx}{\linewidth}{@{}*{2}{L}@{}}
\optlabel{A}\,$k = 3$ and $\theta' = \pi/12$ & \optlabel{B}\,$k = 4$ and $\theta' = \pi/18$ \\[3pt]
\optlabel{C}\,$k = 2$ and $\theta' = \pi/10$ & \correctlabel{D}\,$k = 3$ and $\theta' = \pi/14$
\end{tabularx}
\end{qitem}
\begin{solbox}{D}
\textbf{\color{brand}Concepts required.} An exact search needs $(2k+1)\theta' = \frac\pi2$ with $\theta' \le \theta$, so round $k$ \emph{up}.

\textbf{\color{brand}Relevant formulae.} $k = \lceil\frac{\pi}{4\theta} - \frac12\rceil$; $\;(2k+1)\theta' = \frac\pi2$ with the reduced angle $\theta' \le \theta$.

\textbf{\color{brand}Step-by-step explanation.} $\frac{\pi}{4\theta} - \frac12 = 2.608$, so the ceiling is $k = 3$. Then $7\theta' = \frac\pi2$ gives $\theta' = \frac{\pi}{14} \approx 0.2244$, slightly below $\theta \approx 0.2527$ as required.
One way to reduce the angle is to mark only with amplitude $\sin\theta'/\sin\theta$ by adding an ancilla rotation. The ceiling matters: rounding down to $k = 2$ would need $\theta' = \frac{\pi}{10} > \theta$, which cannot be reached by reducing the marked weight.

\textbf{\color{brand}Why the other options are incorrect.}
\begin{itemize}[leftmargin=2.4em, itemsep=1pt, topsep=2pt]
\item[\textbf{(A)}] $\theta' = \frac{\pi}{12}$ gives $7\theta' = \frac{7\pi}{12} \ne \frac\pi2$.
\item[\textbf{(B)}] $k = 4$ also satisfies the equation with $\theta' = \frac{\pi}{18}$, but it is not the ceiling and costs one more iteration.
\item[\textbf{(C)}] $k = 2$ needs $\theta' = \frac{\pi}{10} > \theta$, and the angle cannot be increased by reducing the marked weight.
\end{itemize}
\textbf{\color{brand}Key takeaway.} Round up, then shrink the angle a little, and the probability becomes exactly $1$.
\end{solbox}
\qsep

\needspace{13\baselineskip}
\begin{qitem}{64}
\qmeta{Single-correct MCQ}{Optimality (BBBV, Zalka)}
Which statement about the query complexity of unstructured search is correct?
\begin{itemize}[leftmargin=2.6em, labelsep=0.4em, itemsep=2pt, topsep=3pt]
\item[\optlabel{A}] Grover's algorithm is optimal only when $N$ is a power of $2$
\item[\optlabel{B}] Quantum algorithms need only $O(\log N)$ queries, because $n = \log_2 N$ qubits hold all $N$ items
\item[\optlabel{C}] The lower bound $\Omega(\sqrt N)$ holds only for classical algorithms
\item[\correctlabel{D}] Every quantum algorithm needs $\Omega(\sqrt N)$ oracle queries, so Grover's algorithm is optimal up to a constant
\end{itemize}
\end{qitem}
\begin{solbox}{D}
\textbf{\color{brand}Concepts required.} BBBV lower bound: $\Omega(\sqrt N)$ queries for black-box search. Zalka: $\frac\pi4\sqrt N$ is exactly optimal.

\textbf{\color{brand}Relevant formulae.} Quantum query lower bound: $\Omega(\sqrt{N/M})$; Grover uses $\frac\pi4\sqrt{N/M}$; classical cost: $\Theta(N/M)$.

\textbf{\color{brand}Step-by-step explanation.} The lower bound of Bennett, Bernstein, Brassard and Vazirani states that any quantum algorithm with black-box access needs $\Omega(\sqrt N)$ queries; Zalka showed that $\frac\pi4\sqrt N$ is optimal even including the constant.
Superposition of $N$ items does not give $O(\log N)$ search, since a measurement returns only one string and the marked amplitude grows slowly.

\textbf{\color{brand}Why the other options are incorrect.}
\begin{itemize}[leftmargin=2.4em, itemsep=1pt, topsep=2pt]
\item[\textbf{(A)}] The result does not depend on $N$ being a power of $2$.
\item[\textbf{(B)}] Holding $N$ amplitudes in $n$ qubits does not allow reading them out, and $\log N$ queries would violate the lower bound.
\item[\textbf{(C)}] Classical search needs $\Theta(N)$; the bound $\Omega(\sqrt N)$ is the quantum one.
\end{itemize}
\textbf{\color{brand}Key takeaway.} No black-box quantum search can beat $\sqrt N$ by more than a constant.
\end{solbox}
\qsep

\needspace{13\baselineskip}
\begin{qitem}{65}
\qmeta{Single-correct MCQ}{Reading gate statistics}
The notebook's full-circuit statistics for $n = 15$ and $k = 142$ report $4275$ \texttt{h} gates, $5964$ \texttt{x} gates and $284$ multi-controlled $Z$ gates. In one iteration the diffusion uses $2n$ \texttt{x} gates and the oracle uses two \texttt{x} gates for each $0$ in $x^\ast$. The number of $0$s in the marked string is
\par\smallskip\noindent\begin{tabularx}{\linewidth}{@{}*{4}{L}@{}}
\optlabel{A}\,$9$ & \optlabel{B}\,$21$ & \correctlabel{C}\,$6$ & \optlabel{D}\,$12$
\end{tabularx}
\end{qitem}
\begin{solbox}{C}
\textbf{\color{brand}Concepts required.} Per iteration the circuit has $2n$ H, $2n + 2z$ X and two multi-controlled $Z$ gates; the first H layer adds $n$ H gates.

\textbf{\color{brand}Relevant formulae.} \texttt{h} count $= n + 2nk$; \texttt{x} count $= k(2n + 2z)$; multi-controlled $Z$ count $= 2k$, with $k$ iterations and $z$ zeros in $x^\ast$.

\textbf{\color{brand}Step-by-step explanation.} \texttt{x} gates per iteration: $\frac{5964}{142} = 42$. Of these, $2n = 30$ belong to the diffusion, which leaves $12 = 2z$ for the oracle, so $z = 6$. The marked string in that run, $011001010011111$, indeed has six $0$s.
\emph{Check of the other counts:} \texttt{h}: $15 + 142\times30 = 4275$; multi-controlled $Z$: $2\times142 = 284$.

\textbf{\color{brand}Why the other options are incorrect.}
\begin{itemize}[leftmargin=2.4em, itemsep=1pt, topsep=2pt]
\item[\textbf{(A)}] $9$ is the number of $1$s in the marked string, not of $0$s.
\item[\textbf{(B)}] $21$ divides $42$ by $2$ without first removing the $2n$ gates of the diffusion.
\item[\textbf{(D)}] $12 = 2z$ counts the oracle's \texttt{x} gates, not the zeros.
\end{itemize}
\textbf{\color{brand}Key takeaway.} Gate statistics can be decoded with the formula $4n + 2z + 2$: here $60 + 12 + 2 = 74$ gates per iteration.
\end{solbox}
\qsep

\needspace{13\baselineskip}
\begin{qitem}{66}
\qmeta{Single-correct MCQ}{Amplitude amplification cost}
Amplitude amplification boosts a procedure $\mathcal{A}$ that succeeds with probability $a = 0.01$ ($\theta \approx 0.10017$). At the optimal $k$, one uses $\mathcal{A}$ or $\mathcal{A}^{-1}$ in total
\par\smallskip\noindent\begin{tabularx}{\linewidth}{@{}*{4}{L}@{}}
\optlabel{A}\,$7$ times & \correctlabel{B}\,$15$ times & \optlabel{C}\,$100$ times & \optlabel{D}\,$8$ times
\end{tabularx}
\end{qitem}
\begin{solbox}{B}
\textbf{\color{brand}Concepts required.} $Q = -\mathcal{A}S_0\mathcal{A}^{-1}S_\chi$: each iteration uses $\mathcal{A}$ and $\mathcal{A}^{-1}$ once, plus one initial $\mathcal{A}$.

\textbf{\color{brand}Relevant formulae.} $\sin^2\theta = a$, the success probability of $\mathcal{A}$; $\;Q = -\mathcal{A}S_0\mathcal{A}^{-1}S_\chi$; total uses of $\mathcal{A}$ and $\mathcal{A}^{-1}$: $2k + 1$.

\textbf{\color{brand}Step-by-step explanation.} Optimal $k$: $\frac{\pi}{4\theta} - \frac12 = 7.34$, so $k = 7$. The initial state uses $\mathcal{A}$ once, and each iteration $Q = -\mathcal{A}S_0\mathcal{A}^{-1}S_\chi$ uses $\mathcal{A}$ and $\mathcal{A}^{-1}$ once each. Total $2k + 1 = 15$ uses.
Classical repetition needs about $1/a = 100$ runs on average.

\textbf{\color{brand}Why the other options are incorrect.}
\begin{itemize}[leftmargin=2.4em, itemsep=1pt, topsep=2pt]
\item[\textbf{(A)}] $7$ is the iteration count, not the number of uses of $\mathcal{A}$.
\item[\textbf{(C)}] $100$ is the classical repetition count.
\item[\textbf{(D)}] $8$ is $k + 1$ and ignores the two uses per iteration.
\end{itemize}
\textbf{\color{brand}Key takeaway.} $2k + 1$ uses of $\mathcal{A}$, compared with $1/a$ classical repetitions.
\end{solbox}
\qsep

\needspace{13\baselineskip}
\begin{qitem}{67}
\qmeta{Single-correct MCQ}{Notebook memory estimate}
On the CPU state-vector backend the notebook estimates the memory as $2^n \times 16$ bytes for the simulator's state plus one equal host copy, because \texttt{RUN\_ITERATION\_SWEEP = True} holds two states at once. Its report for $n = 15$ is $1.05$ MB. With the same settings and $n = 28$ the estimate is
\par\smallskip\noindent\begin{tabularx}{\linewidth}{@{}*{4}{L}@{}}
\optlabel{A}\,$2$ GiB & \optlabel{B}\,$16$ GiB & \optlabel{C}\,$4$ GiB & \correctlabel{D}\,$8$ GiB
\end{tabularx}
\end{qitem}
\begin{solbox}{D}
\textbf{\color{brand}Concepts required.} State vector: $2^n$ amplitudes of $16$ bytes in double precision. The sweep keeps one additional copy.

\textbf{\color{brand}Relevant formulae.} Memory $= (1 + c)\times2^n\times16$ bytes, with $c = 1$ extra copy of the state; $\;1$ GiB $= 2^{30}$ bytes.

\textbf{\color{brand}Step-by-step explanation.} One state: $2^{28}\times16 = 2^{32}$ bytes $= 4$ GiB. Two states: $8$ GiB. \emph{Check against the report:} for $n = 15$, $2\times2^{15}\times16 = 1\,048\,576$ bytes $= 1.05$ MB in decimal units. Each extra qubit doubles the requirement.
Switching off the sweep removes the copy and halves the figure.

\textbf{\color{brand}Why the other options are incorrect.}
\begin{itemize}[leftmargin=2.4em, itemsep=1pt, topsep=2pt]
\item[\textbf{(A)}] $2$ GiB is one state at $n = 27$.
\item[\textbf{(B)}] $16$ GiB is the value for $n = 29$.
\item[\textbf{(C)}] $4$ GiB is one state, without the extra copy.
\end{itemize}
\textbf{\color{brand}Key takeaway.} Memory doubles per qubit; settings such as the sweep can double it again.
\end{solbox}
\qsep

\needspace{13\baselineskip}
\begin{qitem}{68}
\qmeta{Single-correct MCQ}{Key search}
A cipher with a $128$-bit key is attacked by Grover's algorithm (one correct key, cost measured in iterations). The number of iterations is about
\par\smallskip\noindent\begin{tabularx}{\linewidth}{@{}*{4}{L}@{}}
\correctlabel{A}\,$\frac\pi4\cdot2^{64}$ & \optlabel{B}\,$\frac\pi4\cdot2^{127}$ & \optlabel{C}\,$2^{127}$ & \optlabel{D}\,$128$
\end{tabularx}
\end{qitem}
\begin{solbox}{A}
\textbf{\color{brand}Concepts required.} $k \approx \frac\pi4\sqrt N$ with $N = 2^{128}$. Symmetric keys lose half their bits against Grover.

\textbf{\color{brand}Relevant formulae.} $k \approx \frac\pi4\sqrt N$, with $N = 2^{b}$ keys of $b$ bits.

\textbf{\color{brand}Step-by-step explanation.} $N = 2^{128}$, so $k \approx \frac\pi4\sqrt N = \frac\pi4\cdot 2^{64}$. The classical average is $2^{127}$. Doubling the key length restores the original margin.
Each iteration also has to run the cipher reversibly, so the real cost is much larger than the iteration count.

\textbf{\color{brand}Why the other options are incorrect.}
\begin{itemize}[leftmargin=2.4em, itemsep=1pt, topsep=2pt]
\item[\textbf{(B)}] This uses $\sqrt N$ with the wrong exponent.
\item[\textbf{(C)}] $2^{127}$ is the classical average.
\item[\textbf{(D)}] $128$ is the key length, and search is not polynomial in it.
\end{itemize}
\textbf{\color{brand}Key takeaway.} Doubling the key length cancels Grover's advantage.
\end{solbox}
\qsep

\needspace{13\baselineskip}
\begin{qitem}{69}
\qmeta{Single-correct MCQ}{Noise and depth}
A simple noise model multiplies the signal by $f$ per iteration: $P_k \approx f^{k}\sin^2((2k+1)\theta) + (1 - f^{k})/N$. For $n = 10$, $k = 25$, $f = 0.99$ (take $0.99^{25} \approx 0.78$), the success probability is about
\par\smallskip\noindent\begin{tabularx}{\linewidth}{@{}*{4}{L}@{}}
\optlabel{A}\,$0.99$ & \optlabel{B}\,$0.50$ & \correctlabel{C}\,$0.78$ & \optlabel{D}\,$0.25$
\end{tabularx}
\end{qitem}
\begin{solbox}{C}
\textbf{\color{brand}Concepts required.} The signal decays as $f^k$, while the background $(1 - f^k)/N$ is very small.

\textbf{\color{brand}Relevant formulae.} $P_k \approx f^k\sin^2((2k+1)\theta) + (1 - f^k)/N$, with $f$ the signal retained per iteration.

\textbf{\color{brand}Step-by-step explanation.} For $n = 10$, $k = 25$ is optimal and $\sin^2(51\theta) \approx 0.9995$, while $(1 - f^k)/N \approx 0.22/1024$ is negligible. So $P \approx 0.78\times 0.9995 \approx 0.78$. Under this model, a depth of $k$ iterations costs a factor $f^k$;
since $k \propto \sqrt N$, this decay is what limits large searches on noisy hardware.

\textbf{\color{brand}Why the other options are incorrect.}
\begin{itemize}[leftmargin=2.4em, itemsep=1pt, topsep=2pt]
\item[\textbf{(A)}] $0.99$ is the fidelity of one iteration, not of $25$.
\item[\textbf{(B)}] $0.50$ has no basis in the model.
\item[\textbf{(D)}] $0.25$ confuses $k = 25$ with a probability.
\end{itemize}
\textbf{\color{brand}Key takeaway.} Long Grover circuits suffer from noise, since $k \propto \sqrt N$.
\end{solbox}
\qsep

\needspace{13\baselineskip}
\begin{qitem}{70}
\qmeta{Single-correct MCQ}{Toffoli count}
A multi-controlled $Z$ with $m$ controls can be built with $m - 1$ clean ancillas, $2(m-1)$ Toffoli gates and one CZ. For the $n = 10$ Grover register ($m = 9$), one full iteration (oracle and diffusion) needs how many Toffoli gates?
\par\smallskip\noindent\begin{tabularx}{\linewidth}{@{}*{4}{L}@{}}
\correctlabel{A}\,$32$ & \optlabel{B}\,$36$ & \optlabel{C}\,$16$ & \optlabel{D}\,$18$
\end{tabularx}
\end{qitem}
\begin{solbox}{A}
\textbf{\color{brand}Concepts required.} Each iteration has two multi-controlled $Z$ gates. Each costs $2(m-1)$ Toffolis with this construction.

\textbf{\color{brand}Relevant formulae.} Multi-controlled $Z$ with $m$ controls: $m - 1$ clean ancillas, $2(m - 1)$ Toffoli gates, one CZ. Two such gates per iteration.

\textbf{\color{brand}Step-by-step explanation.} One multi-controlled $Z$ with $m = 9$: $2(9 - 1) = 16$ Toffolis. An iteration has two of them, one in the oracle and one in the diffusion, so $32$ Toffolis. The gate count per iteration is linear in $n$.

\textbf{\color{brand}Why the other options are incorrect.}
\begin{itemize}[leftmargin=2.4em, itemsep=1pt, topsep=2pt]
\item[\textbf{(B)}] $36 = 4m$ also uses $m$ instead of $m - 1$.
\item[\textbf{(C)}] $16$ counts only one of the two multi-controlled $Z$ gates.
\item[\textbf{(D)}] $18 = 2m$ uses the wrong count per gate.
\end{itemize}
\textbf{\color{brand}Key takeaway.} The cost is linear in $n$ per iteration, and most of it is in the multi-controlled gates.
\end{solbox}
\qsep

\needspace{13\baselineskip}
\begin{qitem}{71}
\qmeta{Single-correct MCQ}{Simulation cost}
On a state-vector simulator the work grows as (iterations) $\times$ (amplitudes per gate) $\sim 2^{n/2}\cdot2^{n}$. Adding $2$ qubits multiplies the total work by about
\par\smallskip\noindent\begin{tabularx}{\linewidth}{@{}*{4}{L}@{}}
\optlabel{A}\,$16$ & \correctlabel{B}\,$8$ & \optlabel{C}\,$4$ & \optlabel{D}\,$2\sqrt2$
\end{tabularx}
\end{qitem}
\begin{solbox}{B}
\textbf{\color{brand}Concepts required.} Work $\sim 2^{1.5n}$: amplitudes double per qubit and iterations grow by $\sqrt2$ per qubit.

\textbf{\color{brand}Relevant formulae.} Work $\sim k\cdot2^n \sim 2^{n/2}\cdot2^n = 2^{1.5n}$.

\textbf{\color{brand}Step-by-step explanation.} Work $\propto 2^{1.5n}$, so $n \to n + 2$ gives $2^{1.5\times2} = 2^3 = 8$: amplitudes grow by $4$ and iterations by $2$. The number of gates per iteration also grows slowly with $n$, which this estimate ignores.

\textbf{\color{brand}Why the other options are incorrect.}
\begin{itemize}[leftmargin=2.4em, itemsep=1pt, topsep=2pt]
\item[\textbf{(A)}] $16$ double-counts: it would need work $\sim 2^{2n}$.
\item[\textbf{(C)}] $4$ counts only the growth in amplitudes.
\item[\textbf{(D)}] $2\sqrt2$ is the factor for adding \emph{one} qubit.
\end{itemize}
\textbf{\color{brand}Key takeaway.} The classical simulation cost is about $N^{1.5}$.
\end{solbox}
\qsep

\needspace{13\baselineskip}
\begin{qitem}{72}
\qmeta{Single-correct MCQ}{Reading the sweep code}
The notebook's sweep computes\\
\hspace*{1em}\texttt{sweep\_max = 2 * k\_optimal + 1}\\
\hspace*{1em}\texttt{stride = max(1, ceil(sweep\_max / SWEEP\_MAX\_POINTS))}\\
It records the point $k = 0$ and then evolves by \texttt{stride} iterations \texttt{while k\_done + stride <= sweep\_max}, recording one point each time. For $k_{\rm opt} = 142$ and \texttt{SWEEP\_MAX\_POINTS = 60}, the number of recorded points is
\par\smallskip\noindent\begin{tabularx}{\linewidth}{@{}*{4}{L}@{}}
\optlabel{A}\,$57$ & \optlabel{B}\,$285$ & \correctlabel{C}\,$58$ & \optlabel{D}\,$60$
\end{tabularx}
\end{qitem}
\begin{solbox}{C}
\textbf{\color{brand}Concepts required.} Read the loop: the point $k = 0$ is recorded first and then one point per stride while $k$ stays within the range.

\textbf{\color{brand}Relevant formulae.} $\texttt{stride} = \lceil\texttt{sweep\_max}/\texttt{SWEEP\_MAX\_POINTS}\rceil$; points $= 1 + \lfloor\texttt{sweep\_max}/\texttt{stride}\rfloor$.

\textbf{\color{brand}Step-by-step explanation.} $\texttt{sweep\_max} = 285$, and $\texttt{stride} = \lceil 285/60\rceil = \lceil 4.75\rceil = 5$. The loop records $k = 5, 10, \dots, 285$, which is $\frac{285}{5} = 57$ points, and with $k = 0$ the total is $58$. The stored output of the notebook says ``58 points''.
The rounding up means the loop may record fewer points than the requested maximum of $60$, never more.

\textbf{\color{brand}Why the other options are incorrect.}
\begin{itemize}[leftmargin=2.4em, itemsep=1pt, topsep=2pt]
\item[\textbf{(A)}] $57$ forgets the point $k = 0$.
\item[\textbf{(B)}] $285$ is the range of $k$, which would be recorded only with stride $1$.
\item[\textbf{(D)}] $60$ is the requested maximum; rounding the stride up gives fewer.
\end{itemize}
\textbf{\color{brand}Key takeaway.} \texttt{SWEEP\_MAX\_POINTS} is an upper limit: compute the stride to find the real number of points.
\end{solbox}
\qsep

\needspace{13\baselineskip}
\begin{qitem}{73}
\qmeta{Single-correct MCQ}{Interpreting $2000/2000$ shots}
A notebook run has theoretical success probability $P = 0.99998683$ and records $2000$ marked counts out of $2000$ shots. For independent shots, the probability of seeing all $2000$ shots succeed is about
\par\smallskip\noindent\begin{tabularx}{\linewidth}{@{}*{4}{L}@{}}
\optlabel{A}\,$0.99999$ & \optlabel{B}\,$0.026$ & \optlabel{C}\,$0.500$ & \correctlabel{D}\,$0.974$
\end{tabularx}
\end{qitem}
\begin{solbox}{D}
\textbf{\color{brand}Concepts required.} Independent shots multiply: $P(\text{all } S \text{ succeed}) = P^S \approx e^{-S(1-P)}$ for $P$ near $1$.

\textbf{\color{brand}Relevant formulae.} $P(\text{all } S \text{ shots succeed}) = P^S$; $\;\ln P \approx -(1 - P)$ for $P$ near $1$.

\textbf{\color{brand}Step-by-step explanation.} All shots succeed with probability $P^{2000}$. Using $\ln P \approx -(1 - P) = -1.317\times10^{-5}$:
\[ P^{2000} \approx e^{-2000\times1.317\times10^{-5}} = e^{-0.0263} \approx 0.974. \]
So $2000/2000$ is the expected outcome. It does \emph{not} show that $P$ is exactly $1$: the run is consistent with a small failure probability that $2000$ shots cannot resolve.

\textbf{\color{brand}Why the other options are incorrect.}
\begin{itemize}[leftmargin=2.4em, itemsep=1pt, topsep=2pt]
\item[\textbf{(A)}] $0.99999$ is $P$ itself, the value for a single shot.
\item[\textbf{(B)}] $0.026$ is about the probability of at least one failure, the complement of the answer.
\item[\textbf{(C)}] $0.500$ has no basis: $2000(1 - P)$ is much smaller than $\ln 2$.
\end{itemize}
\textbf{\color{brand}Key takeaway.} A perfect count is consistent with $P$ slightly below $1$; finite shots cannot prove $P = 1$.
\end{solbox}
\qsep

\needspace{13\baselineskip}
\begin{qitem}{74}
\qmeta{Single-correct MCQ}{Repeat-until-success cost}
Let $N = 8$, $M = 1$, and repeat a run with a fixed $k$ until the marked item is found, counting only oracle queries in the Grover iterations. Then $P_1 = \frac{25}{32}$ and $P_2 = \frac{121}{128}$. The expected numbers of queries for $k = 1$ and $k = 2$ are about
\par\smallskip\noindent\begin{tabularx}{\linewidth}{@{}*{4}{L}@{}}
\optlabel{A}\,$1$ and $2$ & \optlabel{B}\,$2.12$ and $1.28$ & \correctlabel{C}\,$1.28$ and $2.12$ & \optlabel{D}\,$0.78$ and $1.89$
\end{tabularx}
\end{qitem}
\begin{solbox}{C}
\textbf{\color{brand}Concepts required.} Expected cost of repeat-until-success is $k/P_k$.

\textbf{\color{brand}Relevant formulae.} Number of runs is geometric with mean $1/P_k$; expected queries $= k/P_k$, with $k$ queries per run.

\textbf{\color{brand}Step-by-step explanation.} The number of runs is geometric with mean $1/P_k$, and each run costs $k$ queries, so the expected cost is $k/P_k$. $k = 1$: $\frac{32}{25} = 1.28$. $k = 2$: $2\times\frac{128}{121} \approx 2.12$. So $k = 1$ is cheaper here,
although $k = 2$ is the value that maximises the success probability of a single run: the two criteria are different.

\textbf{\color{brand}Why the other options are incorrect.}
\begin{itemize}[leftmargin=2.4em, itemsep=1pt, topsep=2pt]
\item[\textbf{(A)}] Those are $k$, which ignores repetition after failed runs.
\item[\textbf{(B)}] The values are in the wrong order.
\item[\textbf{(D)}] These are $P_1$ and $k P_2$ rather than $k/P_k$.
\end{itemize}
\textbf{\color{brand}Key takeaway.} The $k$ that maximises $P_k$ need not minimise the expected cost.
\end{solbox}
\qsep

\needspace{13\baselineskip}
\begin{qitem}{75}
\qmeta{Matching lists}{Query complexities}
Match each task in List-I with its query complexity in List-II.
\par\smallskip\noindent{\renewcommand{\arraystretch}{1.35}%
\begin{tabularx}{\linewidth}{@{}c>{\raggedright\arraybackslash}p{0.33\linewidth}@{\hspace{1.2em}}cL@{}}
\toprule
\multicolumn{2}{@{}l}{\textbf{List-I}} & \multicolumn{2}{l@{}}{\textbf{List-II}} \\ \midrule
\textbf{(P)} & Unstructured search, classical & \textbf{(1)} & $\Theta(N^{1/3})$ \\
\textbf{(Q)} & Unstructured search, quantum & \textbf{(2)} & $\Theta(N^{2/3})$ \\
\textbf{(R)} & Collision finding in a 2-to-1 function, quantum & \textbf{(3)} & $\Theta(N)$ \\
\textbf{(S)} & Element distinctness, quantum & \textbf{(4)} & $\Theta(N^{1/2})$ \\
\bottomrule
\end{tabularx}}
\par\smallskip\noindent\begin{tabularx}{\linewidth}{@{}*{2}{L}@{}}
\correctlabel{A}\,P$\to$3, Q$\to$4, R$\to$1, S$\to$2 & \optlabel{B}\,P$\to$4, Q$\to$3, R$\to$1, S$\to$2 \\[3pt]
\optlabel{C}\,P$\to$3, Q$\to$1, R$\to$4, S$\to$2 & \optlabel{D}\,P$\to$3, Q$\to$4, R$\to$2, S$\to$1
\end{tabularx}
\end{qitem}
\begin{solbox}{A}
\textbf{\color{brand}Concepts required.} Known query complexities: search $N$ vs.\ $\sqrt N$, collision $N^{1/3}$, element distinctness $N^{2/3}$.

\textbf{\color{brand}Relevant formulae.} Queries as $N$ grows: $\Theta(N)$, $\Theta(N^{1/2})$, $\Theta(N^{1/3})$, $\Theta(N^{2/3})$, where $\Theta$ means growth up to constant factors.

\textbf{\color{brand}Step-by-step explanation.} Classical search must examine a constant fraction of the items: $\Theta(N)$ (3). Grover: $\Theta(N^{1/2})$ (4). Brassard--H{\o}yer--Tapp collision finding, optimal by Aaronson--Shi: $\Theta(N^{1/3})$ (1). Ambainis's quantum walk for element distinctness: $\Theta(N^{2/3})$ (2).
Hence P$\to$3, Q$\to$4, R$\to$1, S$\to$2.

\textbf{\color{brand}Why the other options are incorrect.}
\begin{itemize}[leftmargin=2.4em, itemsep=1pt, topsep=2pt]
\item[\textbf{(B)}] This exchanges classical and quantum search.
\item[\textbf{(C)}] This assigns $N^{1/3}$ to unstructured search, which contradicts the lower bound.
\item[\textbf{(D)}] This exchanges collision and element distinctness; the promise of 2-to-1 makes collision easier.
\end{itemize}
\textbf{\color{brand}Key takeaway.} More structure leads to lower query complexity.
\end{solbox}
\qsep

\needspace{13\baselineskip}
\begin{qitem}{76}
\qmeta{Matching lists}{Variants of Grover}
Match each method in List-I with its main idea in List-II.
\par\smallskip\noindent{\renewcommand{\arraystretch}{1.35}%
\begin{tabularx}{\linewidth}{@{}c>{\raggedright\arraybackslash}p{0.33\linewidth}@{\hspace{1.2em}}cL@{}}
\toprule
\multicolumn{2}{@{}l}{\textbf{List-I}} & \multicolumn{2}{l@{}}{\textbf{List-II}} \\ \midrule
\textbf{(P)} & Amplitude amplification & \textbf{(1)} & uses randomly chosen $k$ from a growing range when $M$ is unknown \\
\textbf{(Q)} & Quantum counting & \textbf{(2)} & changes the reflection phases so that extra steps cannot over-rotate \\
\textbf{(R)} & BBHT search & \textbf{(3)} & replaces $H^{\otimes n}$ by any state-preparation procedure $\mathcal{A}$ \\
\textbf{(S)} & Fixed-point search & \textbf{(4)} & applies phase estimation to $G$ to find $M$ \\
\bottomrule
\end{tabularx}}
\par\smallskip\noindent\begin{tabularx}{\linewidth}{@{}*{2}{L}@{}}
\optlabel{A}\,P$\to$3, Q$\to$1, R$\to$4, S$\to$2 & \correctlabel{B}\,P$\to$3, Q$\to$4, R$\to$1, S$\to$2 \\[3pt]
\optlabel{C}\,P$\to$3, Q$\to$4, R$\to$2, S$\to$1 & \optlabel{D}\,P$\to$4, Q$\to$3, R$\to$1, S$\to$2
\end{tabularx}
\end{qitem}
\begin{solbox}{B}
\textbf{\color{brand}Concepts required.} Each variant answers one weakness: arbitrary start, unknown $M$, or over-rotation.

\textbf{\color{brand}Relevant formulae.} $G$ has eigenvalues $e^{\pm2i\theta}$; $\;Q = -\mathcal{A}S_0\mathcal{A}^{-1}S_\chi$; $\;\sin^2\theta = M/N$.

\textbf{\color{brand}Step-by-step explanation.} Amplitude amplification generalises $\ket{s}$ to $\mathcal{A}\ket{0}$ (3). Quantum counting reads $\theta$ from the eigenvalues $e^{\pm2i\theta}$ of $G$, which yields $M$ (4). BBHT draws $k$ at random from a range that grows geometrically (1).
Fixed-point search uses reflections with phase $\pi/3$ and similar choices rather than $\pi$ (2). Hence P$\to$3, Q$\to$4, R$\to$1, S$\to$2.

\textbf{\color{brand}Why the other options are incorrect.}
\begin{itemize}[leftmargin=2.4em, itemsep=1pt, topsep=2pt]
\item[\textbf{(A)}] This exchanges Q and R.
\item[\textbf{(C)}] This exchanges R and S; fixed-point search changes the phases and BBHT randomises $k$.
\item[\textbf{(D)}] This exchanges P and Q; quantum counting uses phase estimation.
\end{itemize}
\textbf{\color{brand}Key takeaway.} Learn the problem each variant solves, not only its name.
\end{solbox}
\qsep

\needspace{13\baselineskip}
\begin{qitem}{77}
\qmeta{Assertion--Reason}{Grover and NP-completeness}
\textbf{Assertion (A):} Grover's algorithm solves NP-complete problems in polynomial time.\\[1mm]
\textbf{Reason (R):} For a search over $2^n$ candidates, Grover's algorithm uses about $\frac\pi4 2^{n/2}$ iterations.
\begin{itemize}[leftmargin=2.6em, labelsep=0.4em, itemsep=2pt, topsep=3pt]
\item[\optlabel{A}] Both Assertion (A) and Reason (R) are true, and R is the correct explanation of A
\item[\optlabel{B}] Both A and R are true, but R is \emph{not} the correct explanation of A
\item[\optlabel{C}] A is true, but R is false
\item[\correctlabel{D}] A is false, but R is true
\end{itemize}
\end{qitem}
\begin{solbox}{D}
\textbf{\color{brand}Concepts required.} Quadratic speed-up does not make an exponential polynomial.

\textbf{\color{brand}Relevant formulae.} $k \approx \frac\pi4\sqrt N = \frac\pi42^{n/2}$, with $N = 2^n$ candidates.

\textbf{\color{brand}Step-by-step explanation.} \textbf{R is true:} this is $\frac\pi4\sqrt N$ with $N = 2^n$. \textbf{A is false:} $2^{n/2}$ is still exponential in $n$. The speed-up is quadratic, and the black-box lower bound shows that generic search cannot do better. Answer: (D).

\textbf{\color{brand}Why the other options are incorrect.}
\begin{itemize}[leftmargin=2.4em, itemsep=1pt, topsep=2pt]
\item[\textbf{(A)}] A is false, so this option fails.
\item[\textbf{(B)}] A is false, so this option fails even though R is true.
\item[\textbf{(C)}] A is false and R is true, the opposite of this option.
\end{itemize}
\textbf{\color{brand}Key takeaway.} $2^{n/2}$ is still exponential.
\end{solbox}
\qsep

\needspace{13\baselineskip}
\begin{qitem}{78}
\qmeta{Statement I / II}{Unknown $M$; flat curve at $M = N/2$}
\textbf{Statement I:} The BBHT algorithm finds a marked item with an expected $O(\sqrt{N/M})$ queries without knowing $M$ in advance.\\[1mm]
\textbf{Statement II:} If exactly half of the items are marked, ordinary Grover iterations leave the success probability at $\frac12$.
\begin{itemize}[leftmargin=2.6em, labelsep=0.4em, itemsep=2pt, topsep=3pt]
\item[\correctlabel{A}] Both Statement I and Statement II are true
\item[\optlabel{B}] Both Statement I and Statement II are false
\item[\optlabel{C}] Statement I is true, but Statement II is false
\item[\optlabel{D}] Statement I is false, but Statement II is true
\end{itemize}
\end{qitem}
\begin{solbox}{A}
\textbf{\color{brand}Concepts required.} BBHT handles unknown $M$; and $\theta = 45^\circ$ gives a flat success curve.

\textbf{\color{brand}Relevant formulae.} Expected BBHT queries: $O(\sqrt{N/M})$; $\;\sin^2((2k+1)45^\circ) = \frac12$ for every integer $k$.

\textbf{\color{brand}Step-by-step explanation.} \textbf{I is true:} a growing range for the random $k$ reaches $m \sim \sqrt{N/M}$ after a geometric series of attempts, whose total cost is dominated by the final terms.
\textbf{II is true:} $\theta = 45^\circ$ and $\sin^2((2k+1)45^\circ) = \frac12$ for all $k$ (see Q.50). Answer: (A).

\textbf{\color{brand}Why the other options are incorrect.}
\begin{itemize}[leftmargin=2.4em, itemsep=1pt, topsep=2pt]
\item[\textbf{(B)}] Both statements are true, not false.
\item[\textbf{(C)}] II is also true, as $\sin^2$ of odd multiples of $45^\circ$ is $\frac12$.
\item[\textbf{(D)}] I is also true: the expected cost is $O(\sqrt{N/M})$.
\end{itemize}
\textbf{\color{brand}Key takeaway.} Unknown $M$ is solved, and $M = N/2$ is not a problem since simply measuring works.
\end{solbox}
\qsep

\needspace{20\baselineskip}
\begin{passagebox}{Paragraph for Questions 79 and 80}
\textbf{Quantum counting.} A search space of $N = 64$ items contains an unknown number $M$ of marked items; in fact $M = 16$. The Grover operator $G$ has
eigenvalues $e^{\pm 2i\theta}$ on the plane containing $\ket{s}$, with $\sin^2\theta = M/N$. Phase estimation of $G$ with $t = 3$ counting qubits is run on $\ket{s}$, which has equal weight on the two eigenvectors. The outcome $y \in \{0, \dots, 7\}$ gives the
phase estimate $\varphi \approx y/8$, where $e^{2\pi i\varphi}$ is the eigenvalue. From this one sets $\tilde\theta = \pi y/8$ and $\tilde M = N\sin^2\tilde\theta$.
\end{passagebox}
\needspace{13\baselineskip}
\begin{qitem}{79}
\qmeta{Paragraph based}{Quantum counting outcomes}
The two most likely outcomes are
\par\smallskip\noindent\begin{tabularx}{\linewidth}{@{}*{4}{L}@{}}
\optlabel{A}\,$y = 2$ and $y = 6$ & \correctlabel{B}\,$y = 1$ and $y = 7$ & \optlabel{C}\,$y = 1$ and $y = 2$ & \optlabel{D}\,$y = 0$ and $y = 4$
\end{tabularx}
\end{qitem}
\begin{solbox}{B}
\textbf{\color{brand}Concepts required.} $\varphi = \theta/\pi$ and $1 - \theta/\pi$. QPE peaks at the grid point nearest to $2^t\varphi$.

\textbf{\color{brand}Relevant formulae.} $e^{2\pi i\varphi} = e^{\pm2i\theta}$ gives $\varphi = \frac\theta\pi$ and $1 - \frac\theta\pi$; phase estimation peaks at the integer $y$ nearest to $2^t\varphi$, with $t$ counting qubits.

\textbf{\color{brand}Step-by-step explanation.} $\theta = \arcsin\sqrt{16/64} = \arcsin\frac12 = \frac\pi6$, so the eigenphases are $\varphi = \frac{\theta}{\pi} = \frac16$ and $\frac56$. On the grid $y/8$ these lie at $y = \frac86 = 1.33$ and $y = \frac{40}{6} = 6.67$; the closest integers are $1$ and $7$.
Each has probability about $0.35$. (The next most likely are $y = 2$ and $y = 6$, each about $0.09$.)

\textbf{\color{brand}Why the other options are incorrect.}
\begin{itemize}[leftmargin=2.4em, itemsep=1pt, topsep=2pt]
\item[\textbf{(A)}] $y = 2, 6$ are the next most likely, each with about $0.09$.
\item[\textbf{(C)}] Both eigenphases contribute, and they are symmetric: $y$ and $8 - y$.
\item[\textbf{(D)}] $y = 0, 4$ would be $\theta = 0$ and $\theta = \frac\pi2$.
\end{itemize}
\textbf{\color{brand}Key takeaway.} Outcomes come in pairs $y$ and $2^t - y$, which give the same $\sin^2\tilde\theta$.
\end{solbox}
\qsep

\needspace{13\baselineskip}
\begin{qitem}{80}
\qmeta{Paragraph based}{Quantum counting estimate}
If $y = 1$ is obtained, the estimate of $M$ is
\par\smallskip\noindent\begin{tabularx}{\linewidth}{@{}*{4}{L}@{}}
\optlabel{A}\,$\tilde M = 32$ & \optlabel{B}\,$\tilde M = 8$ & \correctlabel{C}\,$\tilde M \approx 9.4$ & \optlabel{D}\,$\tilde M = 16$
\end{tabularx}
\end{qitem}
\begin{solbox}{C}
\textbf{\color{brand}Concepts required.} $\tilde M = N\sin^2(\pi y/2^t)$; $\sin^2\frac\pi8 = \frac{1 - \cos(\pi/4)}{2}$.

\textbf{\color{brand}Relevant formulae.} $\tilde\theta = \frac{\pi y}{2^t}$; $\;\tilde M = N\sin^2\tilde\theta$; $\;\sin^2\alpha = \frac{1 - \cos2\alpha}{2}$.

\textbf{\color{brand}Step-by-step explanation.} $\tilde\theta = \frac\pi8$ and $\tilde M = 64\sin^2\frac\pi8 = 64\cdot\frac{1 - \cos(\pi/4)}{2} = 32\big(1 - \frac{\sqrt2}{2}\big) \approx 9.37$. The true value $16$ lies between the estimates from $y = 1$ ($9.4$) and $y = 2$ ($32$): three counting qubits are too few for an accurate count.
Each additional counting qubit halves the error in $\theta$.

\textbf{\color{brand}Why the other options are incorrect.}
\begin{itemize}[leftmargin=2.4em, itemsep=1pt, topsep=2pt]
\item[\textbf{(A)}] $32$ is the estimate from $y = 2$.
\item[\textbf{(B)}] $8$ uses $\sin^2\frac\pi8 \approx \frac18$ and is not the correct value.
\item[\textbf{(D)}] $16$ is the true value, but $y = 1$ gives a coarse estimate.
\end{itemize}
\textbf{\color{brand}Key takeaway.} More counting qubits give a finer grid and a better estimate of $M$.
\end{solbox}
\qsep

\end{document}